\documentclass[11pt]{article}
\usepackage[margin=1in]{geometry}
\usepackage{amsmath,amssymb,amsthm,mathtools}
\usepackage{array,booktabs,longtable,tabularx}
\usepackage{enumitem}
\usepackage{xcolor}
\usepackage{listings}
\usepackage{seqsplit}
\usepackage{microtype}
\usepackage[T1]{fontenc}
\usepackage{lmodern}
\usepackage{hyperref}
\hypersetup{colorlinks=false,linkcolor=blue!55!black,citecolor=blue!55!black,urlcolor=blue!55!black,pdfborder={0 0 0}}
\lstset{basicstyle=\ttfamily\small,breaklines=true,columns=fullflexible,keepspaces=true,frame=single,framerule=0.3pt}
\emergencystretch=3em
\hbadness=10000
\raggedbottom
\setlist[itemize]{topsep=3pt,itemsep=2pt,parsep=0pt,leftmargin=1.5em}
\setlist[enumerate]{topsep=3pt,itemsep=2pt,parsep=0pt,leftmargin=1.7em}
\newtheorem{theorem}{Theorem}[section]
\newtheorem{lemma}[theorem]{Lemma}
\newtheorem{proposition}[theorem]{Proposition}
\newtheorem{corollary}[theorem]{Corollary}
\theoremstyle{definition}
\newtheorem{definition}[theorem]{Definition}
\newtheorem{criterion}[theorem]{Criterion}
\newtheorem{protocol}[theorem]{Protocol}
\theoremstyle{remark}
\newtheorem{remark}[theorem]{Remark}
\newcommand{\Lean}[1]{\texttt{#1}}
\newcommand{\LeanName}[1]{{\ttfamily\footnotesize\detokenize{#1}}}
\newcommand{\PosRaw}{\mathsf{Pos}_{\mathsf{raw}}}
\newcommand{\SepBare}{\mathsf{Sep}_{\mathsf{bare}}}
\newcommand{\SepRes}{\mathsf{Sep}_{\mathsf{res}}}
\newcommand{\SepImg}{\mathsf{Sep}_{\mathsf{img}}}
\newcommand{\SepImgOcc}{\mathsf{Sep}_{\mathsf{img}}^{\mathsf{occ}}}
\newcommand{\LocalSepImgOcc}{\mathsf{Sep}_{\mathsf{img}}^{\mathsf{loc}}}
\newcommand{\Asep}{\mathsf{A}_{\mathsf{sep}}}
\newcommand{\SepClay}{\mathsf{Sep}_{\mathsf{Clay}}}
\newcommand{\Apos}{\mathsf{A}_{\mathsf{pos}}}
\newcommand{\PosClay}{\mathsf{Pos}_{\mathsf{Clay}}}
\newcommand{\RawLB}{\mathsf{RawLB}}
\newcommand{\AClass}{\mathsf{AASCClass}}
\newcommand{\AASCCons}{\mathsf{AASCConsequences}}
\newcommand{\ClassOf}[1]{\mathsf{AASCClass}(#1)}
\newcommand{\Kernel}{\mathsf{Kernel}}
\newcommand{\Aplus}{\mathsf{A}^{+}}
\newcommand{\CnfSATP}{\mathsf{CnfSATInPolyTime}}
\newcommand{\CnfLB}{\mathsf{CnfCounterexampleLowerBound}}
\newcommand{\CnfBareSep}{\mathsf{CnfCounterexampleLowerBound}}
\newcommand{\CnfImageSep}{\mathsf{CnfSATImageSeparatorBranch}}
\newcommand{\CnfBareNeg}{\mathsf{CnfSATBareNegativeBranch}}
\newcommand{\KP}{\mathsf{KernelPackage}}
\newcommand{\RSAT}{R_{\mathsf{SAT\text{-}Clay}}}
\newcommand{\UEAP}{\mathsf{UEAP}}
\newcommand{\ATS}{\mathsf{ATS}}
\newcommand{\NonDeg}{\mathsf{NonDegenerate}}
\newcommand{\APlusCert}{\mathsf{KernelAPlusAuditCertificate}}
\newcommand{\EndpointImage}{\mathsf{EndpointImage}}
\newcommand{\OrdNeg}{\mathsf{OrdNeg}}
\newcommand{\OrdNegUn}{\mathsf{OrdNeg}^{\mathsf{unclass}}}
\newcommand{\OrdNegCls}{\mathsf{OrdNeg}^{\mathsf{class}}}
\newcommand{\OrdNegStd}{\mathsf{OrdNeg}^{\mathsf{std}}}
\newcommand{\OrdNegDistinct}{\mathsf{OrdNeg}^{\mathsf{distinct}}}
\newcommand{\Aneg}{\mathsf{A}_{\mathsf{neg}}}
\newcommand{\Aord}{\mathsf{A}_{\mathsf{ord}}}
\newcommand{\EndpointStand}{\mathsf{EndpointStanding}}
\newcommand{\Occupies}{\mathsf{Occupies}}
\newcommand{\BoundaryCross}{\mathsf{BoundaryCross}}
\newcommand{\AMetric}{\mathsf{AMetricBoundary}}
\newcommand{\ProofR}{\mathsf{Proof}_{R}}
\newcommand{\Pclass}{\mathsf{P}}
\newcommand{\NPclass}{\mathsf{NP}}
\newcommand{\ClayCtx}{\mathcal{C}_{\mathsf{SAT}}(R,m)}
\newcommand{\CnfCtx}{\mathsf{CnfSATClayEndpointImageContext}}
\newcommand{\Der}{\operatorname{Der}}
\newcommand{\Constr}{\operatorname{Constr}}
\newcommand{\law}{\mathrel{\simeq_{\mathrm{gov}}}}
\newcommand{\Adm}{\mathrm{Adm}}
\newcommand{\St}{\mathrm{St}}
\newcommand{\RefK}{\mathrm{Ref}}
\newcommand{\Irr}{\mathrm{Irr}}
\newcommand{\ReportClay}[1]{\mathsf{Report}_{\mathsf{Clay}}(#1)}

\newcommand{\RSATctx}{R_{\mathsf{SAT}}^{\mathsf{ctx}}}
\newcommand{\EndpointUse}{\mathsf{EndpointUse}}
\newcommand{\EndpointDisc}{\mathsf{EndpointDisc}}
\newcommand{\IndDisc}{\mathsf{IndependentDisc}}
\newcommand{\GovEquiv}{\mathsf{GovEquivalent}}
\newcommand{\Bookkeep}{\mathsf{Bookkeeping}}
\newcommand{\DomShift}{\mathsf{DomainShift}}
\newcommand{\FixedDom}{\mathsf{FixedDomain}}
\newcommand{\NegThm}{\mathsf{OrdNegThm}}
\newcommand{\RawTrace}{\mathsf{RawTrace}}
\newcommand{\KernelCons}{\mathsf{KernelConsequence}}
\newcommand{\CnfPositiveEndpoint}{\mathsf{CnfPositiveEndpoint}}
\newcommand{\CnfDirectPos}{\mathsf{CnfSATDirectPositiveEndpointOccupation}}
\newcommand{\CnfIndSepDisc}{\mathsf{CnfSATIndependentSeparatorEndpointStatusDiscriminator}}
\newcommand{\CnfNoIndSepClass}{\mathsf{CnfNoIndependentSeparatingClassifier}}
\newcommand{\CnfOfficialNegUse}{\mathsf{CnfSATOfficialNegativeEndpointUse}}
\newcommand{\OfficialEndpointResolution}{\mathsf{OfficialEndpointResolution}}
\newcommand{\CnfNegOccExhaust}{\mathsf{CnfSATNegativeOccupationExhaustion}}
\newcommand{\CnfBookNegOnly}{\mathsf{CnfSATBookkeepingNegativeOnly}}
\newcommand{\CnfCarrierShift}{\mathsf{CnfSATCarrierShift}}
\newcommand{\CnfEndpointStatus}{\mathsf{CnfSATEndpointStatus}}
\newcommand{\CnfEndpointStatusOcc}{\mathsf{CnfSATEndpointStatusOccupation}}
\newcommand{\CnfGovEndpointUse}{\mathsf{CnfSATGovernedEndpointUse}}
\newcommand{\CnfGovEndpointStatusUse}{\mathsf{CnfSATGovernedEndpointStatusUse}}
\newcommand{\CnfSepClassInd}{\mathsf{CnfSeparatingClassifierIsIndependentSameDomain}}
\newcommand{\CnfSepClassifier}{\mathsf{CnfCandidateStatusClassifier}}
\newcommand{\CnfSameDomainSep}{\mathsf{CnfSameDomainSeparator}}
\newcommand{\CnfSeparatesPoly}{\mathsf{CnfClassifierSeparatesPolynomialCandidates}}
\newcommand{\DirPos}{\mathcal{P}_{\mathsf{dir}}}
\newcommand{\SepIndBridge}{\mathcal{I}_{\mathsf{sep}}}
\newcommand{\SepDisc}{\mathcal{D}_{\mathsf{sep}}}
\newcommand{\NoSepDisc}{\mathcal{N}_{\mathsf{sep}}}
\newcommand{\ModelOfCnf}{\mathsf{Model}_{\mathsf{CNF}}}
\newcommand{\ModelAdeq}{\mathsf{Adeq}_{\mathsf{CNF}}}
\newcommand{\ModelCoverage}{\mathsf{Coverage}_{\mathsf{CNF}}}
\newcommand{\ModelSound}{\mathsf{Sound}_{\mathsf{CNF}}}
\newcommand{\FailureAdeq}{\mathsf{FailAdeq}_{\mathsf{CNF}}}
\newcommand{\ImageExact}{\mathsf{ImageExact}_{\mathsf{CNF}}}
\newcommand{\PolyCand}{\mathsf{PolyCand}}
\newcommand{\FailsSAT}{\mathsf{FailsSAT}}
\newcommand{\Decode}{\mathsf{decode}}
\newcommand{\CnfModel}{\mathsf{CnfEncodedCandidateModel}}
\newcommand{\SepRel}{\mathsf{SepRel}}
\newcommand{\CandStatusClass}{\mathsf{CandCls}}
\newcommand{\PolyCandSep}{\mathsf{PolySep}}
\newcommand{\StdLB}{\mathsf{StdLB}}
\newcommand{\CandImageExcl}{\mathsf{CandImageExcl}}
\newcommand{\ThmLevelDisc}{\mathsf{ThmLevelDisc}}
\newcommand{\CandImageGov}{\mathsf{CandImageGov}}
\newcommand{\CandidateUseKind}{\mathsf{CandImageUseKind}}
\newcommand{\ProofSupportOnly}{\mathsf{ProofSupportOnly}}
\newcommand{\EndResolvingNegThm}{\mathsf{EndpointResolvingNegThm}}
\newcommand{\EndResolvingNonGov}{\mathsf{EndpointResolvingNonGov}}
\newcommand{\CarrierChangingLB}{\mathsf{CarrierChangingLB}}
\newcommand{\SameDomClass}{\mathsf{SameDomCls}}
\newcommand{\NonBookSep}{\mathsf{NonBook}}
\newcommand{\NonGovEqSep}{\mathsf{NonGovEq}}

\newcommand{\KSAT}{K_{\mathsf{SAT}}}
\newcommand{\CoreSAT}{C_{0,\mathsf{SAT}}}
\newcommand{\WeakSep}{W_{\mathsf{sep}}}
\newcommand{\RouteSAT}{\mathsf{RouteLocus}_{\mathsf{SAT}}}
\newcommand{\SepNotCoequal}{\mathsf{SepNotCoequalEndpointRole}}
\newcommand{\OrdNegRoleClass}{\mathsf{OrdinaryNegationRoleClassified}}
\newcommand{\NegComplexityLawful}{\mathsf{NegativeComplexityResultsLawful}}
\newcommand{\OfficialSATEval}{\mathsf{OSEval}}
\newcommand{\PosEndpointOutcome}{\mathsf{PosEndpointOutcome}}
\newcommand{\NegEndpointOutcome}{\mathsf{NegEndpointOutcome}}
\newcommand{\SATEndpointImageOcc}{\mathsf{SATEndpointImageOcc}}
\newcommand{\PositiveCandidateOcc}{\mathsf{PositiveCandidateOcc}}
\newcommand{\OutcomePvNP}{\mathsf{Outcome}_{\mathsf{PvNP}}}
\newcommand{\LawfulCoequalImageOcc}{\mathsf{LawfulCoequalImageOccupation}}
\newcommand{\NonPositiveStatusWork}{\mathsf{NonPositiveStatusWork}}
\newcommand{\EndpointStatusWork}{\mathsf{EndpointStatusWork}}
\newcommand{\SingleOccInterfaceSAT}{\mathsf{SingleOccupationInterface}_{\mathsf{SAT}}}
\newcommand{\ImageAdeq}{\mathsf{ImageAdeq}}
\newcommand{\OrdThm}{\mathsf{OrdThm}}
\newcommand{\OrdOfficialRes}{\mathsf{OrdOfficialRes}}
\newcommand{\EndpointRoleDiscipline}{\mathsf{EndpointRoleDiscipline}}
\newcommand{\OrdResolutionBridge}{\mathsf{OrdResolutionBridge}}
\newcommand{\EndpointCompleteEval}{\mathsf{EndpointCompleteEval}}
\newcommand{\UnresolvedEvalStatus}{\mathsf{UnresolvedEvalStatus}}
\newcommand{\EndpointOutcome}{\mathsf{EndpointOutcome}}
\newcommand{\EndpointClosureEval}{\mathsf{EndpointClosureEval}}
\newcommand{\EndpointBranchStanding}{\mathsf{EndpointBranchStanding}}
\newcommand{\EndpointCounterforce}{\mathsf{EndpointCounterforce}}
\newcommand{\EndpointCounterforceDiscipline}{\mathsf{EndpointCounterforceDiscipline}}
\newcommand{\SameModeCounterforce}{\mathsf{SameModeCounterforce}}
\newcommand{\LocalCountercase}{\mathsf{LocalCountercase}}
\newcommand{\LocalCountercaseStanding}{\mathsf{LocalCountercaseStanding}}
\newcommand{\ReductioCountercase}{\mathsf{ReductioCountercase}}
\newcommand{\OfficialCNFSATRoute}{\mathsf{OfficialCNFSATRoute}}
\newcommand{\ProjectionRoleContent}{\mathsf{ProjectionRoleContent}}
\newcommand{\EndpointAdmRel}{\mathsf{EndpointAdmRel}}
\newcommand{\LBNatSAT}{\mathsf{LB}^{\mathsf{nat}}_{\mathsf{SAT}}}
\newcommand{\EndpointImpactSAT}{\mathsf{E}_{\mathsf{SAT}}}
\newcommand{\BridgeNeutralizedSAT}{\mathsf{BridgeNeutralized}_{\mathsf{SAT}}}
\newcommand{\LowerBoundGov}{\mathsf{G}^{-}_{\mathsf{LB}}}

\title{
\(P=NP\) by Exclusion of the SAT Separator Endpoint\\[0.5em]
{\normalsize An AASC Endpoint-Rigidity Proof for Arbitrary Finite CNF-SAT}
}
\author{Amos Jay Maley}
\date{June 2026}

\begin{document}
\maketitle

\begin{abstract}
This manuscript proves \(P=NP\) by excluding the SAT separator endpoint on the unrestricted finite CNF-SAT carrier. The branch-exclusion argument is AASC-internal: AASC supplies the endpoint-rigidity argument that excludes the bare negative branch, while Cook--Levin/Karp supplies only the standard endpoint correspondence from polynomial-time CNF-SAT to \(P=NP\). The positive raw branch is \(\PosRaw := \CnfSATP\), asserting deterministic polynomial-time decidability of arbitrary finite CNF-SAT. The bare negative branch \(\SepBare\) is the raw complement of \(\PosRaw\) on the same carrier.

The proof is kernel-first. Theorem-bearing endpoint use is treated as governed construction, and governed construction is possible only under the kernel of admissibility, standing, reference, and irreversibility. The downstream AASC machinery is therefore not an external classifier or a source of theorem standing; it is the fixed-domain consequence layer forced by kernel necessity. Ordinary theoremhood, positive or negative, is kernel-internal when it is genuine theoremhood rather than raw trace. Ordinary theoremhood alone is distinct from official endpoint resolution: when ordinary mathematical practice uses a theorem as the official resolution of a fixed problem endpoint, the act fixes target, carrier, branch, theorem-status, downstream reuse, act-time finality, and redescription-stable endpoint role. That is exactly theorem-bearing endpoint use.

The decisive step is the SAT-local endpoint asymmetry. The positive branch directly occupies the constructive endpoint carrier. The raw negative branch, by contrast, first transports by the SAT lower-bound/no-positive equivalence into standard lower-bound normal form and candidate endpoint-image exclusion. That support-level lower-bound content is not automatically forbidden. The fixed SAT endpoint image has a single occupation interface: positive encoded non-failing candidate occupation. The candidate-image carrier is nevertheless adequate for the official raw split: the positive raw branch is image occupation, and the bare negative branch is candidate-image exclusion, i.e. non-occupation of that same image. When ordinary mathematical practice chooses the standard CNF-SAT endpoint route for official resolution, that projection fixes the endpoint-role content of both outcomes; a negative resolution cannot coherently accept the CNF-SAT route while refusing the candidate-image non-occupation role that the route assigns. Universal candidate failure is therefore non-occupation of that image, not coequal occupation of it. This single-interface fact alone is non-explosive: it does not forbid lower bounds, non-occupation theorems, or negative complexity results. The contradiction arises only when non-occupation is used as the official same-carrier endpoint-resolving separator occupation. Under that endpoint use, it induces an independent same-domain endpoint-status discriminator, while the same endpoint-use act triggers the kernel-forced no-independent-discriminator closure. Hence the SAT separator route is unavailable.

It follows that separator endpoint occupation is excluded. The final endpoint step does not use the false principle that raw truth implies theoremhood, and it does not route a temporary reductio assumption through a global negative endpoint outcome. It uses ordinary local reductio discipline: in a proof by cases or reductio, the exact complement branch may be introduced with temporary derivational standing as the live countercase. For the official SAT split, \(\SepBare\) enters the reductio route only when introduced as the live exact-complement assumption against \(\PosRaw\). That exact-complement use yields local endpoint-countercase standing, and the local countercase induces local separator occupation \(\LocalSepImgOcc\), which produces the same forbidden separator-discriminator contradiction as global separator occupation while remaining a local subproof role. The local assumption of \(\SepBare\) therefore discharges to \(\neg\SepBare\) by the exact-complement annotation-discharge rule: the proof-role annotation is not an added object-level premise. Raw complementarity then yields \(\PosRaw\). Therefore arbitrary finite CNF-SAT is decidable in deterministic polynomial time. The final passage from \(\CnfSATP\) to \(P=NP\) is the standard Cook--Levin/Karp endpoint correspondence. The proof of branch exclusion is AASC-internal; external complexity theory is used only to identify the official endpoint consequence.
\end{abstract}


\section*{Mathematical-Status Lock and Lean Support Boundary}
\addcontentsline{toc}{section}{Mathematical-Status Lock and Lean Support Boundary}
\paragraph{Mathematical-status lock.}
AASC is used here as a mathematical constraint formalism, not as philosophical commentary, metaphor, optional narration, or an interpretive overlay. The proof is the theorem chain stated in this manuscript. The AASC primitives used below--reference, standing, admissibility, irreversibility, bivalence, AMetric no-selector discipline, unique admissible interior, report preservation, and no-independent-discriminator closure--are theorem-bearing fixed-domain constraints. A same-mode objection must therefore identify a failure in a definition, theorem link, dependency claim, kernel-forcing argument, SAT bridge, endpoint-governance inference, or final endpoint correspondence. It is not a refutation merely to classify the vocabulary as unfamiliar or to demand a circuit-lower-bound, diagonalization, natural-proofs, algebrization, proof-complexity, or algorithm-construction route.

\paragraph{Lean support boundary.}
The manuscript proof is the mathematical theorem chain below. The Lean4 material in the appendix is audit support for the local SAT-operator bridge, the kernel consequence interface, and the proof-spine routing. It is not a substitute for the manuscript proof and does not claim to formalize the Cook--Levin/Karp endpoint correspondence from foundations. A challenge to the Lean layer must identify a concrete mismatch between the cited Lean-facing support surface and the manuscript theorem chain. It is not enough to observe that Lean is not itself the proof, because the manuscript does not claim that it is.

\section*{Read This First: Endpoint Proof Spine}
\addcontentsline{toc}{section}{Read This First: Endpoint Proof Spine}
The live proof spine is intentionally short. In this display, \(O_R^-\) abbreviates official negative endpoint-resolution use. That use is supplied only inside official SAT endpoint evaluation, not by raw lower-bound content alone. The bridge from ordinary mathematical practice to this endpoint-use predicate is explicit: ordinary official endpoint resolution, \(\OrdOfficialRes_R(B)\), is exactly theorem-bearing endpoint use of branch \(B\):
\[
\begin{aligned}
\OrdOfficialRes_R(B)&\Rightarrow \EndpointUse_R(B),\\[2pt]
\OfficialCNFSATRoute(R,model)&\Rightarrow \bigl(\Pclass=\NPclass\leftrightarrow\PosRaw\leftrightarrow\SATEndpointImageOcc(model)\bigr)\\[-1pt]
&\hspace{5em}\wedge\bigl(\Pclass\neq\NPclass\leftrightarrow\SepBare\leftrightarrow\CandImageExcl(model)\bigr),\\[2pt]
\NonDeg_{\mathsf{SAT}}(R,model) &\Rightarrow \KSAT \Rightarrow \Aplus_{\mathsf{SAT}},\\[2pt]
\SepBare &\Rightarrow \StdLB(model) \Rightarrow \CandImageExcl(model),\\[2pt]
\PosRaw &\leftrightarrow \SATEndpointImageOcc(model),\\[2pt]
\SepBare &\leftrightarrow \CandImageExcl(model),\\[2pt]
\CandImageExcl(model) &\Rightarrow \neg\SATEndpointImageOcc(model),\\[2pt]
\CandImageExcl(model)\wedge O_R^-
  &\Rightarrow \SepImgOcc,\\[2pt]
\SepImgOcc &\Rightarrow \EndpointUse_R(\SepImgOcc),\\[2pt]
\SepImgOcc &\Rightarrow \NonPositiveStatusWork_R(model),\\[2pt]
\NonPositiveStatusWork_R(model)&\Rightarrow \SepDisc(model),\\[2pt]
\EndpointUse_R(\SepImgOcc) &\Rightarrow \NoSepDisc(model) \Rightarrow \neg\SepDisc(model),\\[2pt]
\SepDisc(model)\wedge\neg\SepDisc(model) &\Rightarrow \bot.
\end{aligned}
\]
Thus \(\SepImgOcc\) is impossible. The final endpoint step is not an inference from a raw hypothetical negative case to official resolution, does not assume endpoint-complete evaluation, and does not promote a local reductio assumption into a global negative endpoint outcome. It uses raw bivalence plus local countercase discipline in the official endpoint-closure proof:
\[
\begin{aligned}
  &\PosRaw\vee\SepBare,\\[2pt]
  &\ReductioCountercase_R(\SepBare;\PosRaw)\Rightarrow\LocalCountercase_R(\SepBare),\\[2pt]
  &\LocalCountercase_R(\SepBare)\Rightarrow\EndpointCounterforce_R(\SepBare),\\[2pt]
  &\EndpointCounterforce_R(\SepBare)\Rightarrow\LocalSepImgOcc(R,model),\\[2pt]
  &\LocalSepImgOcc(R,model)\Rightarrow\NonPositiveStatusWork_R(model)\Rightarrow\SepDisc(model),\\[2pt]
  &\LocalSepImgOcc(R,model)\Rightarrow\NoSepDisc(model)\Rightarrow\neg\SepDisc(model),\\[2pt]
  &\bigl(\SepBare\text{ assumed as the exact complement reductio countercase}\bigr)\Rightarrow\bot,\\[2pt]
  &\neg\SepBare\Rightarrow\PosRaw.
\end{aligned}
\]
The proof does not infer standing from raw truth; it uses ordinary local assumption standing. When the sole local assumption \([\SepBare]^i\) is opened in a reductio as the exact complement assumption against the official SAT endpoint, the proof act is annotated as \(\ReductioCountercase_R(\SepBare;\PosRaw)\), then as the local endpoint countercase. The annotation is not an added premise; only that live exact-complement use of the discharged assumption has temporary derivational standing. That local standing is enough to induce the local separator route and the separator-discriminator contradiction, without converting the temporary assumption into a global endpoint outcome. The local countercase therefore discharges to \(\neg\SepBare\) by the annotation-discharge theorem: the only discharged object-level assumption is \([\SepBare]^i\). Raw complementarity then yields \(\PosRaw\Rightarrow P=NP\).
\paragraph{No one-step retyping.}
The proof does not infer \(\SepDisc(model)\), \(\SepImgOcc\), or endpoint-status governance directly from \(\SepBare\), \(\neg\PosRaw\), or ordinary lower-bound content. The raw negative branch first becomes standard SAT lower-bound normal form on the context-bound encoded candidate image. That normal form becomes candidate endpoint-image exclusion. Candidate-image exclusion remains lawful lower-bound content unless it is used as the official endpoint-resolving negative branch on the same fixed carrier. Only under that endpoint-use act does candidate-image exclusion become separator endpoint-image occupation and enter the separator-discriminator route. This is the safeguard against collapsing native lower-bound content into its own discriminator.

The candidate-image carrier is not a one-sided surrogate for the official split: on the context-bound carrier, \(\PosRaw\) is image occupation and \(\SepBare\) is image non-occupation. Ordinary theoremhood is proof legitimacy, not endpoint-role equivalence. If ordinary mathematical practice treats a theorem as the official resolution of the fixed endpoint, the act is \(\EndpointUse_R\): it fixes target, carrier, branch, theorem status, downstream reuse, act-time finality, and lawful-redescription stability. The fixed endpoint image has a single occupation interface: encoded non-failing polynomial-time candidate occupation. Universal candidate failure is therefore not a coequal occupant of that image; it is non-occupation of that image. Only when that non-occupation is used as the official endpoint resolution does it become separator endpoint-image occupation and independent endpoint-status work. This is the SAT analogue of the negative-branch bridge audit used throughout the AASC endpoint closure papers.

\paragraph{Role map for the separator route.}
The symbols in the proof spine deliberately mark different roles. Same-domain incompleteness objections are not part of the separator contradiction, but the final bridge records that any such objection matters only if it re-enters the fixed endpoint as standing-bearing gate work; otherwise it is endpoint-inert. The carrier-adequacy theorem pairs the official raw split with the image interface: \(\PosRaw\) is positive image occupation and \(\SepBare\) is candidate-image exclusion on the same exact encoded image. \(\CandImageExcl(model)\) is support-level universal candidate failure over that encoded candidate image. \(\SepImgOcc\) is endpoint-resolution use of that exclusion as separator occupation. \(\SATEndpointImageOcc(model)\) is positive image occupation by an encoded non-failing polynomial-time SAT candidate. \(\NonPositiveStatusWork_R(model)\) is same-carrier non-occupation status determination under official endpoint use. \(\SepDisc(model)\) is the independent separator discriminator induced by that endpoint-status work; after the admissibility bridge below, this means an independent endpoint-admissibility-relevant discriminator, not a mere descriptive label. \(\OrdOfficialRes_R(B)\) records ordinary mathematical use of a theorem as the official resolution of branch \(B\); by the ordinary-resolution bridge it is not a looser status than \(\EndpointUse_R(B)\). \(\EndpointCounterforce_R(B)\) records the narrower role in which a raw branch is asserted with force to defeat endpoint closure. \(\ReductioCountercase_R(\SepBare;\PosRaw)\) records the proof-act annotation generated when the sole local assumption \([\SepBare]^i\) is opened as the live exact-complement assumption against \(\PosRaw\), rather than as inert background syntax; it is not an added object-level premise, and the annotation-discharge theorem ensures that a contradiction in the annotated subproof discharges \([\SepBare]^i\) itself. \(\LocalCountercase_R(\SepBare)\) records the temporary derivational-standing role induced by that exact-complement use. Truth without endpoint standing is allowed as epistemic or raw status, but it has no endpoint-counterexample force. Single-interface typing alone proves none of these forbidden consequences; the contradiction requires either global separator occupation \(\SepImgOcc\) or local separator occupation \(\LocalSepImgOcc(R,model)\) to induce \(\SepDisc(model)\), together with the corresponding no-independent closure yielding \(\neg\SepDisc(model)\). \(\OfficialCNFSATRoute(R,model)\) records that the standard unrestricted finite CNF-SAT correspondence is being used as the route of official resolution, rather than merely as a tool for support-level lower-bound analysis. Ordinary theorem-bearing non-occupation is not a separate endpoint role: it is proof legitimacy for the act whose endpoint role, after the accepted CNF-SAT projection and endpoint use, is candidate-image non-occupation. Once support, bookkeeping, carrier shift, positive occupation, and coequal image occupation are removed, the endpoint-discriminator trichotomy leaves only the independent same-domain separator role. Once that independent role is identified, ordinary mathematical practice has no separate exemption from the consequence layer: ordinary official resolution is \(\EndpointUse_R\), endpoint use instantiates the kernel, and the kernel forces \(\Aplus\), including no-independent-discriminator closure.

\section*{Dependency-Inversion Exclusion}
\addcontentsline{toc}{section}{Dependency-Inversion Exclusion}
The dependency order is not
\[
  \text{AASC framework} \Rightarrow \text{SAT constraints}.
\]
The dependency order is
\[
  \text{non-degenerate fixed-carrier SAT endpoint regime}
  \Rightarrow \text{kernel governance already active}.
\]
The official P versus NP target fixes the carrier of evaluation in complexity-theoretic form: the unrestricted finite CNF-SAT carrier, the positive polynomial-time endpoint slot, and the raw negative complement branch. It does not assert the positive endpoint and does not supply truth of the theorem. Once a proof, counterproof, endpoint resolution, or theorem-bearing act is evaluated as a determinate same-carrier endpoint act at all, it already requires target determinacy, step evaluability, act-time finality, and same-regime fidelity. Those kernel-neutral requirements force the roles later named reference, standing, irreversibility, and admissibility. AASC names and audits the consequence layer of those forced roles; it does not impose an optional proof standard from outside complexity theory.

Thus a reading that treats the kernel as an optional framework premise is a dependency inversion. It changes the argument from a necessity proof into a framework-application proof. A same-mode denial must identify a specific failure in the passage from target adequacy to the kernel roles, exhibit a faithful weaker same-carrier endpoint regime, or break one of the SAT-native bridge links. A refusal of vocabulary alone does not reach the proof.

\tableofcontents
\newpage

\section{Proof Class, Endpoint Validity, and Kernel-First Discipline}

\subsection{Proof-class lock}
The proof class is
\[
  \textbf{kernel-forced endpoint-image discriminator closure.}
\]
It is not a circuit lower bound, diagonalization argument, relativization argument, natural-proofs argument, algebrization argument, proof-complexity lower bound, exhaustive-search lower bound, or empirical solver-hardness claim. It is a main-text mathematical proof about which theorem-bearing endpoint statuses can survive on the fixed SAT endpoint image once the kernel of non-degenerate construction is in force.

\subsection{Necessity, not optional framework}
The argument does not introduce AASC as an optional foundational system layered on top of the P versus NP problem. The kernel of non-degenerate construction, and the downstream consequences used in the proof, are derived as necessities forced by the minimal evaluability conditions required for any raw sequence to count as construction of a determinate same-domain object. These conditions are stated in kernel-neutral form before governed construction is defined. Once they are fixed, the governance roles of admissibility, standing, reference, and irreversibility follow as forced roles; the fixed-domain consequences of those roles include the structural rigidity principles that exclude independent same-domain endpoint-status discriminators. AASC is the systematic development of those necessity constraints. The positive endpoint is therefore not obtained by adopting a proprietary framework, but by applying the consequences of non-degenerate construction to the official unrestricted finite CNF-SAT carrier.

\subsection{Clay-valid as official mathematical endpoint}

\begin{definition}[Clay-valid mathematical endpoint]\label{def:clay-valid-math-endpoint}
In this manuscript, ``Clay-valid'' is used only in the mathematical endpoint sense. A Clay-valid endpoint targets the official P versus NP problem, preserves the standard definitions of \(\Pclass\), \(\NPclass\), polynomial time, SAT, CNF-SAT, polynomial-time reductions, and NP-completeness, and reaches one of the official mathematical branches. The term is used only for this mathematical endpoint property.
\end{definition}

\subsection{Endpoint and method}
The branch-exclusion step is supplied by the necessity constraints already forced by kernel-neutral target adequacy on the official carrier. AASC records those constraints and their fixed-domain consequences; it does not supply an optional foundational layer. The proof is not a diagonalization argument, circuit lower-bound argument, relativization argument, natural-proofs argument, proof-complexity argument, or constructive algorithm search. Classical complexity theory enters only at the endpoint-correspondence layer: unrestricted finite CNF-SAT is the standard NP-complete carrier, and deterministic polynomial-time decidability of that carrier entails \(\Pclass=\NPclass\) by the Cook--Levin/Karp correspondence. The proof of \(\neg\SepBare\) is AASC-internal.

\subsection{Kernel-first ordering}
The central dependency is
\[
  \EndpointUse_R(B)
  \Rightarrow
  \text{governed construction}
  \Rightarrow
  K_R
  \Rightarrow
  \Aplus_R(B).
\]
The reverse dependency is not used. AASC does not create standing, and reportability does not supply standing. Standing is one member of the kernel already presupposed by genuine governed construction. A proof of a proposition, when it is a proof rather than raw syntax, is therefore already kernel-internal.

\begin{proposition}[Downstream status of AASC machinery]\label{prop:aasc-downstream-status}
In this manuscript, AASC, UEAP, ATS, AMetricity, no-selector discipline, no-third-status discipline, unique-interior discipline, standing conservation, and report-preservation discipline are downstream consequence machinery forced by the kernel under fixed-domain theorem-bearing endpoint use. None of them is a primitive source of theorem standing.
\end{proposition}

\begin{proof}
The kernel proof in Section~\ref{sec:self-contained-kernel} begins with raw step-sequences and kernel-neutral target adequacy, then proves that target determinacy, step evaluability, act-time finality, and same-regime fidelity force reference, standing, irreversibility, and admissibility. The later AASC machinery is the endpoint interface of those roles: it preserves target, carrier, admissibility boundary, theorem-status reuse, redescription stability, and report integrity. Therefore it records forced consequences of the kernel rather than adding an independent source of standing.
\end{proof}

\subsection{The SAT context object is neutral}

\begin{definition}[Fixed SAT endpoint context]\label{def:context}
The context object
\[
  \ClayCtx
\]
records only the official SAT/P-vs-NP endpoint setting: the unrestricted finite CNF-SAT carrier, the raw positive branch \(\PosRaw\), the raw negative branch \(\SepBare\), candidate endpoint-image exclusion \(\CandImageExcl(model)\), the separator endpoint-image occupation slot \(\SepImgOcc\), and the fixed model/carrier on which the raw branches are evaluated. It is a carrier-and-slot object. It is not an A+ certificate, not an AMetric premise, not a no-selector axiom, not a separator-classification premise, and not a positive endpoint premise.
\end{definition}

\begin{theorem}[Context non-smuggling]\label{thm:context-non-smuggling}
The context object \(\ClayCtx\) does not assert any of
\[
  \Aplus_R(B),\quad \AMetric(R),\quad \Asep,\quad \neg\SepBare,
  \quad \neg\SepImgOcc,\quad \neg\LocalSepImgOcc(R,model),\quad \PosRaw.
\]
It fixes the official SAT carrier and branch slots only.
\end{theorem}

\begin{proof}
By Definition~\ref{def:context}, the clauses of \(\ClayCtx\) are carrier and slot clauses. They say which endpoint image is under evaluation and which raw branches are being evaluated on that image. They do not classify any branch, exclude any branch, or assert the positive branch. All endpoint-status consequences are derived later from endpoint use plus the kernel.
\end{proof}

\subsection{Endpoint layers}

\begin{definition}[Bare proposition layer]\label{def:bare}
The raw positive proposition is
\[
  \PosRaw := \CnfSATP.
\]
The bare negative proposition \(\SepBare\) is the raw complement branch for the same unrestricted finite CNF-SAT carrier.
\end{definition}

\begin{definition}[Candidate-image exclusion and separator occupation]\label{def:image}
The support-level SAT-native lower-bound image is \(\CandImageExcl(model)\): candidate endpoint-image exclusion on the exact context-bound encoded polynomial-time CNF-SAT candidate image. It is lower-bound content and may remain proof support, local obstruction, or ordinary theorem content.

The separator endpoint-image occupation \(\SepImgOcc\) denotes candidate-image exclusion after it is used as the official endpoint-resolving negative branch on the fixed SAT carrier. Its formal Lean carrier is the object \(\CnfImageSep(R,model)\), but the manuscript keeps the support-level content \(\CandImageExcl(model)\) distinct from endpoint occupation \(\SepImgOcc\). Thus the proof does not infer endpoint use from lower-bound content alone. The local variant \(\LocalSepImgOcc(R,model)\) records the same separator role inside a reductio countercase subproof; it is local countercase occupation, not global endpoint outcome.
\end{definition}

\begin{definition}[SAT endpoint-image occupation]\label{def:sat-endpoint-image-occupation}
The SAT endpoint-image occupation predicate \(\SATEndpointImageOcc(model)\) records occupation of the fixed positive CNF-SAT candidate image. On a context-bound encoded CNF-SAT model, write
\[
  \PositiveCandidateOcc(model)
  \quad\Longleftrightarrow\quad
  \exists c\,\bigl(\PolyCand_{model}(c)\wedge
  \neg\FailsSAT_{model}(\Decode_{model}(c))\bigr).
\]
The fixed endpoint-image occupation slot is typed by that positive candidate occupation:
\[
  \SATEndpointImageOcc(model)\quad\text{means}\quad
  \PositiveCandidateOcc(model).
\]
Thus an occupant of the image is an encoded deterministic polynomial-time CNF-SAT candidate whose decoded procedure has no SAT failure on arbitrary finite CNF inputs.
\end{definition}

\begin{definition}[Symmetric P versus NP outcome carrier]\label{def:symmetric-outcome-carrier}
The symmetric problem-outcome carrier is the two-output presentation
\[
  \OutcomePvNP=\{\Pclass=\NPclass,\;\Pclass\neq\NPclass\}.
\]
It is a legitimate broad problem-outcome presentation. It is not the fixed positive CNF-SAT candidate-image carrier used by the branch-exclusion proof. Moving from \(\SATEndpointImageOcc(model)\) to \(\OutcomePvNP\) changes the endpoint interface from candidate-image occupation to two-outcome classification.
\end{definition}

\begin{definition}[Official CNF-SAT resolution route]\label{def:official-cnf-sat-resolution-route}
\(\OfficialCNFSATRoute(R,model)\) means that the standard unrestricted finite CNF-SAT correspondence is being used as the route of official resolution for the ordinary \(P\) versus \(NP\) endpoint on the context-bound encoded candidate image. It is stronger than using CNF-SAT as a tool for lower-bound support, examples, reductions, or local obstruction analysis: it says that the official resolution is being read through the CNF-SAT endpoint carrier itself. The predicate does not assert the positive branch, does not assert the negative branch, and does not create endpoint standing. It fixes the route by which the symmetric outcome carrier is projected onto the SAT endpoint-image carrier.
\end{definition}

\begin{definition}[Endpoint classification layer]\label{def:clay-layer}
The symbol \(\Asep\) denotes the role-specific separator endpoint-status factor, if such a factor is asserted for endpoint-image separator occupation \(\SepImgOcc\). The Clay-valid separator endpoint is
\[
  \SepClay := \SepImgOcc\wedge\Asep.
\]
The positive report classification is \(\Apos\), and the Clay-valid positive endpoint is
\[
  \PosClay:=\PosRaw\wedge\Apos.
\]
These classification symbols are downstream status factors; they do not create kernel standing.

\end{definition}

\subsection{Official endpoint resolution as governed construction: preview}\label{subsec:official-resolution-preview}
The proof later verifies, after the kernel-neutral adequacy machinery has been established, that ordinary official endpoint resolution through the unrestricted finite CNF-SAT route is itself a target-adequate governed construction act. This subsection records the dependency order. The act of official endpoint resolution fixes the official target, fixes the CNF-SAT carrier and context-bound encoded model, fixes which branch is being used as the endpoint result, and fixes the standing, reportability, and licensed downstream reuse of that result. These are not optional AASC additions. They are the ordinary mathematical features by which a theorem is used as the official resolution of a fixed problem endpoint.

Consequently, the downstream \(\Aplus\) layer is not applied to endpoint resolution from outside. It applies only after the resolution act has satisfied the same kernel-neutral adequacy conditions used below to derive the kernel: target determinacy, step evaluability, act-time finality, and same-regime fidelity. The following bridge lemma states this dependency explicitly; its detailed verification is expanded after the kernel and \(\Aplus\) consequence layer in Theorem~\ref{thm:official-resolution-satisfies-target-adequacy}.

\begin{lemma}[Official endpoint resolution satisfies target adequacy]\label{lem:front-official-resolution-target-adequacy}
Let \(R\) be the official \(P\) versus \(NP\) endpoint context, let \(model\) be a context-bound adequate encoded CNF-SAT model, and suppose the resolution act is conducted through the standard unrestricted finite CNF-SAT route. Then ordinary official endpoint resolution satisfies the four kernel-neutral target-adequacy conditions: target determinacy, step evaluability, act-time finality, and same-regime fidelity. Consequently, for every branch \(B\),
\[
\begin{aligned}
  &\OrdOfficialRes_R(B)\wedge\OfficialCNFSATRoute(R,model)
  \wedge\ModelOfCnf(model,R)\wedge\FixedDom(R)\\
  &\qquad\Rightarrow K_R\wedge\Aplus_R(B).
\end{aligned}
\]
\end{lemma}

\begin{proof}
Target determinacy holds because the official route fixes the \(P\) versus \(NP\) endpoint, the unrestricted finite CNF-SAT carrier, the context-bound encoded candidate image, and the branch role being resolved. The repeatable object of evaluation is the standing-bearing endpoint role of the outcome on that image.

Step evaluability holds because a resolution act is not a bare inscription: it asserts a theorem as the endpoint result, classifies the standing and reportability of the resolved branch, and licenses downstream reuse. Under the accepted CNF-SAT projection, the positive branch has the endpoint-role content of candidate-image occupation and the negative branch has the endpoint-role content of candidate-image non-occupation, so the resolution steps are evaluable as advancing or departing from the fixed endpoint.

Act-time finality holds because a later reclassification of which branch has official standing, reportability, branch exclusion, or licensed downstream reuse produces a distinct endpoint-resolution act. It does not repair the earlier resolution as the same act.

Same-regime fidelity holds because lawful redescription that preserves the official CNF-SAT route and the context-bound encoded model preserves the target, the evaluability of the branch-classification steps, and the act-time status of the resolution. A change of route, carrier, model, branch role, standing assignment, or act-time status is a scope or carrier shift, not the same resolution act.

Thus the four neutral adequacy conditions hold for the resolution act itself. By the kernel-forcing results proved in Section~\ref{sec:self-contained-kernel}, the kernel roles are in force; by the fixed-domain consequence theorem for endpoint use, \(\Aplus\) applies. The full theorem-level verification of this bridge appears as Theorem~\ref{thm:official-resolution-satisfies-target-adequacy}. Consequently, any theorem-bearing endpoint classification within the resolution act that assigns official standing, reportability, branch exclusion, or licensed downstream reuse to one branch rather than the other on the fixed CNF-SAT carrier is admissibility-relevant for the governed resolution act.
\end{proof}

\begin{corollary}[Stability of official endpoint resolution against same-domain incompleteness]\label{cor:front-official-resolution-incompleteness-stability}
Once the kernel and \(\Aplus\) consequence layer are in force for official endpoint resolution on the unrestricted finite CNF-SAT carrier, no same-domain G\"odel-style incompleteness architecture can reopen endpoint standing, reportability, branch exclusion, or licensed downstream reuse for the separator branch.
\end{corollary}

\begin{proof}
By Lemma~\ref{lem:front-official-resolution-target-adequacy}, official endpoint resolution is a kernel-governed endpoint act with the \(\Aplus\) consequence layer in force. A same-domain incompleteness objection has endpoint force only if it re-enters the fixed endpoint domain by changing endpoint standing, reportability, branch exclusion, licensed downstream reuse, admissible continuation, or coherent endpoint trajectory. If it does not, it is endpoint-inert. If it does, it is governance-relevant gate work for the resolution act and is controlled by the same fixed-domain consequence layer. The companion non-instantiability theorem for G\"odel-style incompleteness objections exhausts same-domain code-level readback, semantic import, reflection or infinitary rule-power, second-equivalence routes, and richer same-domain extensions into inert bookkeeping, extensional gate renaming, finite-witness collapse, explicit scope change, or inadmissible authority import \cite{GodelNonInstantiability}. Hence no same-domain incompleteness architecture rehabilitates the separator branch inside the governed endpoint-resolution regime. The endpoint-level form is developed later in Corollary~\ref{cor:endpoint-resolution-stable-against-incompleteness}.
\end{proof}

\section{Official Target and SAT Endpoint Carrier}

\subsection{The official P versus NP target}
The official P versus NP problem asks whether every language accepted by a nondeterministic polynomial-time algorithm is also accepted by a deterministic polynomial-time algorithm. The problem statement is
\[
  \Pclass \stackrel{?}{=} \NPclass.
\]
The proof works through encoded CNF satisfiability as the SAT endpoint carrier.

\subsection{CNF satisfiability}
For a CNF formula \(\varphi\), an assignment \(a\) assigns Boolean values to its variables, and
\[
  SAT(\varphi) \quad\Longleftrightarrow\quad \exists a\, Eval(\varphi,a)=1.
\]
The formal carrier uses finite encodings of CNF formulas, the Boolean characteristic function \(\mathsf{satChar}\), and a Turing-machine polynomial-time certificate for the positive endpoint.

\begin{definition}[Unrestricted finite CNF carrier; bounded-CNF clarification]\label{def:bounded-cnf-clarification}
Throughout this manuscript, ``bounded-CNF'' means the ordinary finitely encoded CNF-SAT carrier: arbitrary finite CNF formulas over a finite Boolean alphabet, with input length equal to the size of the encoded formula. It does not mean fixed-width \(k\)-CNF, fixed clause count, fixed variable count, bounded instance size, bounded search depth, promise-restricted CNF, parameterized tractable fragments, or any other restricted SAT fragment. Thus
\[
  \CnfSATP
\]
is the assertion that arbitrary finite CNF-SAT is decidable by a deterministic polynomial-time machine.

\end{definition}

\subsection{Encoded candidate model adequacy}\label{subsec:encoded-candidate-model-adequacy}

This subsection fixes the ordinary coding burden for the unrestricted finite CNF-SAT carrier. The encoded candidate model is not an AASC source of standing, not a separator, and not a hidden proof device. It is the standard effective presentation of deterministic polynomial-time CNF-SAT candidates used to state the positive endpoint and its raw complement. Its role is to make explicit that the later branch proof quantifies over the full official candidate image.

\begin{definition}[Context-bound encoded CNF model]\label{def:cnf-model-of-context}
\[
  \ModelOfCnf(model,R)
\]
means that \(model\) is the encoded CNF-SAT endpoint model associated with the fixed SAT context \(R\): the carrier is arbitrary finite CNF over finite Boolean alphabets, the input measure is formula length, admissible encodings preserve the same endpoint-image carrier, and the encoded candidate image satisfies the adequacy packet below. The corresponding Lean structure is \LeanName{CnfEncodedCandidateModel}; the adequacy packet is \LeanName{CnfEncodedCandidateModelAdequate}.
\end{definition}

\begin{definition}[Encoded candidate model adequacy]\label{def:encoded-candidate-model-adequacy}
For a context-bound encoded CNF model, \(\ModelAdeq(model)\) is the conjunction of four ordinary coding conditions:
\[
\begin{aligned}
  \ModelAdeq(model) := {}& \ModelCoverage(model)\wedge\ModelSound(model)\\
  &{}\wedge\FailureAdeq(model)\wedge\ImageExact(model).
\end{aligned}
\]
Here:
\begin{enumerate}[label=(\roman*)]
\item \(\ModelCoverage(model)\) says every deterministic polynomial-time CNF-SAT candidate procedure has a code in the model;
\item \(\ModelSound(model)\) says every code marked polynomial-time decodes to a deterministic polynomial-time CNF-SAT candidate procedure;
\item \(\FailureAdeq(model)\) says SAT-failure status is invariant under equal decoded procedures, so no encoding artifact changes endpoint status;
\item \(\ImageExact(model)\) says the positive endpoint \(\CnfSATP\) is equivalent to the existence of an encoded polynomial-time candidate whose decoded procedure has no SAT failure.
\end{enumerate}
These clauses are formalized by \LeanName{CnfSATCandidateModelCoverage}, \LeanName{CnfSATCandidateModelSoundness}, \LeanName{CnfSATCandidateFailureAdequacy}, \LeanName{CnfSATCandidateImageExactness}, and \LeanName{CnfEncodedCandidateModelAdequate} in \LeanName{AnticheckerReduction.lean}.
\end{definition}


\begin{definition}[Standard encoded polynomial-time SAT model]\label{def:standard-encoded-cnf-model}
The standard encoded polynomial-time CNF-SAT model is the ordinary uniform machine model whose codes are pairs
\[
  (e,k),
\]
where \(e\) is an effective code for a deterministic Boolean procedure on finite CNF inputs and \(k\in\mathbb{N}\) is a declared global polynomial-time bound. A code is marked polynomial-time exactly when the decoded procedure is licensed as a uniform deterministic candidate running within its declared polynomial bound on every finite CNF input. The same decoded procedure handles all finite CNF inputs; no advice strings, promise restrictions, fixed-size test sets, nonuniform circuit families, or per-length classifiers are part of this model.
\end{definition}

\begin{proposition}[Standard model instantiation]\label{prop:standard-model-instantiation}
The standard encoded polynomial-time CNF-SAT model of Definition~\ref{def:standard-encoded-cnf-model} instantiates the adequacy packet of Definition~\ref{def:encoded-candidate-model-adequacy}. In particular, it satisfies coverage, soundness, failure adequacy, and candidate-image exactness for the unrestricted finite CNF-SAT carrier.
\end{proposition}

\begin{proof}
Deterministic procedures over finite binary strings, and hence deterministic procedures over finite CNF encodings, admit ordinary effective descriptions. Pairing such a description with a declared global polynomial bound gives the standard candidate code. Coverage follows because every uniform deterministic polynomial-time CNF-SAT decider has some effective description and some polynomial exponent or bound. Soundness follows because a code marked polynomial-time is admitted only when the decoded procedure runs within the declared global polynomial bound on every finite CNF input. Failure adequacy follows because SAT failure is a property of the decoded procedure on lawful finite CNF inputs, not a property of the presentation of its code. Candidate-image exactness follows because \(\CnfSATP\) is precisely the existence of a uniform deterministic polynomial-time procedure deciding arbitrary finite CNF-SAT correctly on all finite CNF inputs. Therefore the standard encoded model supplies the same candidate image as the ordinary definition of \(\Pclass\)-time CNF-SAT decidability.
\end{proof}

\begin{definition}[SAT-failure predicate and uniformity convention]\label{def:sat-failure-uniformity}
For an encoded polynomial-time candidate, \(\FailsSAT_{model}(\Decode_{model}(c))\) means that there exists a lawful finite CNF formula \(\varphi\) such that the decoded uniform deterministic procedure either returns the wrong Boolean value relative to \(SAT(\varphi)\), fails to return a lawful Boolean answer on the encoded input, or fails to meet the declared global polynomial-time bound on that input. The predicate is evaluated over arbitrary finite CNF instances, not over a promise subset, a fixed-width fragment, a bounded-size table, a finite benchmark set, or a nonuniform family.
\end{definition}

\begin{remark}[Duplicate encodings and extensional harmlessness]\label{rem:duplicate-encodings}
The model does not require a unique code for each extensional procedure. Multiple presentations may decode to the same candidate, or to extensionally equivalent candidates on all finite CNF inputs. Duplicate encodings do no endpoint work: endpoint status is read from the decoded procedure and the SAT-failure predicate. The formal Lean adequacy clause records invariance under equal decoded procedures; strengthening this to full extensional equivalence would leave the branch route unchanged, because the proof uses coverage, soundness, and candidate-image exactness rather than uniqueness of code representatives.
\end{remark}

\begin{theorem}[Context-bound model adequacy]\label{thm:context-bound-model-adequacy}
For any context-bound encoded CNF-SAT endpoint model,
\[
  \ModelOfCnf(model,R)\Rightarrow \ModelAdeq(model).
\]
\end{theorem}

\begin{proof}
By Definition~\ref{def:cnf-model-of-context}, \(\ModelOfCnf(model,R)\) fixes the ordinary effective presentation of deterministic polynomial-time CNF-SAT candidate procedures on the unrestricted finite CNF carrier. That presentation includes coverage, sound decoding, failure adequacy, and exactness of the encoded positive endpoint image. In Lean these clauses are assembled by \LeanName{cnfEncodedCandidateModelAdequate} from the model fields and the endpoint exactness theorem.
\end{proof}

\begin{theorem}[Encoded positive endpoint exactness]\label{thm:encoded-positive-endpoint-exactness}
On a context-bound adequate encoded CNF-SAT model,
\[
\begin{aligned}
  \ModelOfCnf(model,R)\Rightarrow{}&
  \Bigl(\CnfSATP \leftrightarrow
    \exists c\,\bigl(\PolyCand_{model}(c)\wedge{}\\
  &\qquad\neg\FailsSAT_{model}(\Decode_{model}(c))\bigr)\Bigr).
\end{aligned}
\]
The Lean anchor is \LeanName{cnfSATInPolyTime_iff_exists_encoded_nonfailing_candidate}.
\end{theorem}

\begin{proof}
The forward direction is the coverage clause applied to a deterministic polynomial-time CNF-SAT decider for the unrestricted carrier: such a candidate has a code marked polynomial-time, and because it decides SAT correctly it has no SAT failure. The reverse direction is soundness plus failure adequacy: if an encoded polynomial-time candidate has no SAT failure, the decoded candidate is a deterministic polynomial-time procedure deciding arbitrary finite CNF-SAT. This is exactly \(\CnfSATP\). The Lean theorem records this equivalence over \LeanName{CnfEncodedCandidateModel}.
\end{proof}

\begin{theorem}[Encoded negative endpoint exactness]\label{thm:encoded-negative-endpoint-exactness}
On a context-bound adequate encoded CNF-SAT model,
\[
\begin{aligned}
  \ModelOfCnf(model,R)\Rightarrow{}&
  \Bigl(\neg\CnfSATP \leftrightarrow
    \forall c\,\bigl(\PolyCand_{model}(c)\Rightarrow{}\\
  &\qquad \FailsSAT_{model}(\Decode_{model}(c))\bigr)\Bigr).
\end{aligned}
\]
The Lean anchor is \LeanName{cnfSATNotInPolyTime_iff_all_encoded_candidates_fail}.
\end{theorem}

\begin{proof}
Negating Theorem~\ref{thm:encoded-positive-endpoint-exactness} gives the standard lower-bound normal form on the exact encoded candidate image: if there is no deterministic polynomial-time CNF-SAT decider, then every encoded polynomial-time candidate fails SAT; conversely, if every encoded polynomial-time candidate fails, no encoded non-failing polynomial-time candidate exists, and by coverage no unencoded deterministic polynomial-time decider has been missed. This is the formal equivalence proved in \LeanName{cnfSATNotInPolyTime_iff_all_encoded_candidates_fail}.
\end{proof}

\begin{proposition}[Standing burden for model-adequacy objections]\label{prop:model-adequacy-objection-burden}
A model-adequacy objection has standing only if it identifies a failure of standard-model instantiation, coverage, soundness, failure adequacy, or image exactness. Equivalently, the objector must exhibit a deterministic polynomial-time CNF-SAT candidate lost by the encoded image, an encoded candidate admitted despite not being a lawful uniform polynomial-time candidate, an encoding artifact that changes SAT-failure status, a mismatch between \(\FailsSAT\) and ordinary failure to decide arbitrary finite CNF-SAT, or a mismatch between \(\CnfSATP\) and the existence of an encoded non-failing polynomial-time candidate. Merely objecting that the proof uses an encoded candidate model does not meet the burden.
\end{proposition}

\begin{proof}
Every standard formulation of \(\Pclass\) quantifies over effective descriptions of uniform deterministic procedures together with polynomial-time bounds. The standard encoded model of Definition~\ref{def:standard-encoded-cnf-model} makes that ordinary candidate image explicit. If standard-model instantiation, coverage, soundness, failure adequacy, and image exactness are preserved, the encoded image is neither smaller nor larger than the deterministic polynomial-time CNF-SAT candidate space relevant to \(\CnfSATP\). Therefore a standing objection must identify a concrete defect in one of those clauses rather than treat encoding itself as a defect.
\end{proof}

\subsection{Raw bivalence is bare-level bivalence}
Raw bivalence is stated at the bare layer:
\[
  \PosRaw\vee\SepBare.
\]
The final proof does not infer reportability from the raw disjunction. It uses the context-bound model adequacy packet of Subsection~\ref{subsec:encoded-candidate-model-adequacy} and the SAT-local lower-bound/no-positive bridge
\[
  \SepBare\Rightarrow\StdLB(model)\Rightarrow\CandImageExcl(model)
\]
to move the raw negative branch first into candidate endpoint-image exclusion. No reportability theorem is used in the bare-branch exclusion. Report acts reappear only after the positive endpoint has been derived.


\section{Self-Contained Kernel Proof for Non-Degenerate Construction}\label{sec:self-contained-kernel}

This section internalizes the kernel proof in the manuscript. The kernel is not cited as an external AASC preference and is not treated as an appendix assumption. The definitions below begin with raw step-sequences and kernel-neutral target adequacy, prove that the roles of reference, standing, irreversibility, and admissibility are forced by determinate same-domain construction, prove non-derivability and no-faithful-lower-generator results, and derive the fixed-domain consequence package used later by the SAT endpoint argument.

\subsection{Fixed-Domain Construction and Minimal Setting}

This section records the minimal structural conditions under which a raw step-sequence is assessable as construction of a determinate object on a fixed domain. They are not introduced as optional definitions for a preferred formalism. They are fixed here because without them object identity, admissible stephood, and the distinction between continuation and repair are not stably assessable at all.

\begin{definition}[Raw step-sequence]\label{def:raw-step-sequence}
A raw step-sequence on a domain $D$ is an ordered sequence of inscriptions, rewrites, transitions, or operations on $D$. As such it carries no automatic claim to determinate target identity, admissible stephood, or preservation of object identity across the sequence.
\end{definition}


\begin{definition}[Kernel-neutral target adequacy]\label{def:target-adequacy}
A raw-sequence regime is target-adequate for the present inquiry iff some of its raw sequences are assessable under the following four requirements, stated without yet using the kernel as a definition of construction:
\begin{enumerate}[label=(A\arabic*)]
\item \emph{Target determinacy:} the sequence has a repeatable object of evaluation rather than merely a presentation-dependent trace.
\item \emph{Step evaluability:} its transitions are assessable as advancing, failing to advance, or departing from that same target rather than merely occurring after one another.
\item \emph{Act-time finality:} a later intervention that changes the status of an earlier performance yields a distinct act or distinct sequence rather than retroactively making the same performance admissible.
\item \emph{Same-regime fidelity:} identity-preserving redescription, encoding, or translation preserves target, step-status, and act-time classification; otherwise the presentation has changed regime or changed the object of evaluation.
\end{enumerate}
These adequacy requirements are neutral with respect to the later vocabulary. They specify what must be preserved if a raw trace is to be assessed as construction of a determinate same-domain object at all. They are not proprietary terms of the present framework. They are the minimal commitments already incurred by any opponent who claims to have a counterexample to the paper's target: there must be something fixed, something done toward it, a stable distinction between the original act and later correction, and a preservation condition under redescription. A regime that denies any of these does not refute the target; it declines to instantiate it. Any regime that denies one or more of these conditions while still claiming to evaluate constructions of determinate same-domain objects has changed the target of inquiry rather than supplied a counterexample within it.
\end{definition}

\begin{lemma}[Target adequacy forces the kernel roles]\label{lem:target-adequacy-forces-kernel}
Any target-adequate raw-sequence regime already requires the governance roles later named reference, standing, irreversibility, and admissibility.
\end{lemma}

\begin{proof}
Target determinacy supplies the role later called reference: without a fixed target, there is no object whose construction could be assessed. Step evaluability supplies the role later called standing: without licit or illicit step-status, a transition is only a transition, not a construction step of the same target. Act-time finality supplies the role later called irreversibility: if later normalization can preserve the same act while changing its status, the difference between continuation and repair is not fixed at the time of performance. Finally, same-regime fidelity requires a boundary that holds those classifications together under identity-preserving redescription; that boundary is the role later called admissibility. Thus the kernel terms are introduced as names for forced roles, not as stipulative entry conditions.
\end{proof}

\begin{remark}[Kernel terms as forced role names]
The kernel terms are introduced here only as names for roles whose necessity has already been fixed by the adequacy conditions. The downstream AASC machinery records consequences of this necessity rather than adding independent structure to it.
\end{remark}

\begin{lemma}[Kernel concession lemma]\label{lem:kernel-concession}
Any assessment that grants determinate target identity, evaluable stephood, act-time distinction between continuation and repair, and same-regime fidelity has already granted the kernel roles, whether or not it uses the vocabulary of AASC. Conversely, any assessment that denies one of those requirements no longer assesses a theorem-bearing construction of the same determinate endpoint. Therefore there is no third position from which one may reject the kernel while preserving ordinary theorem-bearing endpoint use for the same target.
\end{lemma}

\begin{proof}
By Lemma~\ref{lem:target-adequacy-forces-kernel}, the four target-adequacy requirements force the roles later named reference, standing, irreversibility, and admissibility. Granting those requirements therefore grants the corresponding governance roles, independently of the terminology used to name them. If an assessment denies one of the four requirements, then by Definition~\ref{def:target-adequacy} it no longer preserves the conditions under which a raw trace is assessable as construction of a determinate same-domain object. Such an assessment changes or exits the target class rather than preserving theorem-bearing endpoint use for the same target. The two cases are exhaustive: the requirements are either granted or denied.
\end{proof}

\begin{definition}[Act identity]\label{def:act-identity}
Two performances count as the same act only if they agree on the admissibility-relevant invariant bundle, fix the same target on the same fixed domain, and have the same admissibility status at act-time. Any alteration that changes target, repairs inadmissibility, or reclassifies the admissibility status of the original performance yields a numerically distinct act or a distinct construction sequence.
\end{definition}

\begin{definition}[Governed construction and non-degenerate construction]\label{def:governed-construction}
Let $X$ be a target object on a fixed domain $D$. A governed construction of $X$, written $\Constr(X)$, is a raw step-sequence on $D$ that preserves, at each step, determinate reference to $X$, standing of the step as an admissible construction step, and irreversibility of inadmissible steps. A governed construction is \emph{non-degenerate} iff those three conditions hold together on the same fixed domain. A regime is \emph{degenerate relative to the present target} iff it may still compute, transform, symbolize, or produce outputs but fails to preserve target identity, admissible stephood, act-time finality, or redescription fidelity as conditions of determinate same-domain construction.

This definition is classificatory rather than generative: it does not assume that all raw step-sequences already count as governed constructions, but distinguishes exactly those sequences for which admissibility-assessable construction is well-defined on the fixed domain.

Admissibility is not a fifth ingredient placed beside reference, standing, and irreversibility; it names the governance boundary under which those three conditions are jointly enforceable on the same sequence.
\end{definition}

\medskip
\noindent\textbf{Meaningful regime.} A regime is meaningful for the present target iff it is target-adequate in the sense of Definition~\ref{def:target-adequacy} and some of its raw step-sequences are assessable as governed constructions of determinate objects rather than as bare symbol manipulations, encodings, or bookkeeping traces. This is not a stipulative filter. It is the least class in which object identity, admissible stephood, and act-time distinction between continuation and repair are available for assessment at all.

\medskip
\noindent\textbf{Fixed domain.} A fixed domain is the domain on which target identity, identity-preserving transformation, and admissibility-relevant status are all evaluated. Operationally, two constructions remain on the same fixed domain iff there exists a translation between them that preserves (i) object identity under the admissibility-relevant invariant bundle, (ii) the admissibility classification invariant distinguishing admissible continuation from prohibited repair, and (iii) the class of governed constructions up to governance equivalence. A purported continuation remains on the same fixed domain only when these preservation conditions are met.

\medskip
\noindent\textbf{Why this target class is forced.} The target class is not selected because it already contains the conclusion. It is the minimal class in which outputs are evaluable as determinate objects rather than presentation-dependent artifacts. By Lemma~\ref{lem:target-adequacy-forces-kernel}, the roles later named reference, standing, irreversibility, and admissibility are forced by target adequacy before governed construction is defined. If identity-preserving transformation is weakened, object identity becomes redescription-relative; if standing is weakened, there is no licit stephood rather than mere inscription; if irreversibility is weakened, the distinction between admissible continuation and retrospective rescue is unstable at act-time. A regime lacking these conditions may still support transformation, search, or computation, but it does not instantiate a weaker version of the same target phenomenon. The class is therefore fixed by evaluability conditions, not by convenience or stipulation.

\medskip
\noindent\textbf{Adequacy note.} The question is not which extra norms we prefer to impose on an already functioning calculus. The question is what the weakest conditions are under which a raw step-sequence is assessable as construction of a determinate object at all. The next definitions and lemmas argue that stable reference, standing, and irreversibility are forced by that adequacy demand rather than borrowed from a thicker prior theory.

\medskip
\noindent\textbf{Scope restriction.} All universal claims below are restricted to meaningful regimes in this adequacy sense. Systems lacking the adequacy conditions uncovered below are not counterexamples to the present thesis; they are different kinds of formal process because they lack the resources required for governed construction of a determinate same-domain object.

\begin{definition}[Derivation]\label{def:derivation}
When the target object $X$ is a proposition or proof conclusion, the corresponding proposition-targeted special case of governed construction is called a derivation and is written $\Der(X)$.
\end{definition}

\begin{definition}[Governance equivalence]\label{def:governance-equivalence}
Two governance conditions $P$ and $Q$ are governance-equivalent, written $P \law Q$, iff they perform the same governance work for identity-preserving construction on the same fixed domain; that is, replacing $P$ by $Q$ leaves the class of governed constructions unchanged.
\end{definition}

\begin{remark}[Equivalence as identity-level invariance]
For the present target, governance-equivalence is not a loose behavioral similarity relation. Once admissibility-assessable, identity-preserving construction is fixed, admissibility classification is bivalent on acts and standing-bearing continuation is fixed by the same admissibility-relevant invariant bundle. Accordingly, any purported alternative governance structure that differs on any act, or on the admissibility-preserving continuation of any act, is not governance-equivalent on the fixed domain. Conversely, where no such divergence exists, the purported alternative contributes no distinct governance content for the target phenomenon.

The companion equivalence and conservation closures sharpen this point further: admissibility and standing do not form two merely coextensive tracks but collapse to the same standing-bearing classification, and conservation is not an auxiliary constraint but the invariant form of that classification. In that sense, equivalence collapses to identity at the level that matters here: any same-domain invariant that agrees with admissibility and standing on all acts and admissibility-preserving continuations is not merely extensionally similar but identical in governance content for admissibility-assessable construction.
\end{remark}

\begin{definition}[Admissibility-relevant invariant bundle]\label{def:invariant-bundle}
An admissibility-relevant invariant bundle is any family of invariants sufficient to determine when two construction acts on a fixed domain preserve the same target, the same standing-bearing status, and the same admissibility boundary. Two acts or presentations count as belonging to the same fixed domain only when they agree on this invariant bundle up to governance equivalence.
\end{definition}

\medskip
\noindent\textbf{Independent governance work.} A condition performs independent governance work on a fixed domain iff its addition changes the set of governed constructions on that domain.

\medskip
\noindent\textbf{Same fixed domain.} References below to ``the same meaningful regime'' are governed by preservation, not by mere reuse of vocabulary. Two presentations, encodings, or meta-systems count as the same fixed domain for the present target only when translation between them preserves determinate target, standing-bearing step status, the admissibility boundary between admissible continuation and prohibited repair, and the class of governed constructions up to governance equivalence. Failure of any one of these preservation conditions constitutes a change of domain rather than a redescription of the same one.

\begin{definition}[Derivable from below]\label{def:derivable-from-below}
A governance condition $P$ is derivable from below iff there exists a governed construction of $P$ whose admissibility can be certified using only governance conditions strictly independent of $P$, does not already presuppose $P$ or any governance-equivalent condition, and does not rely on any condition whose own governance validity already requires $P$. In particular, approximation, convergence, limit, or emergent characterization of $P$ counts as derivation from below only if the admissibility of the approximating, convergent, or emergent process itself does not presuppose $P$ or any governance-equivalent condition.
\end{definition}

\noindent The phrase ``governed construction'' in this definition is not meant to smuggle in the target conclusion. It marks the minimal condition under which the alleged lower derivation counts as a derivation rather than an unlicensed trace. A purported lower derivation may begin from any formal substrate it likes, but it must eventually explain how its own sequence obtains constructional license without using the role it claims to generate.

\begin{definition}[Governance generation versus definability]\label{def:generation-definability}
A condition is \emph{defined} in a system when the system can write down, encode, or characterize it. A condition is \emph{generated} in the present sense only when the system can produce that condition as new governance work from a lower basis. Definability without new governance work is not generation from below.
\end{definition}

\begin{definition}[Externally proposed construction or derivation system]\label{def:kernel-respecting}
Let $\Sigma$ be any purported construction system, derivation system, meta-theory, logical framework, encoding, or external presentation under which a statement $P$ is said to be derivable. We say that $\Sigma$ is \emph{kernel-respecting with respect to $P$} iff the admissibility conditions under which sequences in $\Sigma$ count as governed constructions already enforce $P$, some governance-equivalent condition, or some condition whose own admissibility already requires $P$. This definition classifies the admissibility conditions of $\Sigma$; whether a non-kernel-respecting system can still count as a system of non-degenerate construction for the present target is a theorem proved below, not a stipulation built into the definition.
\end{definition}

\begin{definition}[Kernel of non-degenerate construction]\label{def:kernel}
The kernel of non-degenerate construction is the governance-role set
\[
K := \{\Adm, \St, \RefK, \Irr\},
\]
where $\Adm$ abbreviates admissibility, $\St$ standing, $\RefK$ reference, and $\Irr$ irreversibility. The minimality and uniqueness claims below concern governance roles and governance work on the fixed domain, not uniqueness of one four-part bookkeeping decomposition.
\end{definition}

\begin{remark}
The point of the kernel definition is not to present four independent axioms, nor to shield four merely stipulated starting points from scrutiny. It is to name the lowest invariant cluster that any genuine governed construction already presupposes. Theorems below show that the members of $K$ are neither optional nor independently eliminable.
\end{remark}


\medskip
\noindent\textbf{Why standing is not redundant with admissibility.} Admissibility marks gate status: whether an act may enter licensed construction. Standing marks act-carried eligibility for later licensed use, continuation, or commitment. The term ``standing'' is used because the relevant property is not merely that an act has passed a gate, but that it can stand as a source for subsequent licensed construction. The two coincide extensionally only after fixed-domain closure is proved in Theorem~\ref{thm:standing-reuse}; they are not introduced as the same role. Before that closure theorem, admissibility answers whether the act is admitted, while standing answers whether the admitted act can function as a standing-bearing source for subsequent construction. The later collapse is therefore a result, not a definition.

\begin{definition}[Same-regime faithful counterexample]\label{def:same-regime-faithful-counterexample}
A same-regime faithful counterexample to the kernel is a target-adequate raw-sequence regime together with a class of candidate constructions such that all of the following are preserved on the same fixed domain: determinate target identity, step evaluability, act-time finality, identity-preserving redescription, and the class of standing-bearing continuations. The counterexample must nevertheless omit at least one member of $K$, and not merely replace it by a governance-equivalent role. Any example that loses one of the preservation conditions is a different formal process; any example that restores the missing work by an auxiliary selector, bridge, ranking, tolerance, or repair channel imports fresh governance rather than omitting the role.
\end{definition}

\begin{proposition}[Burden of a successful objection]\label{prop:burden-of-objection}
To refute the kernel result for the present target, one must do at least one of two things: derive some $P\in K$ from below in the sense of Definition~\ref{def:derivable-from-below}, or exhibit a same-regime faithful counterexample in the sense of Definition~\ref{def:same-regime-faithful-counterexample}. Pointing to a different formal practice, a reversible search regime, a bookkeeping encoding, or a system with weaker target identity does not meet this burden.
\end{proposition}

\begin{proof}
The target is fixed by target adequacy and same-domain fidelity. If an objection claims that a kernel role is generated rather than presupposed, then it must satisfy the strict lower-generation condition; otherwise the alleged derivation already uses the role it claims to produce. If the objection instead claims that the role is unnecessary, then it must preserve the same target phenomenon while omitting that role; otherwise it has changed the object of evaluation. Within the present target, a substantive challenge therefore alleges either lower generation or same-regime dispensability. The former is governed by derivability from below; the latter is governed by same-regime faithful counterexample. Anything else changes subject, imports structure, or supplies only notation.
\end{proof}

\noindent This burden is not an artificial raising of the evidential standard. It follows from target preservation. A counterexample to a necessity claim about determinate same-domain construction must preserve determinate same-domain construction. A regime that abandons target identity, step evaluability, act-time classification, or redescription fidelity is not a counterexample to the necessity claim; it is a different target.

\begin{lemma}[Raw sequences are not yet constructions]\label{lem:raw-not-construction}
A raw step-sequence does not count as a construction of a determinate object merely in virtue of being a sequence.
\end{lemma}

\begin{proof}
A raw sequence may be written, executed, or transformed while leaving open what object, if any, it fixes; whether its steps are admissible steps of the same act; and whether later repair preserves or changes the act being assessed. Those are exactly the governance matters at issue. So raw sequentiality alone yields, at most, a trace. It does not yet yield construction of a determinate same-domain object.
\end{proof}

\begin{lemma}[Same-act repair is impossible]\label{lem:same-act-repair-impossible}
If a later intervention changes an inadmissible performance into an admissible one, then the later intervention does not preserve the same act.
\end{lemma}

\begin{proof}
By Definition~\ref{def:act-identity}, act identity requires preservation of the admissibility-relevant invariant bundle together with preservation of act-time admissibility status. A later intervention that converts inadmissibility into admissibility changes that status and therefore fails to preserve act identity. What results is either a distinct act or a distinct construction sequence, not repair of the same act.
\end{proof}

\begin{lemma}[Reference is necessary for determinate construction output]\label{lem:reference-necessity}
If a purported construction act lacks determinate reference, then it does not yield a determinate object of the regime.
\end{lemma}

\begin{proof}
A construction must be about or directed toward a target in a fixed way in order to preserve identity through its steps. If the act's expressions refer equivocally or vacuously, then no single governed target is fixed: different readings, or none at all, remain equally compatible with the same mark sequence. What remains is a trace, not a determinate object. Hence determinate construction output requires determinate reference.
\end{proof}

\begin{lemma}[Standing is necessary for admissible stephood]\label{lem:standing-necessity}
If a purported step lacks standing, it does not function as an admissible construction step, but only as an unsupported inscription or transition.
\end{lemma}

\begin{proof}
A construction step is not merely something that occurs between two inscriptions. It purports to advance a target under some licensed governance status. Without standing, there is no governed construction act whose role could be assessed. The item may still be written, computed, or recorded, but it does not count as an admissible step. Therefore standing is necessary for constructional stephood.
\end{proof}

\begin{lemma}[Irreversibility is necessary for stable construction status]\label{lem:irreversibility-necessity}
If invalid steps are repairable post hoc without loss, then the distinction between admissible continuation and retrospective rescue is not stably fixed at act-time.
\end{lemma}

\begin{proof}
Suppose invalidity could later be removed while preserving the same act as fully admissible. By Lemma~\ref{lem:same-act-repair-impossible}, that preservation claim is false: any intervention that changes inadmissibility into admissibility changes the act identity of what is being assessed. So if repair is allowed, the later admissible outcome is not the same act but a distinct act or distinct sequence. The original act therefore remains inadmissible, and any regime that treats later repair as preserving the same act loses a stable act-time distinction between continuation and rescue. Stable constructional status consequently requires irreversibility.
\end{proof}

\begin{lemma}[Admissibility is the governance boundary of joint enforceability]\label{lem:admissibility-joint}
Whenever determinate reference, standing, and irreversibility are jointly enforceable on the same construction acts, admissibility is already in force there as the governance boundary of that joint enforceability.
\end{lemma}

\begin{proof}
By Lemmas~\ref{lem:reference-necessity}--\ref{lem:irreversibility-necessity}, a regime that treats its acts as determinate, admissibility-assessable construction acts must already maintain determinate reference, licit step-status, and stable non-repairability of invalid steps. The point is not that one first has three free-standing ingredients and then adds admissibility as an extra fifth clause. Rather, once those three conditions are jointly enforceable on the same acts as conditions of construction, the regime already has a governance boundary distinguishing admissible from non-admissible acts and distinguishing continuation from retrospective rescue. That governance-boundary status is exactly what is here called admissibility. Conversely, a regime lacking that joint governance boundary does not govern construction as non-degenerate construction in the present sense. Hence admissibility is already in force wherever the three conditions are jointly enforceable.
\end{proof}

\begin{proposition}[Boundary closure is forced by admissibility]\label{prop:boundary-closure-forced}
Let a regime support admissibility-assessable, identity-preserving construction with irreversibility. Then any admissibility-relevant discriminator, selector, or classification rule must be invariant under identity-preserving transformation on the fixed domain.

Consequently, any attempt to introduce additional admissibility-relevant structure not preserved under such transformations either:
\begin{enumerate}[label=(\roman*)]
    \item changes admissibility classification and so introduces a competing gate, or
    \item leaves admissibility classification unchanged and is therefore inert bookkeeping.
\end{enumerate}
Hence the admissibility classification invariant is closed under same-domain construction: no additional admissibility-relevant discriminator can be introduced without either collapsing to governance-equivalence or changing the domain of evaluation.
\end{proposition}

\begin{proof}
If an admissibility-relevant discriminator is not invariant under identity-preserving transformation, then two presentations of what purports to be the same target on the same fixed domain receive different admissibility significance. If that difference changes what counts as admissible continuation, standing-bearing reuse, or prohibited repair, then the regime has introduced a second admissibility classifier for the same fixed-domain acts. But the distinction between admissible continuation and prohibited repair is exactly the governance work fixed by admissibility together with irreversibility; reclassifying same-domain acts therefore yields a competing gate rather than a mere refinement. If, by contrast, the additional discriminator does not change any admissibility verdict, then it contributes no independent governance work and is only bookkeeping. There is no third possibility: a same-domain addition either changes admissibility significance or it does not. Accordingly, admissibility-assessable, identity-preserving construction with irreversibility already forces closure of the admissibility classification invariant against further same-domain discriminators.
\end{proof}

\begin{theorem}[Minimal assessability boundary]\label{thm:minimal-assessability-boundary}
Let $R$ be any system whose outputs are treated as determinate objects and whose transitions are evaluated as construction steps. Then determinate reference, admissible stephood, and non-repairable invalidity are jointly necessary and minimal for that evaluability. More exactly:
\begin{enumerate}[label=(\roman*)]
    \item if determinate reference is removed, object identity is lost and the outputs are no longer evaluable as determinate content;
    \item if standing is removed, no transition qualifies as an admissible construction step and the outputs are no longer evaluable as construction steps;
    \item if irreversibility is removed, there is no stable distinction between admissible continuation and retrospective rescue, and the outputs are no longer evaluable as constructions; and
    \item if an additional condition is imposed as an independent governance requirement beyond those three, then the regime is governed more tightly than evaluability itself requires.
\end{enumerate}
Accordingly, these three conditions mark the minimal assessability boundary for identity-preserving construction as such. The minimality claim is a claim about governance work on the fixed domain, not about uniqueness of one bookkeeping decomposition.
\end{theorem}

\begin{proof}
Items (i)--(iii) are exactly the failure modes identified in Lemmas~\ref{lem:reference-necessity}--\ref{lem:irreversibility-necessity}. Without determinate reference there is no fixed object identity and therefore no evaluable content; without standing there is no licensed act and therefore no evaluable construction step; without irreversibility there is no stable admissibility distinction and therefore no evaluable construction rather than retrospective normalization. The point is not merely that these failures are bad outcomes. It is that, once any one of them occurs, the target phenomenon itself is no longer present: what remains is trace-production, transformation, or bookkeeping rather than governed construction of a determinate same-domain object. This also blocks nearby weakenings and substitutes. Any purported weaker replacement for one of the three conditions must either preserve exactly the same governance work, in which case it is merely a notational variant of that condition; or else it must relax some part of that work, in which case one of the same excluded failures reappears---presentation-dependent object identity, indeterminate target, or collapse of the distinction between continuation and repair. If a purported weakening avoids those failures only by adding a new admissibility-relevant ranking, selector, tolerance, or repair-sensitive middle status, then it no longer weakens the package but supplements it by fresh governance structure. For minimality, suppose an additional condition $Q$ is imposed as an independent governance requirement on every admissibility-assessable construction act. Then either $Q$ is already enforced by the joint assessability work of reference, standing, and irreversibility, in which case it is not additional; or else $Q$ excludes some act-sequences that already satisfy those three conditions, in which case $Q$ governs more tightly than evaluability itself requires. There is no third possibility. So the three conditions are jointly necessary and minimal for that evaluability.
\end{proof}

\begin{corollary}[Non-arbitrariness of the target class]\label{cor:non-arbitrariness-target-class}
The class of meaningful regimes is not defined by pre-selecting the kernel. It is exactly the class of systems in which determinate, admissibility-assessable construction is possible for the present target. Restriction to that class is therefore not protective but forced by the target of inquiry.
\end{corollary}

\begin{proof}
By Theorem~\ref{thm:minimal-assessability-boundary}, determinate reference, standing, and irreversibility are not optional entry criteria appended in advance of the argument. They are the minimal conditions under which outputs and transitions are evaluable as construction at all. So to restrict attention to regimes in which those conditions are assessable is simply to restrict attention to the subject matter of the paper rather than to protect the thesis from counterexample by fiat.
\end{proof}

\begin{theorem}[The target class is forced, not stipulated]\label{thm:forced-target-class}
Let $R$ be any regime of raw step-sequences. If $R$ lacks determinate reference, standing, or irreversibility in the senses fixed above, then $R$ does not support governed construction of determinate same-domain objects. It supports only raw transformation, encoding, or bookkeeping processes.
\end{theorem}

\begin{proof}
If determinate reference is lacking, Lemma~\ref{lem:reference-necessity} shows that no determinate target object is fixed. If standing is lacking, Lemma~\ref{lem:standing-necessity} shows that no step counts as an admissible construction step. If irreversibility is lacking, Lemma~\ref{lem:irreversibility-necessity} together with Lemma~\ref{lem:same-act-repair-impossible} shows that act-time distinction between continuation and repair collapses. By Lemma~\ref{lem:raw-not-construction}, raw sequentiality does not by itself restore those losses. So a regime lacking any one of these conditions may still execute sequences, but it does not instantiate governed construction of determinate same-domain objects.
\end{proof}

\begin{lemma}[Exhaustion of failure modes for non-kernel regimes]\label{lem:failure-mode-exhaustion}
Let $R$ be a regime of raw step-sequences. Suppose $R$ lacks at least one member of $K$ while claiming to support determinate, identity-preserving construction on a fixed domain. Then exactly one of the following holds:
\begin{enumerate}[label=(\roman*)]
\item object identity is not fixed, because determinate reference fails;
\item step status is not evaluable as admissible or inadmissible, because standing fails;
\item the distinction between continuation and repair is not stable at act-time, because irreversibility fails; or
\item the regime restores one or more of these distinctions only by adding additional governance-bearing structure not fixed on the original domain, and so proceeds by illicit import rather than same-domain construction.
\end{enumerate}
These cases cover the preservation failures relevant to the present target.
\end{lemma}

\begin{proof}
If $R$ lacks reference, then by Lemma~\ref{lem:reference-necessity} no determinate target is fixed, giving (i). If $R$ lacks standing, then by Lemma~\ref{lem:standing-necessity} no step counts as an admissible construction step, giving (ii). If $R$ lacks irreversibility, then by Lemmas~\ref{lem:same-act-repair-impossible} and \ref{lem:irreversibility-necessity} the distinction between admissible continuation and retrospective rescue is unstable at act-time, giving (iii). The only remaining possibility is that $R$ attempts to preserve determinate identity, admissible stephood, or stable act-status while omitting one of the kernel roles. But then those distinctions are being reintroduced by fresh governance-bearing structure not fixed on the original domain. That is illicit import rather than a same-domain realization below the kernel, giving (iv). Within the present target, every purported non-kernel regime either loses one of the three assessability goods directly or restores it by extra same-domain authority.
\end{proof}

\begin{lemma}[Reduction to kernel failure modes]\label{lem:reduction-kernel-failure}
Let $R$ be any regime that deviates from the kernel while claiming to preserve determinate, identity-preserving construction on a fixed domain. Then every such deviation induces at least one of the following:
\begin{enumerate}[label=(\roman*)]
\item non-uniqueness of target under identity-preserving transformation;
\item indeterminacy of admissible stephood;
\item instability of admissibility under continuation; or
\item introduction of a selector, ranking, or auxiliary structure not invariant under identity-preserving transformation.
\end{enumerate}
Hence every deviation reduces to a kernel failure mode.
\end{lemma}

\begin{proof}
If a deviation changes which target is fixed under identity-preserving transformation, then same-target identity is not preserved, giving (i). If it leaves target fixed but no longer preserves evaluable admissible stephood, then licensed step status becomes indeterminate, giving (ii). If it preserves target and stephood while allowing admissibility status to vary under purported same-act continuation, then the act-time distinction between continuation and repair is unstable, giving (iii). The only remaining possibility is that the regime preserves these goods by introducing some selector, ranking, tolerance class, or other auxiliary governance-bearing discriminator not already fixed by the original admissibility-relevant invariant bundle. But such a discriminator is fresh same-domain authority rather than preservation of the original regime, giving (iv). Since every deviation changes target, step-status, continuation-status, or the authority by which these are fixed, the partition covers the relevant ways a same-domain deviation can proceed. Therefore every deviation reduces to a kernel failure mode.
\end{proof}

\begin{theorem}[Non-degenerate construction forces the kernel]\label{thm:construction-forces-kernel}
If a meaningful regime admits non-degenerate construction in the sense of Definition~\ref{def:governed-construction}, then the full kernel $K$ is already in force in that regime. Equivalently, no meaningful regime supports non-degenerate construction in that sense while lacking admissibility, standing, reference, irreversibility, or some governance-equivalent replacement for one of them.
\end{theorem}

\begin{proof}
Let $R$ be a meaningful regime admitting non-degenerate construction in the sense of Definition~\ref{def:governed-construction}. By Theorem~\ref{thm:minimal-assessability-boundary}, any regime in which outputs are assessable as determinate, admissibility-assessable constructions already requires determinate reference, admissible stephood, and irreversibility as the minimal assessability boundary of that evaluability. By Lemma~\ref{lem:admissibility-joint}, the joint enforceability of those conditions already places admissibility in force as the governance boundary constituted by their coherence on the same acts. Hence $R$ already enforces $\RefK$, $\St$, $\Irr$, and $\Adm$, that is, the full kernel $K$.
\end{proof}

\begin{corollary}[No non-degenerate regime below the kernel]\label{cor:no-regime-below-kernel}
No meaningful regime can both support non-degenerate construction in the sense of Definition~\ref{def:governed-construction} and generate any member of $K$ from a strictly lower governance-free basis while remaining within that same fixed domain.
\end{corollary}

\begin{proof}
Assume some meaningful regime $R$ supports non-degenerate construction in the sense of Definition~\ref{def:governed-construction} and nevertheless generates some $P \in K$ from a strictly lower governance-free basis. By Theorem~\ref{thm:construction-forces-kernel}, the full kernel $K$ is already in force in $R$ before any such generation begins. So the alleged generator does not produce $P$ from below; it operates only within a regime whose meaningfulness already depends on $P$. That contradicts the claim that $P$ was first generated from a lower basis.
\end{proof}

\begin{remark}
The point of Theorem~\ref{thm:construction-forces-kernel} is not to build the package into construction by stipulation. It is to show that, for the target domain fixed in the fixed-domain construction subsection, any regime that already supports non-degenerate construction already instantiates the kernel by necessity. The derivation lemmas of Section~3 therefore record a local reflection of a prior construction-level fact.
\end{remark}

\subsection{Presupposition Structure of Derivation}

With Theorem~\ref{thm:construction-forces-kernel} in place, derivation is the proposition-targeted special case of licensed construction. The present section records the derivation-level presupposition facts. These lemmas do not generate the package. They state what the fixed-domain construction subsection result looks like once the focus narrows from non-degenerate construction in general to derivation in particular: a sequence counts here as a derivation only if it already operates with admissibility, standing, reference, and irreversibility in place.

\begin{lemma}[Derivation presupposes admissibility]
If $\Der(P)$ holds for any proposition $P$, then admissibility is already in force.
\end{lemma}

\begin{proof}
By Definition~\ref{def:derivation}, a governed derivation is not merely a formal trace but a governed construction act whose target is a proposition. Its steps therefore count as steps only under an admissibility regime. So if $\Der(P)$ exists, admissibility has already been presupposed.
\end{proof}

\begin{lemma}[Derivation presupposes standing]
If $\Der(P)$ holds, then standing is already in force.
\end{lemma}

\begin{proof}
A derivation with no standing would be a sequence of acts with no admissible target-bearing status. Such a sequence might be written down, but it would not count as a derivation of anything. Licensed derivation therefore presupposes standing.
\end{proof}

\begin{lemma}[Derivation presupposes reference]
If $\Der(P)$ holds, then determinate reference is already in force.
\end{lemma}

\begin{proof}
A derivation whose expressions fail to refer determinately does not derive a determinate conclusion. It supplies marks, not governed content. Since $\Der(P)$ names a derivation of a proposition rather than a mere syntactic trail, reference is presupposed.
\end{proof}

\begin{lemma}[Derivation presupposes irreversibility]
If $\Der(P)$ holds, then irreversibility is already in force.
\end{lemma}

\begin{proof}
Without irreversibility, invalid steps could be repaired retroactively into acceptable ones, and the distinction between admissible continuation and post hoc rescue would collapse. In that case the derivation would not possess a stable governance status at all. Hence governed derivation presupposes irreversibility.
\end{proof}

\begin{corollary}[Global presupposition lemma]\label{cor:global-presupposition}
For every proposition $P$, if $\Der(P)$ holds, then the full kernel $K$ is already in force.
\end{corollary}

\begin{proof}
Combine the four previous lemmas.
\end{proof}

\begin{theorem}[Definability is not governance generation]\label{thm:definability-not-generation}
Let $P$ be any governance condition. A system may define, encode, or represent $P$ without generating $P$ from below. Generation from below occurs only if the system contributes new governance work sufficient to produce $P$ for the original acts on the original fixed domain.
\end{theorem}

\begin{proof}
A definition or encoding may restate a governance condition, give it a symbol, embed it in a meta-language, or characterize its extension. None of those activities by itself changes the class of governed constructions on the original fixed domain. By definition, governance generation requires independent governance work. So definability, representation, and coding are not by themselves generation from below.
\end{proof}

\begin{remark}
Accordingly, the non-derivability claims below are not claims that one cannot write down or characterize admissibility, standing, reference, or irreversibility. They are claims that no lower same-domain basis generates the governance work those conditions perform for the original acts.
\end{remark}

\subsection{Non-Derivability of the Kernel}

The results of this section have two directions. Internally, they show that no lower same-domain basis generates the governance work performed by the kernel. Externally, they identify the construction-level form of a broader meta-necessity: in any non-degenerate irreversible compositional regime, the invariant required to make irreversibility well-defined is exactly the invariant characterized here under the fixed-domain conditions of identity-preserving construction. The kernel is therefore not an optional refinement layered on top of admissibility. It is the minimal decomposition of the admissibility invariant once the target object is fixed.

The argumentative standard in this section is deliberately adversarial. A proposed defeat of the kernel must either supply a faithful lower generator for one of its roles, or supply a same-regime faithful counterexample that preserves the target phenomenon while lacking that role. The former fails because any licensed construction of the role already relies on the role or a governance-equivalent substitute. The latter fails because any same-regime omission either loses target identity, loses licit stephood, loses act-time irreversibility, imports a replacement selector, or collapses extensionally back to the same governance work.

\begin{theorem}[Non-derivability of admissibility]
Admissibility is not derivable from below.
\end{theorem}

\begin{proof}
Assume for contradiction that admissibility is derivable from below. Then there exists a governed construction of $\Adm$ whose admissibility does not already presuppose $\Adm$ or a governance-equivalent condition. But by the previous section, any licensed construction already presupposes admissibility. So the alleged construction uses what it claims to produce. This is precisely the forbidden circularity in the definition of derivability from below. Therefore admissibility is not derivable from below.
\end{proof}

\begin{theorem}[Non-derivability of standing]
Standing is not derivable from below.
\end{theorem}

\begin{proof}
Suppose standing were derivable from below. Then there would exist a governed construction of $\St$ whose admissibility did not already presuppose standing. By the previous section, however, the very existence of such a construction already presupposes standing. The construction therefore fails the ``from below'' condition by collapsing into the very status it was supposed to generate.
\end{proof}

\begin{theorem}[Non-derivability of reference]
Reference is not derivable from below.
\end{theorem}

\begin{proof}
If reference were derivable from below, there would be a governed construction of $\RefK$ whose admissibility did not already rely on determinate reference. But by the previous section, no licensed construction exists without determinate reference. Hence the putative construction again presupposes what it claims to derive. Reference is therefore not derivable from below.
\end{proof}

\begin{theorem}[Non-derivability of irreversibility]
Irreversibility is not derivable from below.
\end{theorem}

\begin{proof}
Assume irreversibility is derivable from below. Then a governed construction of $\Irr$ exists whose admissibility does not already presuppose irreversibility. The previous section rules this out: a licensed construction already relies on the stable difference between admissible continuation and invalid repair. So irreversibility cannot be derived from beneath itself.
\end{proof}

\begin{corollary}[Kernel non-derivability]\label{cor:kernel-non-derivability}
No member of $K$ is derivable from below.
\end{corollary}

\begin{proof}
Immediate from the four preceding theorems.
\end{proof}

\begin{theorem}[No faithful lower generator / structural non-derivability]\label{thm:kernel-no-lower-generator}
Let $P \in K$. No system supporting non-degenerate construction can generate $P$ from a genuinely lower condition while remaining within the same fixed domain. Equivalently, any such system already has $P$ in force before internal construction begins.
\end{theorem}

\begin{proof}
Assume for contradiction that some system or meaningful regime $R$ supports non-degenerate construction in the sense of Definition~\ref{def:governed-construction} and nevertheless generates some $P \in K$ from a purportedly lower package. There are only two cases.

First, the alleged generator operates on the same fixed domain as the original acts. Then by Theorem~\ref{thm:construction-forces-kernel}, the full package $K$ is already in force in $R$, and hence $P$ is already in force there. So the alleged lower package does not generate $P$ from below; it operates only inside a regime whose meaningfulness already depends on $P$. In this same-domain case the remaining escape routes are exhausted. A purported lower generator must either (i) introduce additional same-domain structure not already licensed on the fixed domain, (ii) collapse to triviality by adding no independent governance work, (iii) alter identity-preserving transformation, reference conditions, or act-identity conditions so that the original target phenomenon is no longer preserved, or (iv) reproduce the original admissibility boundary extensionally under a different presentation. Route (i) is illicit same-domain import rather than lower generation. Route (ii) is not generation at all. Route (iii) changes domain rather than generating the kernel within the same one. Route (iv) is extensional redescription of what was already in force. Any further same-domain route would have to do independent governance work while avoiding all four branches, which is precisely what the preceding results exclude.

Second, the alleged generator operates outside the original fixed domain and then attempts to re-enter it. But then the resulting construction does not govern the original acts on their original fixed domain without an additional transfer principle, selector, or authority connecting the external regime back to the old one. That additional bridge is exactly fresh structure not fixed by the original domain. More sharply: the admissibility boundary together with standing classification fixes the identity of the domain under evaluation. So any purported derivation that changes that boundary or that classification has not generated the kernel for the original domain; it has changed the object of evaluation and is therefore a scope change rather than a lower derivation. The maneuver is consequently scope change or illicit import rather than lower generation of the kernel for the original acts.

Thus neither same-domain nor cross-domain generation succeeds. Re-describing the alleged generator as belonging to a regime that lacks $P$ therefore does not help, because either it fails to govern the original domain or, if it does govern it, $P$ was already in force there. The obstruction is structural rather than a peculiarity of one proof format.
\end{proof}

\begin{theorem}[Regime-invariant non-derivability]\label{thm:regime-invariant}
Let $P \in K$. No system $\Sigma$ can derive $P$ from below unless the admissibility conditions under which a $\Sigma$-construction counts as a construction already enforce $P$, something governance-equivalent to $P$, or a condition that already requires $P$.
\end{theorem}

\begin{proof}
Let $P \in K$, and suppose some system $\Sigma$ is offered as a construction or derivation system for $P$. There are two cases.

First, the admissibility conditions under which a $\Sigma$-construction counts as a licensed construction already enforce $P$, a governance-equivalent condition, or a condition whose own admissibility already requires $P$. Then $\Sigma$ is kernel-respecting with respect to $P$ by Definition~\ref{def:kernel-respecting}, and the alleged construction does not generate $P$ from below.

Second, those conditions do not already enforce $P$ in any of those ways. If $\Sigma$ nevertheless supported non-degenerate construction in the same fixed domain, Theorem~\ref{thm:construction-forces-kernel} would force the whole package, including $P$, already to be in force there. So a non-kernel-respecting $\Sigma$ cannot both lack $P$ and still count as a system of non-degenerate construction for the present target. Its output is therefore, at best, an external coding, a bookkeeping device, or a purely formal trace; it is not a construction of $P$ from below for the target phenomenon at issue.

Changing formal dress, moving to a meta-theory, or recoding the language therefore does not remove the need for the work that $P$ names. By Theorem~\ref{thm:definability-not-generation}, such recodings may define or represent $P$ without generating it from below.
\end{proof}

\begin{theorem}[No faithful lower generator and no faithful same-domain extension]\label{thm:no-faithful-same-domain-extension}
For the admissibility boundary on a fixed domain, there is no faithful lower generator and no faithful same-domain extension that converts admissibility, standing, reference, or irreversibility into lower theorems for the original acts on the original fixed domain. Any purported attempt ends in exactly one of four outcomes: illicit structure import, triviality, violation of the base goods, or extensional collapse back to the gate itself.
\end{theorem}

\begin{proof}
Any genuinely pre-boundary constraining device must be available and meaningful on the fixed pre-admissible side. If it is not, it is illicit import. If it is, then under the same-domain invariance conditions relevant here it either has no independent constraining force on individual acts, destroys reference, identity, or irreversibility, or coincides extensionally with the gate already being determined. These four branches cover the ways a purported pre-boundary device could relate to the original acts. A further case would have to contribute independent same-domain governance work while neither importing structure, trivializing the construction, violating the base goods, nor collapsing extensionally to the gate. But independent same-domain governance work is exactly what the non-derivability and no-faithful-same-domain-extension results exclude. The same exhaustion therefore remains in force under faithful same-domain extension, because fresh same-domain authority for the original acts is itself illicit structure import unless it changes nothing or collapses back to the original gate.
\end{proof}

\begin{proposition}[Cross-domain closure and admissibility transport]\label{prop:cross-domain-closure}
The non-derivability results above exclude not only same-domain lower generators but also cross-domain simulation attempts. In particular, let $\Sigma$ be any external system and let $\phi$ be a mapping from $\Sigma$ into the fixed domain of the original acts. If $\phi$ preserves (i) target identity under identity-preserving transformation, (ii) admissibility classification, and (iii) standing-bearing continuation, then $\phi$ is governance-equivalent to the original domain on the relevant image. If $\phi$ fails to preserve any of (i)--(iii), then the attempted transport changes the admissibility-relevant invariant bundle and constitutes a scope change rather than a derivation. Accordingly, cross-domain simulation does not provide a lower generator for admissibility on the fixed domain: it either collapses to governance-equivalence or exits the domain of evaluation.
\end{proposition}

\begin{proof}
The fixed domain is identified by preservation of target identity, admissibility classification, and the class of governed constructions up to governance equivalence (Definition~\ref{def:governed-construction} and the fixed-domain clause of the fixed-domain construction subsection). By Proposition~\ref{prop:boundary-closure-forced}, admissibility-assessable, identity-preserving construction with irreversibility already forces closure of the admissibility classification invariant against further same-domain discriminators. So if a mapping $\phi$ from an external system preserves target identity, admissibility classification, and standing-bearing continuation for the original acts, then it preserves exactly the invariant bundle that individuates the fixed domain for the present target. On the relevant image, $\phi$ therefore contributes no independent governance work beyond what is already counted as the same domain; it is governance-equivalent rather than a lower generator.

If, by contrast, $\phi$ fails to preserve any one of those items, then the transported construction no longer governs the original acts as the same fixed-domain objects with the same admissibility classification and standing-bearing continuation. That is not derivation of admissibility for the original acts, but change of the object under evaluation. Hence cross-domain transport either collapses to governance-equivalence or exits the domain of evaluation.
\end{proof}

\begin{remark}
The previous results do not show that one may never write down statements asserting admissibility, standing, reference, or irreversibility. They show something sharper: these conditions cannot be generated as foundations by a construction that does not already rely on them. Their non-derivability is thus a minimality result, not an expressive prohibition.
\end{remark}
\begin{theorem}[No governance-effective deeper same-domain invariant]\label{thm:no-deeper-invariant}
Let a system support admissibility-assessable, identity-preserving construction with irreversibility. Then no governance-effective invariant strictly beneath admissibility exists for that same fixed-domain target.

More exactly, any purported invariant that claims to ground, constrain, or underwrite admissibility must fall into exactly one of the following cases:
\begin{enumerate}[label=(\roman*)]
\item it imports structure not invariant under identity-preserving transformation and therefore changes domain rather than grounding admissibility on the fixed domain;
\item it carries no constraining power on admissibility-assessable construction and is therefore vacuous with respect to the present target; or
\item it is governance-equivalent to admissibility or standing under renaming and therefore performs no deeper governance work.
\end{enumerate}
Accordingly, no governance-effective deeper same-domain invariant remains compatible with irreversible, identity-preserving construction on the fixed domain.
\end{theorem}

\begin{proof}
Any purported deeper invariant must, by hypothesis, do one of two things: either it constrains admissibility, or it fails to constrain it. If it fails to constrain admissibility, then it does no governance work for the target phenomenon and is vacuous relative to identity-preserving construction. So only the constraining case remains.

Now any invariant that constrains admissibility must be available wherever admissibility is being fixed for the same acts on the same domain. If it is introduced only by adding structure not preserved under identity-preserving transformation, then it does not underwrite admissibility on the fixed domain but changes the admissibility-relevant invariant bundle and thus changes domain. If instead it is presented as same-domain and same-target while avoiding such imported structure, then by the non-derivability and no-faithful-same-domain-extension results above, it cannot contribute independent deeper governance work beneath or prior to admissibility. In that case either it leaves the admissibility classification unchanged, in which event it is inert bookkeeping, or it reproduces the same governance work already performed by admissibility or standing under a different presentation, in which event it is governance-equivalent under renaming.

A further option would have to constrain admissibility on the same fixed domain while neither importing structure, nor being vacuous, nor collapsing to governance-equivalent admissibility work. But that is exactly the kind of independent deeper same-domain invariant excluded by the preceding non-derivability results together with Proposition~\ref{prop:cross-domain-closure} and the fixed-domain-extension result. Therefore no governance-effective invariant strictly beneath admissibility exists for the present fixed-domain target.
\end{proof}

\begin{remark}[Dependency status of invariant impossibility]
Here and throughout, ``deeper invariant'' means any invariant that purports to ground or constrain admissibility without being reducible to it.

The preceding theorem is eliminative, not generative. It does not serve as a premise for the existence or necessity of admissibility. Rather, it constrains the space of possible alternatives once admissibility is already required by irreversibility and identity-preserving construction and once the fixed-domain kernel has been established in the kernel forcing and non-derivability subsections.

In particular, the impossibility of deeper invariants is not used to derive admissibility itself, but only to exclude competing invariant structures that might otherwise be proposed in its place. Any dependence on architectural commitments such as non-transmissive closure or boundary discipline is therefore downstream: within the argumentative order of the present paper, those commitments are consequences of the admissibility structure already established in the kernel forcing and non-derivability subsections or eliminative restatements of what those sections exclude, not independent assumptions on which the kernel results rely.
\end{remark}

\subsection{Mutual Closure and Minimality}

\begin{proposition}[Admissibility requires standing]\label{prop:admissibility-requires-standing}
Any regime lacking standing is not an admissibility regime.
\end{proposition}

\begin{proof}
A regime without standing does not distinguish licit construction from unsupported utterance or transition. It therefore lacks the governance force required to count as admissibility. At best it describes transitions; it does not govern construction.
\end{proof}

\begin{proposition}[Standing requires reference]\label{prop:standing-requires-reference}
Standing without determinate reference is incoherent.
\end{proposition}

\begin{proof}
Standing is the licitness of a construction act. But where reference is indeterminate, there is no determinate target whose licitness could be assessed. A purported standing relation with no reference relation underneath it licenses emptiness or equivocation, not governed construction.
\end{proof}

\begin{proposition}[Governed reference requires irreversibility]\label{prop:reference-requires-irreversibility}
The reference relevant to non-degenerate construction is not bare denotation but reference fixed for governed use at act-time. That form of reference cannot be maintained if invalid construction acts are repairable post hoc without loss.
\end{proposition}

\begin{proof}
A symbol may still bear a denotational relation under later reinterpretation or repair. That is not yet the kernel notion of reference. The issue here is whether the act, as performed, fixes a determinate target for governed constructive use. If later repair can retroactively convert the same invalid act into a fully admissible one without loss, then the constructive status under which the act's target is fixed is not settled at act-time. What remains may be denotation, but not governed reference in the sense required for non-degenerate construction. Hence kernel-reference requires irreversibility of invalidity.
\end{proof}

\begin{proposition}[Governance irreversibility requires admissibility]\label{prop:irreversibility-requires-admissibility}
The irreversibility relevant here is not mere temporal one-wayness or procedural non-rewindability. It is the governance distinction between licensed continuation and prohibited repair, and that distinction is available only under an admissibility boundary.
\end{proposition}

\begin{proof}
One may define other senses of irreversibility---causal, temporal, computational, or procedural---without invoking admissibility. Those are not the senses at issue here. The kernel member $\Irr$ concerns irreversible invalidity of construction acts. Without admissibility there is no principled governance boundary under which one continuation counts as licensed while another counts as illicit repair. One can still speak of sequence or alteration, but not of irreversible invalidity in the governance sense. Thus governance irreversibility presupposes admissibility.
\end{proof}

\begin{theorem}[Mutual closure of the kernel]\label{thm:mutual-closure-kernel}
The members of $K$ form a mutually closed necessity set: none can be coherently asserted while denying the others' governance role.
\end{theorem}

\begin{proof}
Propositions~\ref{prop:admissibility-requires-standing}--\ref{prop:irreversibility-requires-admissibility} exhibit a governance-specific necessity loop: admissibility requires standing, standing requires reference, governed reference requires irreversibility, and governance irreversibility requires admissibility. Hence no member can be detached from the rest while retaining its intended governance content. What remains after detachment is either vacuous, equivocal, or degenerate.
\end{proof}

\begin{remark}
The last two arrows of the loop concern the governance senses fixed in the fixed-domain construction subsection. Proposition~\ref{prop:reference-requires-irreversibility} is not a claim about bare denotation independently of construction; it concerns reference as fixed for governed constructive use at act-time. Proposition~\ref{prop:irreversibility-requires-admissibility} is not a claim about every temporal or procedural notion of irreversibility; it concerns irreversible invalidity of construction acts. In those senses, the loop is closed.
\end{remark}

\begin{theorem}[Minimality of the kernel]\label{thm:minimality-kernel}
The kernel $K$ is minimal for non-degenerate construction: no proper subset of $K$ suffices for non-degenerate construction.
\end{theorem}

\begin{proof}
Assume a proper subset $K_0 \subsetneq K$ suffices for non-degenerate construction. Let $X \in K \setminus K_0$ be an omitted member. Since non-degenerate construction allegedly remains possible, either $X$ is unnecessary or $X$ is generated from the remaining members. The first option contradicts Theorem~\ref{thm:mutual-closure-kernel}, since the governance role of each member depends on the rest. The second option contradicts Corollary~\ref{cor:kernel-non-derivability}, since $X$ would then be derivable from below from a lower basis. Both possibilities fail. Therefore no proper subset suffices.
\end{proof}

\begin{corollary}[No lower generator]\label{cor:kernel-no-lower-generator}
There exists no governance condition or package strictly below $K$ that generates $K$.
\end{corollary}

\begin{proof}
If such a lower generator existed, at least one member of $K$ would be produced from below by a package not already equivalent to that member. That contradicts Corollary~\ref{cor:kernel-non-derivability}.
\end{proof}

\begin{theorem}[Uniqueness of governance work, not of bookkeeping decomposition]\label{thm:governance-work-uniqueness}
If $K'$ is any governance package sufficient for non-degenerate construction in the sense of Definition~\ref{def:governed-construction}, then $K'$ must cover the same governance roles as $K$ and is governance-equivalent to $K$ at that role level. Coarser or finer decompositions may exist; what cannot exist is a governance-non-equivalent alternative package, where governance-non-equivalent means inducing a different class of governed constructions.
\end{theorem}

\begin{proof}
Any package sufficient for non-degenerate construction must, by Theorem~\ref{thm:construction-forces-kernel}, already secure admissibility, standing, reference, and irreversibility, or governance-equivalent replacements for those roles. If it fails to supply any one of those governance roles, non-degenerate construction collapses by Theorem~\ref{thm:mutual-closure-kernel}. A coarser package may compress multiple roles into one clause, and a finer package may split one role across several clauses. But if the total governance work and the resulting class of governed constructions remain unchanged, then by Definition~\ref{def:governance-equivalence} the package is governance-equivalent to $K$ at the level that matters here. If the governance work changes, the package is governance-non-equivalent and is not sufficient for non-degenerate construction in the present sense. Therefore the uniqueness claim concerns required governance work, not the uniqueness of one four-part bookkeeping decomposition.
\end{proof}

\begin{remark}
This theorem concerns governance-role equivalence, not representational identity and not uniqueness of a particular factorization. It does not exclude alternative formal encodings, type-theoretic or categorical implementations, game-semantic presentations, or coarser/finer decompositions that preserve the same governance work. It excludes only packages that alter the class of governed constructions and therefore perform different governance work.
\end{remark}

\subsection{Fixed-Domain Closure Results}

In the present section, ``same domain'' always means preservation of the admissibility-relevant invariant bundle and hence preservation of the class of governed constructions up to governance equivalence. The point is not to classify arbitrary statuses or metrics, but to describe what fixed-domain admissibility can and cannot look like once the kernel is in place.

The bivalence and AMetricity claims below concern only gate status for standing-bearing commitment, licensed reuse, and admissible continuation. They do not deny object-language many-valued semantics, paraconsistent valuations, probabilistic confidence, proof-search rankings, approximation orders, proof-state metadata, type refinements, or graded measures internal to an already admitted construction regime. Such structures are admissible only when they do not themselves decide standing-bearing entry. If they do decide entry, they are part of the admissibility gate and fall under the present results; if they do not, they are downstream bookkeeping or internal structure.

These results do more than describe one convenient interface. Taken together, they close the remaining underdetermination question for the admissibility invariant itself. Once a fixed domain already requires an admissibility boundary, every same-domain realization must either reproduce the interface proved in this section up to governance-equivalence or else fall into one of the failure modes already excluded: hidden selector work, repair-sensitive middle status, illicit structure import, domain change, or inert bookkeeping. The point of the section is therefore structural rigidity: necessity plus uniqueness of same-domain realization.

\begin{lemma}[No intermediate admissibility state]\label{lem:no-intermediate-status}
Let $S$ be any admissibility-relevant status distinct from admitted and not-admitted on a fixed domain. Then exactly one of the following holds:
\begin{enumerate}[label=(\roman*)]
\item $S$ affects admissible continuation, commitment, or licensed reuse, in which case $S$ performs hidden selector work or reopens a repair-sensitive middle status; or
\item $S$ affects none of those, in which case $S$ is inert bookkeeping rather than an admissibility-relevant status.
\end{enumerate}
Hence no third admissibility-relevant status exists.
\end{lemma}

\begin{proof}
Take any status $S$ distinct from admitted and not-admitted. If possession of $S$ changes what may later count as licensed continuation, commitment, or reuse, then $S$ adds governance-relevant discrimination beyond the gate itself. In that case it either selects among otherwise eligible acts or preserves a repair-sensitive middle state between rejection and admission. If instead possession of $S$ changes nothing about later licensed use, then $S$ contributes no governance work and is merely descriptive bookkeeping. Since every admissibility-relevant status must either affect licensed use or fail to affect it, these two cases are exhaustive. Therefore no third admissibility-relevant status survives.
\end{proof}

\begin{theorem}[Necessity and bivalence of admissibility]\label{thm:bivalence}
On every fixed domain supporting licensed construction, a prior admissibility predicate is forced, and any admissibility-relevant status interface on that same fixed domain is extensionally equivalent to a binary one. Equivalently, every admissibility-relevant status falls into exactly one of two classes: admitted for commitment, reuse, and standing-bearing continuation, or not admitted.
\end{theorem}

\begin{proof}
Without a prior admissibility predicate, entry into licensed composition would have to occur before evaluation, by post hoc repair, by scope transport, or by injected parameterization. Each route breaks reference, composition, or irreversibility. So a prior gate is forced. Now suppose the resulting admissibility interface were not extensionally binary. Then some additional same-domain status would have to survive between, or alongside, admitted and not admitted. If that status changes later admissible use, it is hidden gate work or a repair-sensitive middle state and therefore contradicts the fail-closed consequences of irreversibility. Any state whose admissibility status can change without constituting a numerically distinct act violates irreversibility and is therefore excluded. If it changes no later admissible use, it is inert bookkeeping and not an admissibility-relevant status at all. There is no third possibility. So every admissibility-relevant interface collapses extensionally to admitted or not admitted.
\end{proof}

\begin{theorem}[AMetric boundary and non-parameterization]\label{thm:ametric}
At the admissibility boundary, no admissibility-relevant selector, ranking, ordering, grading, metric, or discrete indexing survives. In particular, there is no admissibility-relevant family of distinct positive statuses beyond the binary admitted / not-admitted split.
\end{theorem}

\begin{proof}
On the pre-boundary side, no representative, coordinate, or labeling choice may carry admissibility-relevant authority for the same fixed domain. This label-neutrality is not a free modeling preference. If admissibility were allowed to depend on labels, coordinates, or presentation indices as such, then two identity-preserving presentations of the same candidate would receive different admissibility significance solely by representation. That would make representation itself a same-domain selector and would violate the requirement that admissibility be invariant under identity-preserving transformation. So any admissibility-relevant differentiator must be invariant under governance-preserving relabeling. Under that invariance condition, only equality-pattern structure survives: there is no nontrivial unary grading, no admissibility-relevant order, no selector of a privileged positive route, and no meaningful metric beyond discrete equality. Suppose, then, that a parameterized or ranked boundary nevertheless existed. If the parameter or ranking makes a governance-relevant difference, it selects among otherwise admissible candidates or routes and therefore introduces a same-domain selector not preserved by the relabeling symmetry just described. If it makes no governance-relevant difference, it is bookkeeping rather than boundary structure. Likewise, if the extra structure is used to encode an intermediate positive status rather than a selector, it reopens a repair-sensitive middle state between admitted and not admitted, contrary to the fail-closed consequence of irreversibility. Thus every purported admissibility-relevant parameterization or metric either contradicts the boundary invariance conditions or collapses to bookkeeping. Accordingly, any purported admissibility-relevant parameterization at the boundary is forbidden.
\end{proof}

\begin{theorem}[Standing equals reuse-stable admissibility]\label{thm:standing-reuse}
Let $\mathrm{Stand}_D$ be any predicate claimed to mark those evaluated acts on a fixed domain $D$ that are eligible for standing-bearing continuation, commitment, or licensed reuse. Then for every act $a \in D$,
\[
\mathrm{Stand}_D(a)=1 \iff a \text{ is reuse-stably admissible on } D.
\]
Equivalently, standing is not an extra primitive or interface-level classifier beyond admissibility on a fixed scope. It is the persistence condition of admissibility under licensed same-scope continuation.
\end{theorem}

\begin{proof}
If $\mathrm{Stand}_D(a)=1$, then standing is already doing admissibility-relevant gate work on the fixed scope, so by Theorem~\ref{thm:bivalence} there is no additional governance-relevant classifier beyond the unique gate. Thus $a$ is admitted, and if it were not reuse-stable some licensed same-scope continuation would silently retarget or reclassify the act. Conversely, if $a$ is reuse-stably admissible but denied standing, then admitted acts are split into two substantive positive subclasses. If that split affects later continuation, it is forbidden same-side gate work; if it affects nothing, it is inert bookkeeping. So the split is impossible as a substantive classification.
\end{proof}

\begin{theorem}[Extensional uniqueness of the admissible interior]\label{thm:unique-interior}
For a fixed system $S$ on a fixed domain, the admissible interior is unique up to extensional, or lawful, equivalence. More precisely, if $I_1$ and $I_2$ are predicates on acts such that each contains exactly the standing acts of $S$, then for every act $a$,
\[
I_1(a) \iff I_2(a).
\]
Equivalently: in the standing-instantiated regime, there are no plural admissible interiors in extension; any adequate encoding of the admissible interior collapses pointwise to the standing predicate.
\end{theorem}

\begin{proof}
Assume that $I_1$ and $I_2$ both capture exactly the standing acts on the same fixed domain but that they disagree on some act $a$. Then the same act receives two different standing verdicts on the same fixed scope. If one now adds a selector deciding which interior governs $a$, the selector itself becomes a new same-domain admissibility-relevant parameter, contrary to Theorems~\ref{thm:bivalence} and \ref{thm:ametric}. If one adds no such selector, then the fixed-domain standing classification is indeterminate, contrary to the role of standing in Theorem~\ref{thm:standing-reuse}. So two extensional interiors capturing exactly the standing acts cannot disagree on any act. Their apparent plurality therefore collapses extensionally to one admissible interior.
\end{proof}

\begin{theorem}[Conservation of standing]\label{thm:conservation}
For each fixed domain, no internal operation of the licensed construction regime creates, amplifies, transfers, or recovers standing. Standing can only be preserved or destroyed.
\end{theorem}

\begin{proof}
Recovery would be same-act repair; creation would be standing from a failed act without fresh evaluation; amplification would require standing to admit degrees despite Theorem~\ref{thm:bivalence}; and transfer would require an independent carrier of standing distinct from the gate. Each branch is excluded by the preceding kernel packet. So only preservation or destruction remain.
\end{proof}

\begin{theorem}[Structural rigidity of the admissibility invariant]\label{thm:rigidity}
Let $D$ be any fixed domain supporting governed construction. Let $J_D$ be any candidate same-domain realization of the admissibility invariant sufficient to fix eligibility, standing-bearing continuation, and reference-preserving commitment on $D$. Then exactly one of the following holds:
\begin{enumerate}[label=(\roman*)]
\item $J_D$ is governance-equivalent to the interface already fixed by the present results: a binary admissibility gate with an AMetric boundary, standing equal to reuse-stable admissibility, a unique admissible interior, and conservation of standing;
\item $J_D$ introduces fresh same-domain selector work, repair-sensitive middle status, or other illicit structure and therefore fails as a faithful realization on $D$;
\item $J_D$ changes the admissibility-relevant invariant bundle and therefore changes domain rather than realizing the same invariant on the same domain; or
\item $J_D$ contributes no independent governance work and collapses to bookkeeping.
\end{enumerate}
In particular, the admissibility invariant required for standing-bearing construction admits no governance-non-equivalent same-domain realization.
\end{theorem}

\begin{proof}
By Theorem~\ref{thm:governance-work-uniqueness}, any package sufficient for non-degenerate construction must already cover the same governance roles as $K$ up to governance-equivalence; there is no further governance-non-equivalent package that preserves the same class of governed constructions. Theorems~\ref{thm:bivalence} and \ref{thm:ametric} fix the interface shape of any such package on a fixed domain: admissibility-relevant status collapses extensionally to admitted or not admitted, and no same-domain selector, ranking, metric, or parameter survives at the boundary. Theorem~\ref{thm:standing-reuse} then eliminates any additional positive-side classifier beyond reuse-stable admissibility; Theorem~\ref{thm:unique-interior} eliminates plural same-domain admissible interiors; and Theorem~\ref{thm:conservation} eliminates creation, amplification, transfer, and same-act recovery of standing.

Suppose, then, that $J_D$ is a candidate realization of the required invariant on the same fixed domain. If it preserves all of the governance work just listed, then by Theorem~\ref{thm:governance-work-uniqueness} it is governance-equivalent to the interface already proved, giving (i). If it differs by introducing fresh selector work, middle statuses, repair channels, or imported authority, then it contradicts Theorems~\ref{thm:bivalence}, \ref{thm:ametric}, and \ref{thm:conservation}, yielding (ii). If it preserves some other boundary only by altering what counts as the same target, the same standing-bearing continuation, or the same admissibility-relevant invariant bundle, then it is no longer realizing the same invariant on the same fixed domain, yielding (iii). If it changes no governance verdict at all, then it is an extensional redescription or bookkeeping layer, yielding (iv). Any further candidate would have to change governance work while avoiding the listed branches, which is exactly what the rigidity result excludes.
\end{proof}

\subsection{Consequences for Admissible Structure}

The previous sections establish that the foundational layer is fixed, minimal, and non-derivable. The present section derives several structural consequences once that package is in place. These are not additional assumptions; they are strict consequences of the fixed package itself. All consequences in this section follow strictly from the impossibility of same-domain selector introduction, repair-sensitive status, or invariant violation established in Sections~4--6. The forcing pattern is uniform. A purported same-domain alternative to the results below must either (i) introduce an additional admissibility-relevant parameter or selector, (ii) reopen a repair channel or middle status between admission and rejection, or (iii) preserve all same-domain governance work and therefore collapse extensionally to the original interface. The first two branches contradict the kernel results already established; the third yields only bookkeeping redescription. The arguments below therefore proceed by contradiction against those three possibilities rather than by stipulating an interface in advance.

\begin{definition}[Invariant bundle and same target]
An invariant bundle is a family of admissibility-relevant invariants sufficient to determine when two acts have the same governed target. For acts $a$ and $b$, write $a \sim b$ to mean that they agree on every invariant in the bundle.
\end{definition}

\begin{definition}[Scope-preserving continuation]
A continuation $C$ is scope-preserving iff it remains within one admissibility-relevant frame and does not change target by illicit redescription.
\end{definition}

\begin{definition}[Admissible operator]
An operator or relation is admissible iff its application respects the governance of standing and admissibility rather than introducing distinctions by redescription alone.
\end{definition}

\begin{theorem}[Scope-preserving invariance]\label{thm:scope-preserving-invariance}
If $C$ is scope-preserving, then for every admissible act $a$, one has $a \sim C(a)$.
\end{theorem}

\begin{proof}
Assume $C$ is scope-preserving but changes an admissibility-relevant invariant of $a$. Then $a$ and $C(a)$ do not fix the same target. The continuation would therefore have introduced a governance-bearing difference while still claiming to remain within the same scope. That is precisely a silent redescription or hidden reparameterization, which admissibility excludes. Hence every scope-preserving continuation preserves the invariant bundle pointwise.
\end{proof}

\begin{theorem}[Transport closure]\label{thm:transport-closure}
Let $R$ be an admissible relation. If $b_1 \sim b_2$, then for every $q$,
\[
R(q,b_1) \Longleftrightarrow R(q,b_2).
\]
\end{theorem}

\begin{proof}
If $R$ could distinguish $b_1$ from $b_2$ despite $b_1 \sim b_2$, then the distinction would arise from presentation rather than governed target. But admissible construction cannot allow standing-bearing content to depend on mere redescription. Once target is fixed, admissible transport must be closed under that fixing. Therefore admissible relations cannot distinguish same-target presentations.
\end{proof}

\begin{theorem}[Domain-defined operators]
An admissible operator is defined exactly on the admissible domain, and in particular must be defined on every standing point applied to itself.
\end{theorem}

\begin{proof}
If an admissible operator were undefined on a standing point applied to itself, then a standing act would fail to lie in its own governed domain. The operator would not govern the admissible act as such. Conversely, if the operator had governance-bearing force outside the admissible domain, it would confer structure where standing is absent and thereby violate the admissibility boundary. Hence admissible operators are defined exactly on the admissible domain.
\end{proof}

\begin{theorem}[Quotient identity]
Identity of admissible presentation is exactly same-target identity. Equivalently, the admissible redescription orbit of an act coincides with its invariant bundle class.
\end{theorem}

\begin{proof}
Suppose two presentations lie in the same admissible orbit but differ in same-target identity. Then admissible redescription would change governed target, violating Theorem~\ref{thm:scope-preserving-invariance}. Conversely, suppose two presentations agree on every admissibility-relevant invariant but fail to count as identical for admissible purposes. Then admissibility would preserve a surplus distinction not grounded in target, violating Theorem~\ref{thm:transport-closure}. Therefore orbit identity and invariant identity coincide.
\end{proof}

\begin{corollary}[No silent redescription]
Any purported continuation that changes admissibility-relevant status while claiming same scope is illicit.
\end{corollary}

\begin{proof}
A same-scope continuation must preserve the invariant bundle by Theorem~\ref{thm:scope-preserving-invariance}. A status-changing continuation would therefore either leave scope or introduce illicit redescription. In either case it is not admissible as same-scope transport.
\end{proof}

\begin{remark}
The formal consequences listed here---transport closure, quotient identity, domain-defined operators, and scope-preserving invariance---therefore do not stand as fresh assumptions. They are structural fallout from the fixedness of the kernel. They belong downstream of foundational closure, not alongside it.
\end{remark}

\subsection{Closure Against Additional Foundational Conditions}

In this section, an \emph{additional foundational condition} means a putative further bearer of irreducible constructional work. It does not mean a merely declared symbol, shorthand, or presentational starting point. The question is whether anything beyond $K$ can function as a further base-level condition of non-degenerate construction.

Any such purported further condition can enter only in one of three ways: it can refine admissibility classification on the same fixed domain, purport to generate admissibility-relevant governance from a lower basis, or serve as an independent same-domain invariant alongside the kernel. These possibilities cover the relevant ways an additional condition could relate to fixed-domain admissibility: it may change admissibility classification, supply or generate it, or leave it untouched. If it leaves admissibility untouched, then it performs no independent governance work and is derivative rather than foundational.

\begin{theorem}[Exhaustion of foundational conditions]\label{thm:exhaustion-foundational-conditions}
Let $Q$ be any purported additional foundational condition. Then exactly one of the following holds:
\begin{enumerate}[label=(\roman*)]
\item $Q$ performs no independent governance work beyond the kernel $K$, in which case $Q$ is derivative rather than foundational;
\item $Q$ purports to generate admissibility-relevant governance from a lower same-domain basis; or
\item $Q$ purports to persist as an independent same-domain classifier or invariant alongside the kernel.
\end{enumerate}
Cases (ii) and (iii) are impossible: generation from below is excluded by the non-derivability and no-faithful-lower-generator results, while independent same-domain persistence is excluded by the no-faithful-same-domain-extension results together with the absence of any faithful same-domain counterexample. Therefore no additional foundational condition survives on the same fixed domain as an independent bearer of governance work.
\end{theorem}

\begin{proof}
Take any proposed foundational candidate $Q$. If $Q$ makes no difference to the class of governed constructions once $K$ is fixed, then by definition it is not foundational; it is descriptive, not load-bearing. Suppose instead that $Q$ performs independent governance work on the same fixed domain. Then, by the classification just stated, that extra work must enter in one of two nontrivial ways: either $Q$ is supposed to generate admissibility-relevant governance from a lower same-domain basis, or it is supposed to survive as an independent same-domain classifier or invariant alongside the kernel.

If $Q$ is claimed to be generated from a lower same-domain basis, then Corollary~\ref{cor:kernel-non-derivability} together with Theorem~\ref{thm:kernel-no-lower-generator} excludes the claim: no member of the kernel, and no genuinely foundational governance role, is generated from below while remaining on the same fixed domain.

If $Q$ is instead claimed to persist as an independent same-domain classifier or invariant, then it would constitute an additional same-domain bearer of governance work. But this is excluded by Theorem~\ref{thm:no-faithful-same-domain-extension}: there is no faithful same-domain extension that converts the original acts into a second independent governance architecture rather than collapsing back to the original admissibility boundary, importing illicit structure, violating the base goods, or proving trivial.

A purported fourth option would have to leave admissibility untouched while still doing independent foundational work. But a condition that leaves admissibility untouched leaves the class of governed constructions untouched as well, and so is derivative rather than foundational. Therefore the listed cases are complete. Hence any such condition either does no independent work or falls into one of the two excluded branches. Therefore no additional foundational condition survives as an independent bearer of governance work on the same fixed domain.
\end{proof}

\begin{corollary}[Foundational closure]
The foundational layer of non-degenerate construction is complete in the following sense: no admissible extension below or alongside $K$ survives on the same fixed domain as an independent bearer of governance work.
\end{corollary}

\begin{proof}
Immediate from Theorem~\ref{thm:exhaustion-foundational-conditions}.
\end{proof}

\subsection{Mechanization Boundary and Proof-Theoretic Scope}

\begin{theorem}[Internal mechanization boundary]
No proof system can establish the package $K$ as a genuinely new internal theorem from a lower layer while preserving the target domain fixed in the fixed-domain construction subsection.
\end{theorem}

\begin{proof}
Any proof system that already counts its outputs as admissibility-assessable constructions belongs to the target class of the fixed-domain construction subsection. By Theorem~\ref{thm:construction-forces-kernel}, any such system already operates with admissibility, standing, reference, and irreversibility in place. An ``internal construction'' of $K$ therefore begins under the very conditions it purports to generate. By Theorems~\ref{thm:kernel-no-lower-generator}, \ref{thm:regime-invariant}, and \ref{thm:no-faithful-same-domain-extension}, this is not generation from below but restatement, encoding, or circular recapture.
\end{proof}

\begin{remark}
The point is not anti-formal. It is to mark where formal work properly begins. Once the fixed package is in place, the remaining consequence layer is ordinary derivable mathematics and is an entirely appropriate target for mechanization.
\end{remark}

\begin{corollary}[No lower internal substrate]
A faithful mechanization of the present framework begins at the consequence layer rather than beneath admissibility and standing.
\end{corollary}

\begin{proof}
If a lower internal layer still counted as part of the same construction regime, then Theorem~\ref{thm:construction-forces-kernel} would already place admissibility and standing there. If it did not count as part of that regime, it would not amount to a lower constructional substrate for the present target. So the correct starting point for mechanization is the downstream consequence layer.
\end{proof}

\begin{remark}
A companion formalization is included only as corroboration of this downstream layer. It is not a premise in the argument of the self-contained kernel proof, and the non-derivability results do not depend on repository-level or file-level detail. For present purposes it is enough that representative families---split-collapse, standing/reuse closure, no same-act repair, no generators, no independent carriers, uniqueness of the admissible interior, conservation, and the supporting transport/invariance lemmas---have been machine-checked. That mechanized layer matters here not as a substitute for the present proofs, but as an external check that the equivalence-level uniqueness claims do not rest on merely verbal factorization.
\end{remark}



\subsection{Endpoint-facing kernel wrappers}\label{subsec:endpoint-facing-kernel-wrappers}

\begin{definition}[Theorem-bearing endpoint use]\label{def:endpoint-use}
For a fixed endpoint context \(R\) and a branch \(B\),
\[
  \EndpointUse_R(B)
\]
means that \(B\) is offered as a theorem-bearing resolution of the official endpoint image of \(R\). Endpoint use is not raw inscription and not mere report language; it is a governed mathematical act with target, carrier, theorem-status, redescription stability, and downstream reuse.
\end{definition}

\begin{theorem}[Endpoint use is governed construction]\label{thm:endpoint-use-governed}
For every fixed endpoint context \(R\) and branch \(B\),
\[
  \EndpointUse_R(B)\Rightarrow \Constr(B).
\]
Consequently endpoint use is subject to the kernel of non-degenerate construction.
\end{theorem}

\begin{proof}
A theorem-bearing endpoint must fix the target it resolves, preserve the carrier on which the problem is stated, distinguish theorem-status from raw trace or failed proof, permit downstream mathematical reuse, and remain stable under admissible redescription. These are precisely the features of governed construction for a determinate object on a fixed domain. Therefore endpoint use is governed construction.
\end{proof}

\begin{theorem}[Kernel necessity for theorem-bearing endpoints]\label{thm:kernel-necessity}
If \(\EndpointUse\), then the full kernel
\[
  K=\{\Adm,\St,\RefK,\Irr\}
\]
is already in force for that endpoint use.
\end{theorem}

\begin{proof}
By Theorem~\ref{thm:endpoint-use-governed}, endpoint use is governed construction. By Theorem~\ref{thm:construction-forces-kernel}, every meaningful non-degenerate governed construction already instantiates admissibility, standing, reference, and irreversibility. Thus the kernel is already in force for theorem-bearing endpoint use.
\end{proof}

\begin{theorem}[Ordinary theoremhood is kernel-internal]\label{thm:ordinary-theoremhood-kernel-internal}
For any proposition \(P\), if \(\ProofR(P)\) is a genuine proof rather than raw trace, then
\[
  \ProofR(P)\Rightarrow \Der(P)\Rightarrow K_R.
\]
In particular, an ordinary theorem of \(\neg\PosRaw\), if it exists as a theorem, is not rejected merely because it is negative.
\end{theorem}

\begin{proof}
A genuine proof is a proposition-targeted governed derivation. By Definition~\ref{def:derivation}, derivation is the proposition-targeted special case of governed construction. Corollary~\ref{cor:global-presupposition} states that every governed derivation already presupposes the full kernel. Therefore ordinary theoremhood is kernel-internal. The sign or polarity of the proposition does not change this conclusion.
\end{proof}

\begin{definition}[Ordinary official endpoint resolution]\label{def:ordinary-official-endpoint-resolution}
For a fixed endpoint context \(R\) and branch \(B\),
\[
  \OrdOfficialRes_R(B)
\]
means that ordinary mathematical practice uses a theorem-bearing act as the official resolution of the mathematical endpoint represented by \(R\) in branch \(B\). This ordinary-resolution act fixes the official target, fixes the carrier on which the problem is stated, fixes \(B\) as the branch being resolved, treats the result as theorem-bearing rather than support-level or heuristic, licenses downstream mathematical reuse as the settled endpoint result, is final at act-time rather than repairable as the same act, and is stable under lawful redescription of the same endpoint. This definition describes the ordinary mathematical role of official resolution; it is not an additional AASC status assignment.
\end{definition}

\begin{theorem}[Ordinary official endpoint resolution is endpoint use]\label{thm:ordinary-official-resolution-is-endpoint-use}
For every fixed endpoint context \(R\) and branch \(B\),
\[
  \OrdOfficialRes_R(B)\Rightarrow \EndpointUse_R(B).
\]
\end{theorem}

\begin{proof}
An ordinary official endpoint resolution is not mere theoremhood considered in isolation. It is theoremhood used to settle the official endpoint of a fixed problem. By Definition~\ref{def:ordinary-official-endpoint-resolution}, such an act fixes the target, carrier, branch, theorem-status, downstream reuse, act-time finality, and redescription-stable endpoint role. These are exactly the features of theorem-bearing endpoint use in Definition~\ref{def:endpoint-use}. If one of these features is absent, the theorem may remain ordinary mathematics, proof support, background content, or local obstruction, but it is not the official endpoint resolution of that fixed problem. Therefore \(\OrdOfficialRes_R(B)\Rightarrow\EndpointUse_R(B)\).

\end{proof}

\begin{theorem}[Ordinary theoremhood does not exempt official resolution from endpoint-role discipline]\label{thm:ordinary-theoremhood-not-shield}
Ordinary theoremhood certifies proof legitimacy; it does not decide endpoint-role equivalence. A theorem-bearing act that is not used as the official endpoint resolution may remain proposition-level theoremhood, support, or background mathematics. A theorem-bearing act used as the official endpoint resolution enters endpoint-role discipline:
\[
  \ProofR(B)\wedge\OrdOfficialRes_R(B)\Rightarrow \EndpointUse_R(B).
\]
\end{theorem}

\begin{proof}
The premise \(\ProofR(B)\) states ordinary theorem-bearing legitimacy. It does not by itself say whether the theorem occupies a positive endpoint image, supports a route, changes carrier, remains bookkeeping, or resolves the official endpoint. The premise \(\OrdOfficialRes_R(B)\) supplies the additional use: the theorem is being used as the official resolution of the fixed endpoint. By Theorem~\ref{thm:ordinary-official-resolution-is-endpoint-use}, that use is \(\EndpointUse_R(B)\). Thus ordinary theoremhood does not shield official endpoint resolution from the kernel/A+ endpoint-role discipline; it supplies the theorem-bearing act to which that discipline applies.
\end{proof}

\section{A+ as Kernel-Forced Consequence Layer}\label{sec:aplus-endpoint-consequence-layer}

\begin{definition}[A+ endpoint consequence package]\label{def:aplus-certificate}
For an endpoint context \(R\) and branch \(B\), \(\Aplus\) denotes the fixed-domain consequence package forced by kernel-instantiated theorem-bearing endpoint use. It consists of boundary fixation, endpoint bivalence, AMetric boundary, standing as reuse-stable admissibility, conservation of standing, unique admissible interior, no raw trace sufficiency, no selector import, no carrier/domain transfer, no third endpoint status, no auxiliary endpoint-status refinement, and report-preserving output discipline.
\end{definition}

\begin{theorem}[Kernel forces A+ consequences]\label{thm:a-plus}
For every fixed endpoint context \(R\) and branch \(B\),
\[
  \EndpointUse_R(B)\wedge \FixedDom(R)\Rightarrow \Aplus_R(B).
\]
\end{theorem}

\begin{proof}
By Theorem~\ref{thm:kernel-necessity}, endpoint use instantiates the kernel. Fixed-domain closure then applies the consequences proved in the kernel section. Boundary fixation is admissibility as the joint governance boundary. Endpoint bivalence follows from Theorem~\ref{thm:bivalence}. The AMetric boundary follows from Theorem~\ref{thm:ametric}. Standing as reuse-stable admissibility follows from Theorem~\ref{thm:standing-reuse}. Unique admissible interior follows from Theorem~\ref{thm:unique-interior}. Conservation follows from Theorem~\ref{thm:conservation}. No raw trace sufficiency follows from Lemma~\ref{lem:raw-not-construction}. No selector import and no carrier/domain transfer follow from Proposition~\ref{prop:cross-domain-closure} and Theorem~\ref{thm:rigidity}. No third endpoint status follows from Lemma~\ref{lem:no-intermediate-status}. No auxiliary endpoint-status refinement follows from structural rigidity. Report-preserving output discipline is the claim-report interface of the same reference, standing, admissibility, and irreversibility roles. Thus A+ is a consequence package, not an independent premise.
\end{proof}

\begin{theorem}[Detailed verification of official endpoint resolution target adequacy]\label{thm:official-resolution-satisfies-target-adequacy}
Let \(R\) be the official \(P\) versus \(NP\) endpoint context. Suppose ordinary endpoint resolution proceeds through the standard unrestricted finite CNF-SAT route on a context-bound adequate encoded model:
\[
  \OrdOfficialRes_R(B)\wedge \OfficialCNFSATRoute(R,model)
  \wedge \ModelOfCnf(model,R).
\]
Then the act of official endpoint resolution satisfies the four kernel-neutral target-adequacy conditions of Definition~\ref{def:target-adequacy}. Consequently, the resolution act instantiates the kernel and the fixed-domain \(\Aplus\) consequence layer:
\[
\begin{aligned}
  &\OrdOfficialRes_R(B)\wedge \OfficialCNFSATRoute(R,model)\\
  &\qquad\wedge\ModelOfCnf(model,R)\wedge\FixedDom(R)\\
  &\Rightarrow K_R\wedge \Aplus_R(B).
\end{aligned}
\]
\end{theorem}

\begin{proof}
We verify the four neutral adequacy conditions for the resolution act itself.

\emph{Target determinacy.} The target is the determinate official endpoint question, read through the standard unrestricted finite CNF-SAT correspondence, of which raw branch occupies or excludes the fixed positive candidate-image endpoint. The object of evaluation is repeatable: the standing-bearing endpoint role of the outcome on the context-bound encoded candidate image. The route predicate \(\OfficialCNFSATRoute(R,model)\) fixes the projection, and \(\ModelOfCnf(model,R)\) fixes the adequate encoded candidate image.

\emph{Step evaluability.} The steps of the resolution act are not merely inscriptions. They include asserting a theorem as the resolving outcome, classifying which branch has standing, fixing reportability, and licensing downstream reuse as the settled endpoint result. These steps are evaluable as advancing or departing from the fixed endpoint because the accepted SAT projection gives determinate branch roles: positive resolution is candidate-image occupation, while negative resolution is candidate-image non-occupation under official use.

\emph{Act-time finality.} If a later intervention changes which branch has official standing, reportability, or licensed downstream reuse, the later intervention produces a distinct endpoint-resolution act. It does not retroactively repair the original resolution act as the same act. This is the endpoint-resolution instance of act-time finality and of the no same-act repair discipline proved in the kernel section.

\emph{Same-regime fidelity.} Lawful redescription of the resolution act that preserves the official CNF-SAT route, the context-bound encoded model, and the fixed candidate-image carrier also preserves the target, the evaluability of the classification steps, and the act-time status of the resolution. If a reformulation changes the route, carrier, model, branch role, standing assignment, or act-time status, it is not the same endpoint-resolution regime but a scope or carrier shift.

Thus all four conditions of Definition~\ref{def:target-adequacy} hold for official endpoint resolution through the fixed CNF-SAT route. By Lemma~\ref{lem:target-adequacy-forces-kernel}, reference, standing, irreversibility, and admissibility are forced for the resolution act. By Theorem~\ref{thm:construction-forces-kernel}, the full kernel is in force for non-degenerate governed construction; and by Theorem~\ref{thm:a-plus}, endpoint use on the fixed domain yields the \(\Aplus\) consequence layer. Finally, by Theorem~\ref{thm:ordinary-official-resolution-is-endpoint-use}, the ordinary official-resolution act is \(\EndpointUse_R(B)\), so the displayed implication follows.

Consequently, any theorem-bearing endpoint classification within that resolution act that assigns official standing, reportability, branch exclusion, or licensed downstream reuse to one branch rather than the other on the fixed CNF-SAT carrier is an admissibility-relevant act within the governed endpoint-resolution regime. This is the precise bridge used later by the endpoint-admissibility relevance theorems; it is not an additional predicate imposed outside the kernel.
\end{proof}


\begin{corollary}[AMetricity from endpoint use]\label{cor:context-ametric}
On the fixed SAT endpoint image,
\[
  \ClayCtx\wedge\EndpointUse_R(B)\wedge\FixedDom(R)\Rightarrow \AMetric(R).
\]
\end{corollary}

\begin{proof}
The context supplies the carrier and branch slots. Endpoint use plus fixed-domain closure supplies \(\Aplus\) by Theorem~\ref{thm:a-plus}; the AMetric boundary is one of the listed kernel-forced A+ consequences.
\end{proof}



\section{Local SAT Kernel Packet and Weakening Resistance}\label{sec:local-sat-kernel-packet}

The self-contained kernel proof above establishes the general necessity of the kernel for non-degenerate construction. The present section records the SAT-local consequences actually used by the endpoint proof and audits the possible weakening routes. The claim is local: the manuscript does not require every formalism to use the same vocabulary. It requires that, on the fixed unrestricted finite CNF-SAT carrier, a theorem-bearing endpoint regime that preserves the same carrier, candidate image, branch slots, endpoint role, and act-time finality already performs the governance work named by the local kernel packet.

\subsection{Local SAT kernel packet}

\begin{longtable}{>{\raggedright\arraybackslash}p{0.10\linewidth}>{\raggedright\arraybackslash}p{0.25\linewidth}>{\raggedright\arraybackslash}p{0.56\linewidth}}
\toprule
Node & General necessity & Local SAT form \\
\midrule
K1 & Fixed-domain kernel & A theorem-bearing SAT endpoint requires the fixed unrestricted finite CNF carrier, encoded candidate image, branch slots, admissibility boundary, and act-time irreversibility. \\
K2 & Reference fixation & A SAT report cannot be licensed by silently replacing arbitrary finite CNF-SAT with fixed-width k-SAT, a promise fragment, nonuniform circuits, finite benchmarks, oracle behavior, per-length advice, or a different encoding carrier. \\
K3 & Standing-readout separation & \(\CnfSATP\) is the positive endpoint readout. Lower-bound data, solver hardness, failure traces, or candidate obstruction data are not that readout by definition. \\
K4 & Bivalence and fail-closed status & On the fixed endpoint role, a theorem-bearing branch is admitted or not admitted. Unknown, hard, unproved, heuristic, benchmark-failed, partial, or open status is not a third endpoint truth branch. \\
K5 & Unique admissible interior & The same SAT endpoint slot cannot carry plural independent endpoint-governing interiors. A second endpoint-status classifier must collapse to the bridge interface, reset carrier, or fail. \\
K6 & Standing-admissibility identity & Separator standing cannot float outside admissibility. A negative endpoint branch with standing must be admissibly typed; otherwise it is proof support, bookkeeping, or carrier drift. \\
K7 & No generators & Writing that all encoded polynomial-time candidates fail does not self-authorize endpoint standing. \\
K8 & No carriers & Finite tests, solver hardness, oracle separations, nonuniform fragments, proof-complexity fragments, lower-bound templates, or circuit families cannot carry CNF-SAT endpoint status without a lawful bridge back to the fixed carrier. \\
K9 & No post-selection repair & Later candidate discovery, encoding shift, oracle shift, reclassification, or benchmark extension creates a new certified act; it does not repair an earlier unbridged endpoint report as the same act. \\
K10 & AMetric no-selector boundary & No preferred encoding, candidate order, search depth, benchmark set, ranking, heuristic metric, or finite table decides endpoint status before admissibility is fixed. \\
K11 & Report preservation & A report may not exceed its carrier, candidate image, and bridge support. Candidate failure evidence alone cannot replace endpoint proof; candidate-image exclusion becomes endpoint governance only under endpoint-resolving use. \\
K12 & ATS skin non-authority & Labels, syntax, model names, code orderings, proof sketches, and bookkeeping do not become theorem-bearing by being written. \\
K13 & Endpoint exhaustion & Once the carrier, branch slots, endpoint role, and admissibility boundary are fixed, endpoint-bearing branches are exhausted by positive endpoint occupation or governed separator use. Proof support, bookkeeping, unknown status, deferred status, repair, carrier shift, unsupported separator status, and hidden-fifth statuses do not occupy endpoint truth. \\
\bottomrule
\end{longtable}

\begin{theorem}[Local SAT kernel packet]\label{thm:local-sat-kernel-packet}
For every theorem-bearing endpoint use on the fixed unrestricted finite CNF-SAT carrier, the local SAT consequences K1--K13 are in force.
\end{theorem}

\begin{proof}
By Theorem~\ref{thm:kernel-necessity}, endpoint use instantiates reference, standing, admissibility, and irreversibility. K1 records their joint fixed-domain role. K2 and K3 are reference and role-separation consequences. K4--K6 are fail-closed, unique-interior, and standing-admissibility consequences. K7--K9 block self-authorization, unlicensed carrier transfer, and post-selection repair. K10 blocks selector authority at the endpoint boundary. K11 is the SAT report-preservation rule. K12 prevents skin-level notation from becoming theorem-bearing. K13 is endpoint exhaustion: after the SAT endpoint coordinates are fixed, a purported negative branch is positive occupation, proof support, bookkeeping, repair, carrier reset, selector/surrogate import, lawful external comparison, or forbidden second endpoint authority. No additional same-domain endpoint truth occupation remains.
\end{proof}

\subsection{Weakening-resistance of the local SAT packet}

\begin{definition}[Core SAT endpoint adequacy packet]\label{def:core-sat-endpoint-adequacy}
The core SAT endpoint adequacy packet is
\[
  \CoreSAT(R,model)=\{\RefK_R,\St_R,\Adm_R,\Irr_R\},
\]
together with target determinacy, step evaluability, act-time finality, same-regime fidelity, the fixed unrestricted finite CNF-SAT carrier, and the context-bound encoded candidate image of Subsection~\ref{subsec:encoded-candidate-model-adequacy}.
\end{definition}

\begin{definition}[Separator-governance-permissive weakening]\label{def:separator-governance-permissive-weakening}
A separator-governance-permissive weakening is a regime \(\WeakSep\) such that:
\begin{enumerate}
\item \(\WeakSep\) preserves the same unrestricted finite CNF-SAT carrier;
\item \(\WeakSep\) preserves the same encoded candidate image, positive endpoint slot, and raw negative branch slot;
\item \(\WeakSep\) preserves theorem-bearing target-slot fixation, minimal report evaluability, same-regime fidelity, and act-time finality;
\item \(\WeakSep\) preserves the core packet \(\CoreSAT(R,model)\);
\item \(\WeakSep\) permits endpoint-resolving separator governance that is not positive endpoint occupation, proof support only, bookkeeping, carrier shift, repair, or a lawful coequal target role.
\end{enumerate}
Target-slot fixation in this definition is weaker than K13: it fixes what is being evaluated but does not itself exclude unknown, deferred, bookkeeping, or hidden-fifth endpoint statuses. Minimal report evaluability is weaker than K11: it attaches reports to the fixed target and classifies their role, but it does not itself impose no-overreporting. K11 and K13 are proved below as local necessities; they are not smuggled into the definition of \(\WeakSep\).
\end{definition}

\begin{lemma}[Plural SAT data are permitted; plural endpoint interiors are not]\label{lem:plural-sat-data-not-plural-interiors}
K5 does not prohibit multiple encodings, candidate procedures, proof routes, reductions, lower-bound attempts, solver traces, oracle comparisons, circuit families, proof-complexity systems, or heuristic studies. It prohibits only plural endpoint-governing admissible interiors for the same fixed SAT endpoint slot when those interiors assign endpoint status while preserving the same carrier, candidate image, branch roles, and endpoint target.
\end{lemma}

\begin{proof}
Plural SAT data may be proof support, auxiliary structure, representation, or bridge material without deciding endpoint status. If two alleged endpoint interiors agree on all endpoint-bearing acts, the second is endpoint bookkeeping. If they disagree while preserving the same carrier and endpoint slot, then the same SAT endpoint has two governing authorities. That is not harmless mathematical plurality; it is plural endpoint governance.
\end{proof}

\begin{lemma}[Local necessity of K5]\label{lem:sat-k5-necessity}
Under \(\CoreSAT(R,model)\), a weakening that permits two distinct endpoint-governing admissible interiors for the same SAT endpoint slot either collapses to bookkeeping, splits the domain, or creates a second same-domain endpoint classifier.
\end{lemma}

\begin{proof}
Let \(I_1\) and \(I_2\) be endpoint-governing interiors for the same fixed endpoint slot. If they agree on all endpoint-bearing acts, they are extensionally identical for endpoint purposes. If they disagree, then the disagreement must be resolved by a higher selector, a domain split, or a second classifier. A higher selector is endpoint governance above the gate; a domain split exits same-carrier use; a second classifier is precisely the forbidden plural endpoint authority. Thus K5 is locally necessary for non-degenerate SAT endpoint determinacy.
\end{proof}

\begin{lemma}[Local necessity of K6]\label{lem:sat-k6-necessity}
Under \(\CoreSAT(R,model)\), separator standing outside admissibility is either support-only, carrier shift, or a second endpoint-status gate.
\end{lemma}

\begin{proof}
Suppose a separator branch has endpoint standing while not being admissibility-bound. If it does not affect endpoint status, it is support, bookkeeping, or lower-status report material. If it affects endpoint status on the same carrier, it performs endpoint-status work outside the admissibility boundary and therefore acts as a second gate. If it avoids that conclusion by changing the carrier, encoding image, branch role, or endpoint slot, it is carrier or domain shift.
\end{proof}

\begin{lemma}[Local necessity of K11]\label{lem:sat-k11-necessity}
Under \(\CoreSAT(R,model)\) plus minimal report evaluability, a theorem-bearing SAT endpoint report that exceeds its carrier, encoded candidate image, bridge position, and support is unsupported endpoint authority, silent redescription, or no-carrier transfer.
\end{lemma}

\begin{proof}
Minimal report evaluability says that a report is attached to the fixed target and can be classified as theorem-bearing, support-level, bookkeeping, or out-of-domain. K11 adds no-overreporting: endpoint force may not exceed the fixed carrier, candidate image, bridge position, and tensor support. A finite obstruction, heuristic lower-bound trace, or candidate failure datum may support a proof; it does not by itself settle the endpoint. If it is reported as endpoint failure on the same carrier without lawful bridge support, it overstates its support. The overstatement is unsupported endpoint authority, silent redescription of a surrogate as the endpoint, or transfer from another carrier.
\end{proof}

\begin{lemma}[Local necessity of K13]\label{lem:sat-k13-necessity}
K13 does not deny epistemic uncertainty, open problem status, incomplete proof search, heuristic evidence, or review status. It states that unknown, deferred, unsupported, bookkeeping, support-level, and hidden-fifth statuses are report or epistemic statuses, not endpoint truth branches on the fixed SAT endpoint slot.
\end{lemma}

\begin{proof}
Target-slot fixation fixes what is being evaluated. It does not by itself exhaust the endpoint truth slot. K13 supplies the fail-closed endpoint-exhaustion rule: once carrier, branch slots, endpoint role, and admissibility boundary are fixed, an endpoint branch is admitted for endpoint standing or not. Unknown status may describe the state of a problem; support status may describe evidence; bookkeeping may organize notation. None of these occupies endpoint truth. Treating any of them as a third theorem-bearing endpoint branch degenerates the fixed endpoint regime.
\end{proof}

\begin{theorem}[No strict same-carrier weakening permits independent SAT separator governance]\label{thm:no-strict-sat-weakening-permits-separator-governance}
Under the core SAT endpoint adequacy packet \(\CoreSAT(R,model)\), no separator-governance-permissive weakening \(\WeakSep\) preserves non-degenerate same-carrier SAT endpoint use.
\end{theorem}

\begin{proof}
Assume such a weakening exists, and let \(G_{\mathsf{sep}}\) be the endpoint-resolving separator act it permits. Since \(G_{\mathsf{sep}}\) is endpoint-resolving, it changes the endpoint status of the fixed SAT slot. Since it is not proof support or bookkeeping, it is theorem-bearing. Since it is not carrier shift or repair, it acts on the same carrier and same endpoint act. Since it is not positive endpoint occupation, it does not occupy the constructive SAT endpoint. Since it is not a lawful coequal target role, it cannot be a separately declared alternative endpoint inside the same positive endpoint-image carrier.

If \(G_{\mathsf{sep}}\) shares the same admissible interior as the positive bridge, it collapses into bridge occupation and cannot supply separator exclusion. If it has a distinct admissible interior, Lemma~\ref{lem:sat-k5-necessity} yields plural endpoint governance, a second classifier, or domain split. If it has standing without admissibility, Lemma~\ref{lem:sat-k6-necessity} makes it support-only, carrier shift, or second gate. If its report force exceeds carrier and bridge support, Lemma~\ref{lem:sat-k11-necessity} gives unsupported endpoint authority, silent redescription, or no-carrier transfer. If it is neither admitted nor rejected while remaining endpoint-resolving, Lemma~\ref{lem:sat-k13-necessity} gives third-status degeneration. In every case the alleged weakening either preserves the original endpoint behavior extensionally, lowers the separator to support or bookkeeping, shifts carrier, or reproduces forbidden second endpoint authority. Hence no strict same-carrier weakening survives.
\end{proof}

\begin{corollary}[Strict SAT weakening either preserves the packet extensionally or exits endpoint use]\label{cor:sat-weakening-extension-or-exit}
Let \(\WeakSep\) be a weaker same-carrier SAT regime preserving core endpoint adequacy. If \(\WeakSep\) does not permit independent endpoint governance, it is extensionally equivalent to the local SAT packet for endpoint purposes. If it does permit independent endpoint governance, it fails reference, standing, admissibility, irreversibility, target-slot fixation, minimal report evaluability, no-overreporting, endpoint exhaustion, or same-carrier fidelity.
\end{corollary}

\begin{proof}
If the weakening does not alter endpoint-bearing statuses, its differences are support-level or presentation-level for the endpoint. If it alters endpoint status by permitting independent separator governance, Theorem~\ref{thm:no-strict-sat-weakening-permits-separator-governance} shows that it exits non-degenerate same-carrier endpoint use by one of the listed failures.
\end{proof}


\section{Endpoint-Status Discriminator Discipline}\label{sec:discriminator-discipline}

\begin{definition}[Endpoint-status discriminator]\label{def:endpoint-discriminator}
An endpoint-status discriminator for a branch \(B\) on context \(R\), written
\[
  \EndpointDisc_R(C,B),
\]
is a status factor, rule, classifier, role-marker, selector, or admissibility-relevant distinction by which \(B\) is asserted to occupy an endpoint status on the fixed endpoint image of \(R\).
\end{definition}

\begin{definition}[Independent same-domain discriminator]\label{def:independent-discriminator}
A discriminator \(C\) is independent same-domain governance work for \(B\), written \(\IndDisc_R(C,B)\), iff it changes endpoint-status, admissibility, standing-bearing continuation, or reportable theorem use on the same fixed domain while not being governance-equivalent to the already fixed admissibility/standing interface.
\end{definition}

\begin{definition}[Endpoint-status work]\label{def:endpoint-status-work}
\(\EndpointStatusWork_R(C,B)\) means that the theorem-bearing use of \(C\) for branch \(B\) changes, settles, or assigns the endpoint status of the fixed endpoint image of \(R\). This is a functional role, not a separate source of standing: if the act does not determine endpoint status, it is not endpoint-resolving.
\end{definition}


\begin{definition}[Endpoint-admissibility relevance]\label{def:endpoint-admissibility-relevance}
For a fixed endpoint context \(R\), a classifier \(C\) attached to a branch \(B\) is endpoint-admissibility-relevant, written
\[
  \EndpointAdmRel_R(C,B),
\]
when varying or removing \(C\) changes at least one admissibility-relevant feature of the endpoint-resolution act: whether \(B\) has endpoint standing, whether \(B\) may be reported as the endpoint result, whether \(B\) may be reused downstream as resolved theorem content, whether a continuation of the endpoint-resolution act remains admissible on the same carrier, or whether the competing branch is excluded from endpoint standing. This is the endpoint-resolution instance of the admissibility-relevant invariant bundle of Definition~\ref{def:invariant-bundle}; it is not a new primitive or a new source of standing.
\end{definition}

\begin{theorem}[Endpoint resolution classifications are governed-construction data]\label{thm:endpoint-resolution-classifications-governed}
Let \(B\) be used as a theorem-bearing official endpoint resolution on a fixed carrier. Any classification that determines the endpoint standing, reportability, or downstream reuse of \(B\) is data of the governed endpoint-resolution act and is therefore subject to the same kernel and fixed-domain consequence layer as the act itself.
\end{theorem}

\begin{proof}
Official endpoint resolution is theorem-bearing endpoint use. By Theorem~\ref{thm:endpoint-use-governed}, endpoint use is governed construction. The resolution act has a fixed target, a fixed carrier, theorem-status rather than support-status, act-time finality, and redescription-stable downstream reuse. A classification that determines which branch may stand, be reported, or be reused as the endpoint result changes the admissibility-relevant invariant bundle of that resolution act. It is therefore governed-construction data for the act and cannot be treated as an inert descriptive label outside the kernel consequence layer.
\end{proof}

\begin{theorem}[Endpoint-status work is endpoint-admissibility-relevant]\label{thm:endpoint-status-work-adm-relevant}
For theorem-bearing endpoint use on a fixed carrier, endpoint-status work is endpoint-admissibility-relevant:
\[
  \EndpointUse_R(B)\wedge\FixedDom(R)\wedge\EndpointStatusWork_R(C,B)
  \Rightarrow \EndpointAdmRel_R(C,B).
\]
\end{theorem}

\begin{proof}
Endpoint-status work changes, settles, or assigns the status of a branch at the fixed endpoint. Such a status assignment determines whether the branch may stand as the endpoint result, whether it may be reported, whether it may be reused downstream, and whether the competing branch is excluded from endpoint standing. Those are precisely admissibility-relevant features of the endpoint-resolution act in Definition~\ref{def:endpoint-admissibility-relevance}. Hence endpoint-status work under endpoint use is endpoint-admissibility-relevant.
\end{proof}

\begin{theorem}[Endpoint discriminator trichotomy with domain exit]\label{thm:endpoint-discriminator-trichotomy}
Let \(C\) be any endpoint-status discriminator attached to branch \(B\) on a fixed domain \(R\). Exactly one of the following governance cases applies:
\[
\begin{array}{ll}
\Bookkeep_R(C,B), & \GovEquiv_R(C,B),\\[2pt]
\IndDisc_R(C,B), & \DomShift_R(C,B).
\end{array}
\]
That is, \(C\) either changes no endpoint-governance verdict, reproduces the existing kernel-forced interface, performs independent same-domain governance work, or changes the admissibility-relevant invariant bundle and leaves the domain.
\end{theorem}

\begin{proof}
If \(C\) changes no admissibility, standing, continuation, endpoint-status, or reportability verdict, it is bookkeeping. If it changes verdicts only in a way already fixed by the kernel-forced admissibility/standing interface, it is governance-equivalent. If it changes endpoint-status verdicts on the same invariant bundle without being governance-equivalent, it performs independent same-domain governance work. If it changes the invariant bundle by which target, carrier, standing-bearing continuation, or admissibility classification are individuated, it changes domain. These cases exhaust whether a discriminator changes governance work and whether the change remains on the fixed domain.
\end{proof}

\begin{theorem}[Endpoint-status work criterion for independence]\label{thm:endpoint-status-work-criterion-independence}
Let \(C\) be endpoint-status work for branch \(B\) on a fixed endpoint carrier. If the act is neither bookkeeping, nor governance-equivalent to the admitted endpoint interface, nor a domain shift, then it is independent same-domain endpoint-status governance:
\[
\begin{aligned}
  &\EndpointStatusWork_R(C,B)\wedge\FixedDom(R),\\
  &\neg\Bookkeep_R(C,B)\wedge\neg\GovEquiv_R(C,B),\\
  &\neg\DomShift_R(C,B)
  \quad\Rightarrow\quad \IndDisc_R(C,B).
\end{aligned}
\]
\end{theorem}

\begin{proof}
By Definition~\ref{def:endpoint-status-work}, endpoint-status work is a discriminator for the endpoint status of the fixed image. The trichotomy of Theorem~\ref{thm:endpoint-discriminator-trichotomy} says that such a discriminator is bookkeeping, governance-equivalent, independent same-domain work, or domain-shifting. The hypotheses exclude bookkeeping, governance-equivalence, and domain shift. The remaining trichotomy branch is independent same-domain endpoint-status governance.
\end{proof}

\begin{theorem}[Independent same-domain discriminator is forbidden]\label{thm:independent-discriminator-forbidden}
For theorem-bearing endpoint use on a fixed domain,
\[
  \EndpointUse_R(B)\wedge\FixedDom(R)\Rightarrow
  \neg\exists C\,\IndDisc_R(C,B).
\]
\end{theorem}

\begin{proof}
Endpoint use plus fixed-domain closure yields the A+ consequence package by Theorem~\ref{thm:a-plus}. The AMetric boundary excludes selector, ranking, metric, or parameter work at the endpoint boundary. Bivalence excludes third endpoint status. Standing as reuse-stable admissibility and the unique admissible interior exclude a second same-domain endpoint gate. Conservation excludes creation, recovery, transfer, or amplification of standing by an extra carrier. Transport closure excludes same-domain redescription changes. Therefore no independent same-domain endpoint-status discriminator survives. By Theorem~\ref{thm:rigidity}, any attempted realization is governance-equivalent, illicit same-domain selector work, domain-changing, or bookkeeping; the independent same-domain case is the illicit case.
\end{proof}


\begin{corollary}[A+ excludes independent endpoint-admissibility-relevant discriminators]\label{cor:aplus-excludes-independent-endpoint-admrel}
For theorem-bearing endpoint use on a fixed domain, an independent same-domain discriminator that is endpoint-admissibility-relevant is excluded:
\[
  \EndpointUse_R(B)\wedge\FixedDom(R)\wedge
  \EndpointAdmRel_R(C,B)\wedge\IndDisc_R(C,B)
  \Rightarrow \bot.
\]
\end{corollary}

\begin{proof}
By Theorem~\ref{thm:independent-discriminator-forbidden}, endpoint use on a fixed domain excludes every independent same-domain endpoint discriminator attached to the branch. The additional hypothesis \(\EndpointAdmRel_R(C,B)\) records that the discriminator lies in the admissibility-relevant class controlled by the kernel: it changes branch standing, reportability, downstream reuse, or admissible continuation for the endpoint-resolution act. Thus the contradiction is not triggered by a mere descriptive label; it is triggered by independent same-domain admissibility-relevant endpoint governance.
\end{proof}


\begin{corollary}[A+ consequences bind ordinary official endpoint resolution]
\label{cor:aplus-binds-ordinary-official-resolution}
Let \(B\) be a branch used as the ordinary official resolution of a fixed endpoint on a fixed domain. If an endpoint-status discriminator \(C\) for that use is independent same-domain work, then the same endpoint use also triggers the \(\Aplus\) no-independent-discriminator closure against \(C\):
\[
\begin{aligned}
  &\OrdOfficialRes_R(B)\wedge\FixedDom(R)\wedge\IndDisc_R(C,B)\\
  &\qquad\Rightarrow
  \Aplus_R(B)\wedge\neg\IndDisc_R(C,B).
\end{aligned}
\]
\end{corollary}

\begin{proof}
By Theorem~\ref{thm:ordinary-official-resolution-is-endpoint-use}, ordinary official endpoint resolution is \(\EndpointUse_R(B)\). Theorem~\ref{thm:a-plus} gives \(\Aplus_R(B)\) from \(\EndpointUse_R(B)\wedge\FixedDom(R)\). Theorem~\ref{thm:independent-discriminator-forbidden} says that no independent same-domain endpoint-status discriminator survives under theorem-bearing endpoint use on the fixed domain. Hence the identified independent discriminator is excluded. The conclusion is not an optional overlay on ordinary mathematics; it is the fixed-domain consequence of the endpoint-use role already supplied by ordinary official resolution.
\end{proof}

\begin{remark}[A+ is not an optional overlay on ordinary practice]
\label{rem:aplus-not-optional-overlay}
The manuscript does not require ordinary mathematical practice to adopt AASC vocabulary. It requires only that ordinary practice, when it performs official endpoint resolution on a non-degenerate fixed carrier, cannot consistently maintain both that the act is \(\EndpointUse_R\) and that it is exempt from the fixed-domain consequences of \(\EndpointUse_R\). By the kernel concession lemma, any assessment that grants determinate target identity, step evaluability, act-time finality, and same-regime fidelity has already granted the kernel roles under whatever terminology it prefers. The \(\Aplus\) layer is downstream of those roles; it is not a discretionary second standard.
\end{remark}

\begin{corollary}[EndpointUse plus trichotomy binds the independent role]
\label{cor:endpointuse-trichotomy-aplus-binding}
Suppose an ordinary official endpoint-resolution act on a fixed carrier is classified by the endpoint-discriminator trichotomy as independent same-domain endpoint-status work. Then the no-independent-discriminator consequence applies to that very act. In symbols, for any discriminator \(C\) attached to branch \(B\),
\[
  \OrdOfficialRes_R(B)\wedge\FixedDom(R)\wedge\IndDisc_R(C,B)
  \Rightarrow \neg\IndDisc_R(C,B).
\]
\end{corollary}

\begin{proof}
This is the second component of Corollary~\ref{cor:aplus-binds-ordinary-official-resolution}. The trichotomy identifies the role; ordinary official resolution supplies endpoint use; the kernel-forced \(\Aplus\) consequence layer excludes the independent role. Ordinary theoremhood as proof legitimacy does not suspend that consequence once the theorem is used as official endpoint resolution.
\end{proof}

\begin{corollary}[Role-specific classification is consequence, not source]\label{cor:role-specific-classification}
If a branch has a role-specific endpoint status, that status is admissible only as a consequence of endpoint use under the kernel-forced fixed-domain interface. It does not create standing for the branch.
\end{corollary}

\begin{proof}
By Theorem~\ref{thm:kernel-necessity}, endpoint use already presupposes the kernel. By Theorem~\ref{thm:a-plus}, the endpoint-status interface is downstream of that kernel. Hence role-specific classification may record the status a branch has under the fixed interface, but cannot be a source of kernel standing.
\end{proof}

\section{SAT Branch Asymmetry and Separator Occupation}\label{sec:sat-branch-asymmetry}

\begin{definition}[Ordinary negative theorem]\label{def:ordinary-negative-theorem}
\[
  \NegThm_R := \ProofR(\neg\PosRaw).
\]
This denotes ordinary theoremhood of the negative SAT proposition inside a governed theorem regime. It is not, by itself, a distinct separator endpoint-status discriminator.
\end{definition}

\begin{theorem}[Negative theoremhood is not automatically defective]\label{thm:negative-theorem-not-defective}
\[
  \NegThm_R\Rightarrow K_R.
\]
Thus the manuscript does not reject \(\ProofR(\neg\PosRaw)\) merely for being negative.
\end{theorem}

\begin{proof}
This is Theorem~\ref{thm:ordinary-theoremhood-kernel-internal} applied to \(P=\neg\PosRaw\). Genuine theoremhood is governed derivation and is already kernel-internal. The excluded object below is not theoremhood of a negative proposition; it is the same-domain separator discrimination induced by the SAT image-separator branch.
\end{proof}


\begin{theorem}[Ordinary negative theoremhood is not the rejected discriminator]\label{thm:ordinary-negative-not-discriminator}
Ordinary theoremhood of the negative proposition is not, merely as theoremhood, the independent same-domain separator discriminator excluded below:
\[
  \NegThm_R \;\text{by theorem status alone does not equal}\; \SepDisc(model).
\]
When the content proved is the raw negative branch \(\SepBare\), the SAT-specific lower-bound/no-positive bridge transports that content to candidate endpoint-image exclusion. Only when the negative content is offered as the official endpoint resolution is that support-level exclusion promoted to \(\SepImgOcc\). The resulting contradiction targets the discriminator induced by endpoint occupation, not negative theoremhood as such.
\end{theorem}

\begin{proof}
Ordinary theoremhood is a proposition-targeted governed derivation, hence kernel-internal by Theorem~\ref{thm:negative-theorem-not-defective}. The discriminator excluded below is not a theorem-status predicate. Theorem~\ref{thm:bare-negative-forces-image} supplies only candidate endpoint-image exclusion; endpoint-resolution use is required before that support-level content becomes separator endpoint-image occupation. Thus an ordinary negative theorem is not rejected for polarity or for being theoremhood; the rejected object is the same-domain separator endpoint-status discrimination induced only after official endpoint-resolution use promotes the negative content to \(\SepImgOcc\).
\end{proof}



\begin{theorem}[Ordinary SAT falsity is not lost; endpoint-bearing negation is role-classified]\label{thm:ordinary-sat-falsity-role-classified}
The manuscript does not deny that \(\neg\PosRaw\) is the ordinary logical negation of the positive SAT endpoint. It denies only that this negation can occupy theorem-bearing endpoint status on the same fixed SAT carrier without taking one of the audited roles. If it is proposition-targeted theoremhood only, it is kernel-internal and not rejected. If it is proof support, it does not resolve the official endpoint. If it changes the carrier, it is not the same endpoint. If it is offered as the official negative endpoint resolution, it becomes endpoint use of \(\SepBare\), and the SAT lower-bound/no-positive bridge transports it first into \(\CandImageExcl(model)\) and then, under that endpoint-resolution use, into separator endpoint-image occupation \(\SepImgOcc\).
\end{theorem}

\begin{proof}
Theorem~\ref{thm:negative-theorem-not-defective} grants ordinary negative theoremhood as kernel-internal theoremhood. Theorem~\ref{thm:ordinary-negative-not-discriminator} separates that theoremhood from the rejected separator discriminator. Theorem~\ref{thm:ordinary-theoremhood-not-shield} records the missing distinction: ordinary theoremhood is proof legitimacy, not endpoint-role exemption. The additional act of offering the negative theorem as the official endpoint resolution is classified by Definition~\ref{def:official-negative-endpoint-resolution}. Under that use, Theorem~\ref{thm:official-negative-theoremhood-endpoint-use} yields endpoint use of \(\SepBare\); the SAT-native bridge then routes \(\SepBare\) through lower-bound normal form to \(\CandImageExcl(model)\), and official endpoint-resolution use promotes that content to \(\SepImgOcc\). Thus ordinary negation remains ordinary negation, while endpoint-bearing use of that negation is role-classified.
\end{proof}

\begin{theorem}[The SAT separator branch is not a lawful coequal endpoint role]\label{thm:separator-not-lawful-coequal-role}
Inside the fixed unrestricted finite CNF-SAT endpoint-image carrier, \(\SepBare\) is not a separately declared lawful coequal target role alongside positive endpoint occupation. The positive branch is direct constructive occupation of \(\CnfSATP\). A same-carrier negative branch that resolves the official endpoint is proof support, bookkeeping, carrier/domain shift, or endpoint-status governance by candidate-image exclusion. It is not an additional coequal truth slot inside the same positive endpoint-image carrier.
\end{theorem}

\begin{proof}
The official problem admits the ordinary mathematical alternatives \(P=NP\) and \(P\ne NP\). The present proof spine, however, evaluates the unrestricted finite CNF-SAT endpoint carrier whose positive endpoint image is polynomial-time CNF-SAT decidability. On that carrier, Theorem~\ref{thm:positive-direct-occupation} identifies the positive branch as direct constructive endpoint occupation. The negative CNF branch is the raw complement of that positive image and, by the SAT-native lower-bound/no-positive bridge, becomes candidate endpoint-image exclusion and, under endpoint-resolution use, separator endpoint-image occupation. If this negative branch is not endpoint-resolving, it is support or bookkeeping. If it changes the carrier, it is not the same endpoint. If it resolves the same carrier negatively, it governs endpoint status from outside positive occupation and therefore enters independent-discriminator analysis. Hence it is not a lawful coequal endpoint role inside the same endpoint-image carrier.
\end{proof}

\begin{theorem}[Separator exclusion is not a ban on negative complexity theory]\label{thm:negative-complexity-not-globally-forbidden}
The proof does not forbid negative complexity-theoretic mathematics. Circuit lower bounds, proof-complexity lower bounds, relativization or oracle separations, natural-proofs and algebrization barriers, failed candidate algorithms, finite solver-hardness evidence, local obstructions, and proposition-targeted negative theoremhood are legitimate when they remain in their lawful roles. The excluded structure is narrower: same-carrier endpoint-resolving separator governance that does not occupy the positive CNF-SAT endpoint, does not remain proof support or bookkeeping, does not change carrier, and nevertheless governs endpoint status through an independent same-domain discriminator.
\end{theorem}

\begin{proof}
Negative mathematical data may be proof support, local obstruction, barrier analysis, domain diagnosis, or ordinary theoremhood. None of those roles is excluded by the no-independent-discriminator theorem. The excluded case is the residue after lawful roles have been removed: a same-carrier endpoint-resolving negative branch that does endpoint-status work outside positive occupation while preserving the same carrier. By Theorem~\ref{thm:ordinary-sat-falsity-role-classified} this is endpoint-bearing negative governance, and by Theorem~\ref{thm:separator-not-lawful-coequal-role} it is not a lawful coequal endpoint role inside the same positive endpoint-image carrier. It is therefore routed to the independent-discriminator closure.
\end{proof}

\begin{definition}[Official negative endpoint resolution]\label{def:official-negative-endpoint-resolution}
\[
  \OfficialEndpointResolution_R(\neg\PosRaw)
\]
means that an ordinary proof of the negative proposition is offered not merely as a proposition-targeted derivation, but as the official negative resolution of the P versus NP endpoint on the unrestricted finite CNF-SAT carrier. Equivalently for the SAT branch under audit, it records the ordinary official-resolution use \(\OrdOfficialRes_R(\SepBare)\) of the raw negative branch. It is not a source of theorem standing and does not make negative theoremhood defective.
\end{definition}

\begin{theorem}[Official negative theoremhood becomes endpoint use]\label{thm:official-negative-theoremhood-endpoint-use}
There is no standing intermediate position between ordinary negative theoremhood and negative endpoint use when the theorem is offered as a resolution of the official P versus NP endpoint. Formally, on the fixed SAT carrier,
\[
\begin{aligned}
  &\ClayCtx\wedge\FixedDom(R)\wedge\ModelOfCnf(model,R)\wedge\NegThm_R\\
  &\qquad\wedge\OfficialEndpointResolution_R(\neg\PosRaw)
    \Rightarrow \EndpointUse_R(\SepBare).
\end{aligned}
\]
\end{theorem}

\begin{proof}
An ordinary proof of \(\neg\PosRaw\), considered only as a proposition-targeted derivation, is governed theoremhood and is not rejected for being negative. The additional premise \(\OfficialEndpointResolution_R(\neg\PosRaw)\) changes the use of that theorem: the proof is now being offered as the official negative resolution of the unrestricted finite CNF-SAT endpoint. By Definition~\ref{def:official-negative-endpoint-resolution}, this is the SAT-specific ordinary official-resolution use \(\OrdOfficialRes_R(\SepBare)\). Theorem~\ref{thm:ordinary-official-resolution-is-endpoint-use} then gives \(\EndpointUse_R(\SepBare)\). If the proof refuses that endpoint-use status, then it is no longer being used as the official endpoint resolution of the same target; it is merely non-endpoint theoremhood, bookkeeping negation, carrier shift, or a degenerate act relative to the fixed endpoint.
\end{proof}

\begin{corollary}[No floating official negative endpoint]\label{cor:no-floating-official-negative-endpoint}
A proof of \(\neg\PosRaw\) cannot simultaneously be offered as the official negative resolution of the unrestricted finite CNF-SAT endpoint and refuse negative endpoint-use status. If it is not endpoint use, it does not resolve the official endpoint; if it is endpoint use, it enters the SAT negative-occupation and endpoint-status-bivalence results below.
\end{corollary}

\begin{proof}
Immediate from Theorem~\ref{thm:official-negative-theoremhood-endpoint-use} and the ordinary-resolution bridge, Theorem~\ref{thm:ordinary-official-resolution-is-endpoint-use}. The only harmless form of ordinary negative theoremhood is theoremhood considered without official endpoint-resolution use. Once the theorem is offered as the endpoint resolution, refusing endpoint-use discipline withholds the very status required for the act to be the same theorem-bearing resolution act.
\end{proof}

\begin{theorem}[Ordinary official negative SAT resolution is non-positive status work]\label{thm:ordinary-official-negative-sat-resolution-nonpositive-status-work}
On the context-bound fixed CNF-SAT endpoint image, ordinary official negative resolution of the P versus NP endpoint is non-positive endpoint-status work:
\[
\begin{aligned}
  &\ClayCtx\wedge\FixedDom(R)\wedge\ModelOfCnf(model,R)\wedge\SepBare\\
  &\qquad\wedge\OfficialEndpointResolution_R(\neg\PosRaw)
    \Rightarrow \NonPositiveStatusWork_R(model).
\end{aligned}
\]
\end{theorem}

\begin{proof}
By Definition~\ref{def:official-negative-endpoint-resolution}, official negative resolution is the SAT-specific ordinary official-resolution use of the raw negative branch. By Theorem~\ref{thm:ordinary-official-resolution-is-endpoint-use}, it is endpoint use rather than mere theoremhood or support. The carrier-adequacy bridge, Theorem~\ref{thm:candidate-image-carrier-adequacy-raw-split}, identifies \(\SepBare\) with \(\CandImageExcl(model)\), and Theorem~\ref{thm:candidate-image-exclusion-to-separator-occupation} promotes that exclusion to \(\SepImgOcc\) only under official endpoint-resolution use. Theorem~\ref{thm:sepocc-is-non-positive-status-work} then yields \(\NonPositiveStatusWork_R(model)\). Thus ordinary negative theoremhood remains ordinary theoremhood as proof legitimacy, but its use as the official endpoint resolution has the non-positive status role fixed by the candidate-image carrier.
\end{proof}

\begin{theorem}[Positive branch is direct constructive endpoint occupation]\label{thm:positive-direct-occupation}
The positive CNF-SAT branch directly occupies the constructive endpoint carrier:
\[
  \CnfPositiveEndpoint\Rightarrow\CnfDirectPos .
\]
Equivalently, when \(\PosRaw\) is obtained as \(\CnfSATP\), its endpoint role is the direct polynomial-time CNF-SAT endpoint. It is not a same-domain separator classifier and does not require an independent endpoint-status discriminator.
\end{theorem}

\begin{proof}
The positive branch asserts the existence of a deterministic polynomial-time decision procedure for arbitrary finite CNF-SAT. That assertion directly occupies the constructive endpoint carrier named by \(\CnfSATP\). No further separator relation over the endpoint image is introduced: the branch does not classify the positive route from outside; it is the positive route. The corresponding audited formal theorem is listed in Appendix~A.
\end{proof}

\begin{theorem}[Raw positive branch equals the formal CNF positive endpoint]\label{thm:raw-positive-formal-positive}
On any context-bound CNF-SAT endpoint model,
\[
\begin{aligned}
  &\ClayCtx\wedge\FixedDom(R)\wedge\ModelOfCnf(model,R)\\
  &\qquad\Rightarrow
  (\PosRaw\leftrightarrow\CnfPositiveEndpoint).
\end{aligned}
\]
\end{theorem}

\begin{proof}
The premise \(\ModelOfCnf(model,R)\) fixes the formal CNF-SAT endpoint model associated with the manuscript's raw carrier. By Definition~\ref{def:bare}, \(\PosRaw:=\CnfSATP\), the assertion that arbitrary finite CNF-SAT is decidable in deterministic polynomial time. In the formal SAT operator layer, \(\CnfPositiveEndpoint\) names that same positive endpoint on the encoded CNF model. Thus the raw positive branch and the formal positive endpoint are equivalent on the fixed model.
\end{proof}

\begin{corollary}[Bare negative branch equals the formal bare negative branch]\label{cor:bare-negative-formal}
On any context-bound CNF-SAT endpoint model,
\[
\begin{aligned}
  &\ClayCtx\wedge\FixedDom(R)\wedge\ModelOfCnf(model,R)\wedge\SepBare\\
  &\qquad\Rightarrow
  \CnfBareNeg(R,model).
\end{aligned}
\]
\end{corollary}

\begin{proof}
By Theorem~\ref{thm:bare-complement}, \(\SepBare\leftrightarrow\neg\PosRaw\). By Theorem~\ref{thm:raw-positive-formal-positive}, \(\PosRaw\leftrightarrow\CnfPositiveEndpoint\) on the context-bound CNF model. Hence \(\SepBare\) gives \(\neg\CnfPositiveEndpoint\), which is exactly \(\CnfBareNeg(R,model)\) in the formal SAT operator notation.
\end{proof}
\begin{theorem}[Canonical CNF model binding is instantiated, not arbitrary]\label{thm:canonical-cnf-model-binding}
The fixed SAT context supplies its own canonical encoded CNF model, denoted \(m_{\mathsf{SAT}}\), rather than making an arbitrary \(model\) context-bound:
\[
  \ClayCtx_{\RSATctx}\wedge\FixedDom(\RSATctx)\Rightarrow\ModelOfCnf(m_{\mathsf{SAT}},\RSATctx).
\]
All branch-local theorems below keep \(\ModelOfCnf(model,R)\) as an explicit premise until the final instantiated SAT context supplies \(m_{\mathsf{SAT}}\).
\end{theorem}

\begin{proof}
The CNF model is fixed by the official SAT endpoint carrier: arbitrary finite CNF instances over finite Boolean alphabets, evaluated by formula size, with admissible encodings preserving the same endpoint-image carrier. By Definition~\ref{def:cnf-model-of-context}, context binding includes the model-adequacy packet of Subsection~\ref{subsec:encoded-candidate-model-adequacy}: coverage, soundness, failure adequacy, and image exactness. This determines the canonical context-bound model for \(\RSATctx\). It does not license the stronger and false reading that every arbitrary \(model\) is context-bound.
\end{proof}


\subsection{SAT-native lower-bound normal form and endpoint-status governance}\label{subsec:sat-native-lower-bound-governance}

The remaining SAT-local burden is to make the raw negative branch visible in ordinary lower-bound language before the endpoint-status exclusion theorem is applied. The point of this subsection is not to import Cook--Levin/Karp machinery and not to add a new source of AASC standing. It records the SAT-native bridge formalized across \LeanName{AnticheckerReduction.lean} and \LeanName{SATOperatorProofQueue.lean}: the context-bound candidate model is adequate for ordinary deterministic polynomial-time CNF-SAT candidates; the raw negative branch is standard lower-bound normal form over that exact encoded candidate image; that normal form is candidate endpoint-image exclusion; candidate endpoint-image exclusion supplies theorem-level endpoint-status discrimination; and an endpoint-resolving negative use of that exclusion is endpoint-status governance, not a harmless proof-support observation.

\begin{definition}[Standard SAT lower-bound normal form]\label{def:standard-lower-bound-normal-form}
For a context-bound encoded CNF-SAT candidate model, \(\StdLB(model)\) denotes the standard lower-bound normal form:
\[
  \forall c\,\bigl(\mathsf{PolyCand}_{model}(c)\Rightarrow\mathsf{FailsSAT}_{model}(c)\bigr).
\]
Here \(c\) ranges over encoded deterministic polynomial-time candidate SAT procedures, \(\mathsf{PolyCand}_{model}(c)\) says that the code is polynomial-time in the fixed encoded CNF model, and \(\mathsf{FailsSAT}_{model}(c)\) says that the decoded procedure fails to decide arbitrary finite CNF-SAT on that carrier. The corresponding Lean predicate is \LeanName{CnfSATStandardLowerBoundNormalForm}.
\end{definition}

\begin{definition}[Candidate endpoint-image exclusion]\label{def:candidate-endpoint-image-exclusion}
\(\CandImageExcl(model)\) denotes universal exclusion of every encoded polynomial-time candidate from occupying the positive SAT endpoint image. In Lean this is \LeanName{CnfSATCandidateEndpointImageExclusion}. It is a theorem-level statement about the fixed encoded candidate image; it is not the assertion that a practical witness selector or uniform failing-instance generator has been constructed.
\end{definition}

\begin{definition}[Theorem-level endpoint-status discriminator]\label{def:theorem-level-endpoint-status-discriminator}
\(\ThmLevelDisc(model)\) denotes the theorem-level discriminator induced by candidate endpoint-image exclusion: the global status assignment under which each encoded polynomial-time candidate is classified as a non-occupant of the positive CNF-SAT endpoint image. The corresponding Lean predicate is \LeanName{CnfSATTheoremLevelEndpointStatusDiscriminator}. This is not an algorithmic separator and not witness selection; it is endpoint-status classification at theorem level.
\end{definition}

\begin{theorem}[Bare negative branch has standard lower-bound normal form]\label{thm:bare-negative-standard-lower-bound-form}
On the fixed context-bound CNF-SAT endpoint carrier,
\[
  \ClayCtx\wedge\FixedDom(R)\wedge\ModelOfCnf(model,R)\wedge\SepBare
  \Rightarrow \StdLB(model).
\]
The audited Lean anchor is \LeanName{cnfSATBareNegativeBranch_standardLowerBoundNormalForm}.
\end{theorem}

\begin{proof}
By Corollary~\ref{cor:bare-negative-formal}, \(\SepBare\) gives the formal bare negative branch:
\[
  \CnfBareNeg(R,model).
\]
This is the negation of the positive CNF-SAT endpoint on the context-bound encoded model. In ordinary machine-language form, it is precisely the universal lower-bound normal form: every encoded deterministic polynomial-time candidate fails to occupy the positive SAT endpoint. The cited Lean anchor records the same route from the formal bare negative branch to standard lower-bound normal form.
\end{proof}

\begin{theorem}[Lower-bound normal form is candidate endpoint-image exclusion]\label{thm:standard-lower-bound-candidate-image-exclusion}
On the fixed encoded CNF-SAT candidate image,
\[
  \StdLB(model)\Rightarrow \CandImageExcl(model).
\]
The Lean counterpart is \LeanName{cnfSATImageSeparatorBranch_standardLowerBoundNormalForm} together with \LeanName{CnfSATCandidateEndpointImageExclusion}.
\end{theorem}

\begin{proof}
The standard lower-bound normal form says that every encoded polynomial-time candidate fails SAT on the fixed unrestricted finite CNF carrier. By Theorem~\ref{thm:encoded-positive-endpoint-exactness}, the positive endpoint image is exactly the image of encoded polynomial-time candidates with no SAT failure. Universal failure of those candidates is therefore universal exclusion from the positive endpoint image. This is candidate endpoint-image exclusion.
\end{proof}

\begin{theorem}[Candidate endpoint-image exclusion induces theorem-level candidate-image classification]\label{thm:candidate-image-exclusion-theorem-level-discriminator}
For the context-bound CNF-SAT endpoint model,
\[
  \CandImageExcl(model)\Rightarrow \ThmLevelDisc(model).
\]
The Lean anchor is \LeanName{cnfSATTheoremLevelEndpointStatusDiscriminator_of_candidateImageExclusion}; the direct bare-negative form is \LeanName{cnfSATBareNegativeBranch_theoremLevelDiscriminator}.
\end{theorem}

\begin{proof}
Candidate endpoint-image exclusion globally assigns each encoded polynomial-time candidate the endpoint-relevant status ``non-occupant of the positive CNF-SAT endpoint image.'' This is theorem-level candidate-image classification. It does not compute a failing formula uniformly, does not select a canonical witness, and does not claim an algorithmic separator. It becomes endpoint-status governance only when the classification is used as the endpoint-resolving negative branch on the fixed carrier; otherwise it remains lower-bound or proof-support content. The cited Lean anchor records the theorem-level classification surface.
\end{proof}

\begin{definition}[Legitimate invariant use versus endpoint-status governance]\label{def:invariant-use-kind}
The Lean type \LeanName{CnfSATInvariantUseKind} separates three uses of invariant or obstruction material:
\[
  \mathsf{auxiliaryInvariant},\qquad \mathsf{localObstruction},\qquad \mathsf{endpointStatusGovernance}.
\]
Auxiliary invariants and local obstructions are legitimate proof-support material. Endpoint-status governance on the fixed SAT carrier is the forbidden same-domain endpoint classifier. This is recorded by \LeanName{CnfSATInvariantUseLegitimate} and \LeanName{cnfSATEndpointStatusGovernance_not_legitimate}.
\end{definition}

\begin{center}
\small
\begin{tabularx}{0.94\linewidth}{>{\raggedright\arraybackslash}p{0.30\linewidth}>{\centering\arraybackslash}p{0.14\linewidth}>{\raggedright\arraybackslash}X}
\toprule
Use kind & Legitimate? & Reason \\
\midrule
Auxiliary invariant & Yes & Supports a proof without governing the endpoint-status surface. \\
Local obstruction & Yes & Records local incompatibility or proof support inside an already fixed route. \\
Endpoint-status governance & No & Acts as independent same-domain authority over endpoint occupation on the fixed SAT carrier. \\
\bottomrule
\end{tabularx}
\end{center}

\begin{definition}[Candidate-image use classification]\label{def:candidate-image-exclusion-use-classification}
The Lean candidate-image-use type distinguishes four uses of candidate-image exclusion:
\[
\ProofSupportOnly,\quad
\EndResolvingNegThm,\quad
\EndResolvingNonGov,\quad
\CarrierChangingLB .
\]
The corresponding Lean classifier is \LeanName{CnfSATCandidateImageExclusionUseClassification}. In manuscript language: proof-support-only use does not resolve the endpoint; endpoint-resolving negative theorem use is official negative endpoint use; endpoint-resolving non-governance is the hidden fifth case; and carrier-changing lower-bound use exits the fixed carrier.
\end{definition}

\begin{theorem}[Endpoint-resolving non-governance is the hidden fifth case]\label{thm:endpoint-resolving-nongovernance-hidden-fifth-impossible}
There is no fixed-carrier use of candidate-image exclusion that both resolves the official endpoint negatively and avoids endpoint-status governance:
\[
  \neg\EndResolvingNonGov .
\]
The Lean anchor is \LeanName{cnfSATEndpointResolvingNonGovernance_hiddenFifthCase_impossible}.
\end{theorem}

\begin{proof}
If candidate-image exclusion is only proof support, it does not resolve the official endpoint. If it resolves the official endpoint negatively on the fixed carrier, it governs the endpoint-status surface by classifying candidate positive endpoint occupants as non-occupants. If it avoids that governance by changing the carrier, it is a carrier/domain shift. The alleged fourth possibility--endpoint-resolving but not endpoint-governing on the same carrier--is exactly the hidden fifth case, and the Lean classifier assigns it no inhabitant.
\end{proof}

\begin{definition}[Candidate-image endpoint governance]\label{def:endpoint-status-governance-candidate-image-exclusion}
\(\CandImageGov_R(model)\) means that candidate-image exclusion is used as endpoint-status governance on the fixed CNF-SAT carrier:
\[
\begin{aligned}
  \CandImageGov_R(model) := {}& \FixedDom(R)\wedge \neg\CnfDirectPos\\
  &{}\wedge\ThmLevelDisc(model).
\end{aligned}
\]
This matches the Lean predicate \LeanName{CnfSATEndpointStatusGovernanceByCandidateImageExclusion}: the carrier is fixed, the use does not occupy the positive endpoint, and theorem-level candidate-status discrimination is present.
\end{definition}

\begin{theorem}[Official negative endpoint use is candidate-image endpoint-status governance]\label{thm:official-negative-use-candidate-image-governance}
For the context-bound finite CNF-SAT carrier,
\[
  \CnfOfficialNegUse(R,model)\Rightarrow \CandImageGov_R(model).
\]
The Lean anchor is \LeanName{cnfSATOfficialNegativeEndpointUse_endpointStatusGovernance}.
\end{theorem}

\begin{proof}
Official negative endpoint use contains the fixed endpoint domain, the context-bound CNF model, and the bare negative branch. The bare negative branch yields candidate-image exclusion and then the theorem-level endpoint-status discriminator by Theorems~\ref{thm:bare-negative-standard-lower-bound-form}--\ref{thm:candidate-image-exclusion-theorem-level-discriminator}. Since the use is negative, it does not occupy the direct positive endpoint. Therefore it is endpoint-status governance by candidate-image exclusion, exactly as stated by the Lean theorem.
\end{proof}

\begin{corollary}[Endpoint-resolving negative theorem use is endpoint-status governance]\label{cor:endpoint-resolving-negative-theorem-governance}
If candidate-image exclusion is used as the endpoint-resolving negative theorem on the fixed CNF-SAT carrier, then it is endpoint-status governance:
\[
  \EndResolvingNegThm\Rightarrow \CandImageGov_R(model).
\]
The Lean anchor is \LeanName{cnfSATEndpointResolvingNegativeTheorem_is_endpointStatusGovernance}.
\end{corollary}

\begin{proof}
By Definition~\ref{def:candidate-image-exclusion-use-classification}, the endpoint-resolving negative-theorem case is official negative endpoint use. Theorem~\ref{thm:official-negative-use-candidate-image-governance} then gives endpoint-status governance by candidate-image exclusion. Proof-support-only use does not resolve the endpoint, and carrier-changing use exits the fixed carrier; the alleged endpoint-resolving non-governance use is ruled out by Theorem~\ref{thm:endpoint-resolving-nongovernance-hidden-fifth-impossible}.
\end{proof}

\begin{theorem}[Raw negative branch forces candidate endpoint-image exclusion]\label{thm:bare-negative-forces-image}
On the fixed context-bound CNF-SAT endpoint carrier,
\[
  \ClayCtx\wedge\FixedDom(R)\wedge\ModelOfCnf(model,R)\wedge\SepBare
  \Rightarrow \CandImageExcl(model).
\]
The audited Lean support is the route through \LeanName{cnfSATBareNegativeBranch_standardLowerBoundNormalForm} and the candidate-image exclusion bridge.
\end{theorem}

\begin{proof}
The implication is not a reportability principle and does not assert that raw truth creates a Clay report act. By Theorem~\ref{thm:bare-negative-standard-lower-bound-form}, \(\SepBare\) gives standard SAT lower-bound normal form over the fixed encoded candidate image. By Theorem~\ref{thm:standard-lower-bound-candidate-image-exclusion}, that normal form is candidate endpoint-image exclusion. Therefore \(\CandImageExcl(model)\) follows. No endpoint-image separator occupation is inferred at this stage.
\end{proof}

\begin{theorem}[Candidate-image exclusion becomes separator occupation only under official endpoint resolution]\label{thm:candidate-image-exclusion-to-separator-occupation}
On the fixed context-bound CNF-SAT endpoint carrier,
\[
\begin{aligned}
  &\ClayCtx\wedge\FixedDom(R)\wedge\ModelOfCnf(model,R)\wedge\CandImageExcl(model)\\
  &\qquad\wedge\OfficialEndpointResolution_R(\neg\PosRaw)
  \Rightarrow \SepImgOcc .
\end{aligned}
\]
\end{theorem}

\begin{proof}
Candidate-image exclusion by itself is lower-bound content: it may be proof support, a local obstruction, ordinary theorem content, or carrier-diagnostic information. The additional hypothesis \(\OfficialEndpointResolution_R(\neg\PosRaw)\) changes the act under assessment. It says that the negative content is being used as the official endpoint-resolving branch for the same unrestricted finite CNF-SAT endpoint. Under that use, the support-level image exclusion is promoted to separator endpoint-image occupation. This is precisely \(\SepImgOcc\). No support-level lower-bound statement is promoted merely by being true or written.
\end{proof}

\begin{lemma}[Separator occupation contains candidate-image exclusion]\label{lem:separator-occupation-contains-candidate-image-exclusion}
On the fixed context-bound CNF-SAT endpoint carrier, separator endpoint-image occupation includes the support-level candidate-image exclusion content:
\[
  \ClayCtx\wedge\FixedDom(R)\wedge\ModelOfCnf(model,R)\wedge\SepImgOcc
  \Rightarrow \CandImageExcl(model).
\]
\end{lemma}

\begin{proof}
By Definition~\ref{def:image}, \(\SepImgOcc\) is candidate-image exclusion after that content is used as the official endpoint-resolving negative branch on the fixed carrier. The occupation status therefore contains, rather than replaces, the support-level content \(\CandImageExcl(model)\). This lemma is an unpacking lemma only: it does not promote lower-bound content into endpoint occupation and does not assert endpoint use from \(\CandImageExcl(model)\) alone.
\end{proof}

\subsection{Ordinary Native Lower-Bound Falsification Collapse}\label{subsec:ordinary-native-lower-bound-collapse}

This subsection isolates the strongest ordinary-practice objection to the separator route. The objection says that a SAT lower-bound theorem is simply native mathematical falsification of \(P=NP\), not separator endpoint-status governance. The response is role-sensitive. Lower-bound syntax, proof support, ordinary theoremhood, and local obstruction data are preserved. What is classified is the stronger act: using the lower-bound package as the official same-carrier endpoint defeat of the unrestricted finite CNF-SAT endpoint.

\begin{definition}[Native lower-bound falsification package]\label{def:native-lower-bound-falsification-package}
For a fixed SAT endpoint context \(R\) and context-bound encoded candidate model \(model\), \(\LBNatSAT(R,model)\) denotes the proposed native lower-bound falsification package carrying all of the following data: the same unrestricted finite CNF-SAT carrier; the same raw positive branch \(\PosRaw:=\CnfSATP\); the same raw negative branch \(\SepBare\); the same context-bound encoded polynomial-time candidate image; the same model-adequacy packet of coverage, soundness, failure adequacy, and image exactness; the same standard lower-bound normal form
\[
\forall c\,(\mathsf{PolyCand}_{model}(c)\Rightarrow\mathsf{FailsSAT}_{model}(c));
\]
the same candidate endpoint-image exclusion \(\CandImageExcl(model)\); claimed endpoint-defeating force against \(\PosRaw\); and the claim that this use is merely ordinary native lower-bound falsification rather than separator endpoint-status governance.
\end{definition}

\begin{definition}[SAT endpoint-impact bundle]\label{def:sat-endpoint-impact-bundle}
\(\EndpointImpactSAT(R,model)\) is the endpoint-impact bundle for a same-carrier lower-bound use. It consists of target fixation, unrestricted finite CNF-SAT carrier fixation, context-bound encoded candidate-image fixation, branch fixation, endpoint standing, endpoint reportability, downstream reuse, exclusion of positive candidate-image occupation, act-time finality, and lawful-redescription stability. A use is endpoint-impacting when changing or removing it changes at least one component of this bundle.
\end{definition}

\begin{lemma}[Endpoint impact versus bookkeeping for SAT lower-bound use]\label{lem:sat-lb-endpoint-impact-vs-bookkeeping}
If a use of lower-bound or candidate-exclusion content changes no component of \(\EndpointImpactSAT(R,model)\), then it is bookkeeping for endpoint purposes.
\end{lemma}

\begin{proof}
The endpoint-impact bundle lists the features by which a fixed endpoint act differs from inert lower-bound notation: target, carrier, candidate image, branch role, standing, reportability, reuse, exclusion of the positive candidate-image occupant, act-time status, and lawful-redescription behavior. If a lower-bound use changes none of these features, it performs no endpoint work. It may remain a useful theorem label, proof organization device, or explanatory marker, but with respect to endpoint resolution it is bookkeeping.
\end{proof}

\begin{lemma}[Native lower-bound falsification below endpoint use]\label{lem:native-lb-below-endpoint-use}
If \(\LBNatSAT(R,model)\) is not deployed to defeat the official \(P\) versus \(NP\) endpoint through the standard unrestricted finite CNF-SAT route, then it remains below endpoint use: raw theoremhood, lower-bound support, local obstruction, proof-complexity support, candidate-failure evidence, bookkeeping, or carrier-diagnostic material.
\end{lemma}

\begin{proof}
Without deployment as endpoint defeat, the package does not exclude the positive endpoint as the settled official branch, does not license downstream reuse as a negative endpoint result, and does not occupy the separator endpoint-image role. It may still be meaningful mathematics: a formula, an ordinary theorem, a lower-bound lemma, an obstruction, a proof-support datum, or a diagnostic of a different carrier. Those roles are below endpoint use.
\end{proof}

\begin{lemma}[Deployment fixes official negative endpoint force]\label{lem:deployed-native-lb-official-negative-force}
If \(\LBNatSAT(R,model)\) is deployed to defeat the official endpoint through the standard unrestricted finite CNF-SAT route while preserving the fixed carrier and candidate image, then it is official negative endpoint-resolution use of \(\SepBare\).
\end{lemma}

\begin{proof}
Deployment to defeat the endpoint means that the package is not merely recorded as lower-bound content. It is used against \(\PosRaw\) as the negative resolution of the fixed official endpoint. Because the package preserves the unrestricted finite CNF-SAT carrier, the context-bound encoded candidate image, the standard lower-bound normal form, and candidate-image exclusion, the use is same-carrier negative endpoint force. This is official negative endpoint-resolution use of \(\SepBare\).
\end{proof}

\begin{lemma}[Deployed native lower-bound falsification is endpoint use]\label{lem:deployed-native-lb-endpoint-use}
A deployed native lower-bound falsification package is endpoint use:
\[
  \LBNatSAT(R,model)\text{ deployed as endpoint defeat}
  \Rightarrow \EndpointUse_R(\SepBare).
\]
\end{lemma}

\begin{proof}
By Lemma~\ref{lem:deployed-native-lb-official-negative-force}, deployed native lower-bound falsification is official negative endpoint-resolution use on the fixed carrier. By Theorem~\ref{thm:ordinary-official-resolution-is-endpoint-use}, ordinary official endpoint resolution is endpoint use. Therefore the deployed package is \(\EndpointUse_R(\SepBare)\).
\end{proof}

\begin{lemma}[Endpoint-use lower-bound falsification enters the trichotomy]\label{lem:endpoint-use-lb-enters-trichotomy}
A deployed native lower-bound falsification package is bookkeeping, governance-equivalent to positive candidate-image occupation, independent same-domain endpoint-status work, or a carrier/domain/interface shift.
\end{lemma}

\begin{proof}
Endpoint use fixes the branch role and places the act on the endpoint-status surface. A status factor attached to that use either changes no endpoint-impact component, in which case Lemma~\ref{lem:sat-lb-endpoint-impact-vs-bookkeeping} gives bookkeeping; reproduces the admitted positive candidate-image interface, in which case it is governance-equivalent to positive occupation; changes endpoint status on the same carrier without being equivalent, in which case it is independent same-domain endpoint-status work; or changes the carrier, candidate image, route, or interface, in which case it is domain or interface shift. These cases exhaust whether the deployed use changes endpoint work and whether it remains on the fixed carrier.
\end{proof}

\begin{lemma}[Universal candidate failure is not positive-equivalent]\label{lem:deployed-universal-failure-not-positive-equivalent}
A deployed native lower-bound falsification package cannot be governance-equivalent to positive candidate-image occupation while retaining endpoint-defeating negative force.
\end{lemma}

\begin{proof}
Positive candidate-image occupation is the existence of an encoded deterministic polynomial-time SAT candidate whose decoded procedure has no SAT failure. Universal candidate failure says that every encoded polynomial-time candidate fails SAT on the fixed carrier. If the deployed lower-bound package were governance-equivalent to positive occupation, it would not exclude the positive endpoint branch. If it excludes the positive branch, it performs a different endpoint-status role. Thus it is not governance-equivalent to positive occupation while retaining endpoint-defeating force.
\end{proof}

\begin{definition}[SAT bridge-neutralized lower-bound package]\label{def:sat-bridge-neutralized-lower-bound}
\(\BridgeNeutralizedSAT(model)\) means that the apparent lower-bound counterforce is neutralized by lawful positive candidate-image occupation:
\[
  \exists c\,(\mathsf{PolyCand}_{model}(c)\wedge
  \neg\mathsf{FailsSAT}_{model}(\mathsf{decode}_{model}(c))).
\]
Equivalently, the positive endpoint image \(\SATEndpointImageOcc(model)\) is occupied. In the live negative countercase, bridge-neutralization is incompatible with \(\CandImageExcl(model)\) on the same context-bound candidate image.
\end{definition}

\begin{lemma}[Bridge-neutralized lower-bound package is incompatible with live candidate-image exclusion]\label{lem:bridge-neutralized-lb-incompatible-live-exclusion}
On the fixed context-bound candidate image,
\[
  \CandImageExcl(model)\wedge\BridgeNeutralizedSAT(model)
  \Rightarrow\bot .
\]
\end{lemma}

\begin{proof}
Candidate-image exclusion says that every encoded polynomial-time candidate fails to occupy the positive endpoint image. Bridge-neutralization supplies an encoded polynomial-time candidate with no SAT failure, equivalently positive candidate-image occupation. The carrier, model, and failure predicate are fixed, so universal failure and a non-failing occupant cannot both hold as the same local endpoint branch.
\end{proof}

\begin{theorem}[Ordinary native lower-bound falsification collapse]\label{thm:ordinary-native-lb-falsification-collapse}
For the fixed unrestricted finite CNF-SAT endpoint carrier, a native lower-bound falsification package \(\LBNatSAT(R,model)\) has exactly the following same-mode possibilities: it is raw theoremhood, lower-bound support, local obstruction, or bookkeeping below endpoint use; it is bridge-neutralized by positive candidate-image occupation; it is bookkeeping for endpoint purposes; it is a carrier/domain/interface shift; or it is the independent SAT separator discriminator \(\SepDisc(model)\). There is no autonomous ordinary-native lower-bound endpoint role outside this list.
\end{theorem}

\begin{proof}
If the package is not deployed to defeat the endpoint, Lemma~\ref{lem:native-lb-below-endpoint-use} places it below endpoint use. If positive candidate-image occupation is supplied, Definition~\ref{def:sat-bridge-neutralized-lower-bound} gives bridge-neutralization. If the use changes no component of \(\EndpointImpactSAT(R,model)\), Lemma~\ref{lem:sat-lb-endpoint-impact-vs-bookkeeping} gives bookkeeping. If it changes the unrestricted finite CNF carrier, the encoded candidate image, the model-adequacy packet, the uniformity convention, the official route, or the endpoint interface, it is carrier/domain/interface shift.

Suppose none of those alternatives applies. Then the deployed use is endpoint-impacting, same-carrier, non-positive endpoint-status work. By Lemma~\ref{lem:deployed-native-lb-endpoint-use}, it is endpoint use. By Lemma~\ref{lem:endpoint-use-lb-enters-trichotomy}, the remaining possibilities are governance-equivalence to positive occupation or independent same-domain endpoint-status work. Lemma~\ref{lem:deployed-universal-failure-not-positive-equivalent} excludes governance-equivalence to positive occupation. Corollary~\ref{cor:candidate-image-exclusion-not-coequal-occupation} and Theorem~\ref{thm:negative-coequal-requires-interface-change} exclude coequal occupation of the same positive candidate-image carrier. The remaining role is independent same-domain separator-discriminator work, namely \(\SepDisc(model)\).
\end{proof}

\begin{corollary}[No sixth lower-bound role in the live SAT subproof]\label{cor:no-sixth-lower-bound-role-live-subproof}
Inside the exact-complement reductio subproof, the opened assumption \([\SepBare]^i\) cannot be deployed as ordinary native lower-bound endpoint falsification without yielding \(\SepDisc(model)\).
\end{corollary}

\begin{proof}
The opened assumption is the live exact-complement countercase against \(\PosRaw\), not inert syntax. The SAT projection sends it to standard lower-bound normal form and candidate-image exclusion. Endpoint counterforce removes support-only and bookkeeping readings; fixed-domain evaluation removes carrier and interface shift; and bridge-neutralization by positive candidate-image occupation contradicts the live candidate-image exclusion by Lemma~\ref{lem:bridge-neutralized-lb-incompatible-live-exclusion}. The collapse theorem therefore leaves \(\SepDisc(model)\).
\end{proof}

\paragraph{Technical closure.}
After this subsection, the objection that a SAT lower-bound theorem is simply ordinary native falsification of \(P=NP\) is not an unaddressed case. It is a challenge to Theorem~\ref{thm:ordinary-native-lb-falsification-collapse}: the objector must produce a same-carrier endpoint-defeating lower-bound role outside raw support, bridge-neutralized positive occupation, bookkeeping, carrier shift, and \(\SepDisc(model)\).


\subsection{SAT Negative Occupation Exhaustion and Endpoint-Status Bivalence}\label{subsec:sat-negative-occupation-bivalence}

The preceding bridge is not a representational convention. It is the SAT-local lower-bound/no-positive equivalence on the context-bound unrestricted finite CNF-SAT endpoint model. The audited Lean formalization separates two locks that make the negative-occupation step explicit. First, official negative endpoint use is exhausted into four cases: image-separator occupation, ordinary negative alternative, bookkeeping-only negative, or carrier/domain shift. Second, governed endpoint use on the fixed finite CNF-SAT carrier has only two endpoint statuses: positive occupation or separator occupation. These locks are included to block the claim that the official raw negative may remain a third theorem-bearing endpoint status that is neither positive nor separator.

The phrase ``official negative endpoint use'' is not a new escape hatch and not an AASC-added source of standing. By Theorem~\ref{thm:official-negative-theoremhood-endpoint-use}, ordinary theoremhood of \(\neg\PosRaw\) becomes negative endpoint use exactly when it is offered as the official negative resolution of the unrestricted finite CNF-SAT endpoint. If it is not so offered, it does not resolve the official target. If it is so offered, it enters the occupation-exhaustion and bivalence locks in this subsection.

\begin{definition}[Official negative occupation alternatives]\label{def:negative-occupation-alternatives}
For the context-bound CNF-SAT endpoint model, \(\CnfOfficialNegUse(R,model)\) denotes theorem-bearing official negative endpoint use of the finite CNF-SAT carrier. The Lean surface distinguishes four exhaustive negative-occupation alternatives:
\[
\begin{aligned}
  &\SepImgOcc,\qquad \OrdNeg_R^{\mathsf{alt}},\\
  &\CnfBookNegOnly(R,model),\qquad \CnfCarrierShift(R,model).
\end{aligned}
\]
The first is separator endpoint-image occupation. The second is a putative ordinary negative endpoint alternative that remains governed but does not occupy separator status. The third is bookkeeping-only negation, not theorem-bearing endpoint use. The fourth is a change of carrier or domain. The corresponding Lean objects are \LeanName{CnfSATOfficialNegativeEndpointUse}, \LeanName{CnfSATNegativeOccupationExhaustion}, \LeanName{CnfSATBookkeepingNegativeOnly}, and \LeanName{CnfSATCarrierShift}.
\end{definition}

\begin{theorem}[SAT negative occupation exhaustion]\label{thm:sat-negative-occupation-exhaustion}
Official negative endpoint use on the context-bound finite CNF-SAT carrier is exhausted by image-separator occupation, ordinary negative alternative, bookkeeping-only negative, or carrier/domain shift. In Lean this is the theorem \LeanName{cnfSATNegativeOccupation_exhaustion}.
\end{theorem}

\begin{proof}
The exhaustion is branch-complete by construction of the fixed SAT endpoint image. A theorem-bearing negative endpoint either occupies the image-separator branch, claims a distinct governed ordinary negative endpoint status, remains a bookkeeping-only negation without endpoint use, or changes the carrier/domain on which the official endpoint is being evaluated. These alternatives cover whether the negative is endpoint-bearing and, if endpoint-bearing, whether it remains on the fixed carrier and whether it occupies separator status. The cited Lean anchor records this exact exhaustion.
\end{proof}

\begin{theorem}[Non-optionality of SAT negative occupation]\label{thm:sat-negative-occupation-nonoptional}
On the fixed context-bound finite CNF-SAT endpoint carrier, official negative endpoint use together with fixed endpoint domain, context-bound CNF model, and the AMetric boundary entails image-separator occupation:
\[
\begin{aligned}
  &\CnfOfficialNegUse(R,model)\wedge \FixedDom(R)
    \wedge \ModelOfCnf(model,R)\\
  &\qquad\wedge \AMetric(R)\Rightarrow \SepImgOcc .
\end{aligned}
\]
The Lean anchor is \LeanName{cnfSATNegativeOccupation_nonoptional}.
\end{theorem}

\begin{proof}
Apply the occupation exhaustion of Theorem~\ref{thm:sat-negative-occupation-exhaustion}. Fixed endpoint domain excludes carrier/domain shift; official negative endpoint use excludes bookkeeping-only negation; and the AMetric boundary excludes the ordinary-negative-alternative branch. The corresponding Lean anchors are:
\begin{lstlisting}[basicstyle=\ttfamily\scriptsize]
cnfSATNoCarrierShift_of_fixedEndpointDomain
cnfSATOfficialNegativeEndpointUse_not_bookkeepingOnly
cnfSATNoOrdinaryNegativeAlternative_of_ametricBoundary
cnfSATNegativeOccupation_nonoptional
\end{lstlisting}
The remaining exhaustive branch is separator endpoint-image occupation.
\end{proof}

\begin{definition}[SAT endpoint-status bivalence]\label{def:governed-sat-endpoint-status-bivalence}
The Lean status type \LeanName{CnfSATEndpointStatus} has exactly two constructors, \LeanName{positive} and \LeanName{separator}. The predicates \LeanName{CnfSATGovernedEndpointUse} and \LeanName{CnfSATGovernedEndpointStatusUse} mark governed endpoint use and governed endpoint-status use on the context-bound finite CNF-SAT carrier.
\end{definition}

\begin{theorem}[Endpoint-status bivalence for governed CNF-SAT endpoint use]\label{thm:governed-endpoint-status-bivalence}
Governed endpoint use on the fixed finite CNF-SAT carrier is endpoint-status bivalent: it has positive status or separator status, and no third governed endpoint status. The Lean anchors are \LeanName{cnfSATEndpointStatus_exhaustive} and \LeanName{cnfSATGovernedEndpointUse_bivalent}.
\end{theorem}

\begin{proof}
The endpoint-status surface contains only the positive and separator constructors. A governed endpoint use must occupy one of these statuses. Any alleged third governed endpoint status would either be bookkeeping, a carrier shift, or an independent same-domain status refinement; those alternatives are excluded by the fixed-domain endpoint interface already established in Sections~\ref{sec:discriminator-discipline} and \ref{sec:sat-branch-asymmetry}. The Lean anchors record the status-exhaustive and governed-use bivalence forms.
\end{proof}

\begin{theorem}[Negative governed endpoint use has separator status]\label{thm:negative-governed-endpoint-use-separator}
For the governed finite CNF-SAT endpoint,
\[
\begin{aligned}
  &\CnfGovEndpointUse(R,model)\wedge \neg\CnfPositiveEndpoint\\
  &\qquad\Rightarrow \CnfImageSep(R,model).
\end{aligned}
\]
Equivalently, a governed endpoint use that is not positive occupies separator status and therefore the image-separator branch. The Lean anchors are \LeanName{cnfSATNegativeGovernedEndpointUse_has_separatorStatus} and \LeanName{cnfSATImageSeparatorBranch_of_negativeGovernedEndpointUse}.
\end{theorem}

\begin{proof}
By Theorem~\ref{thm:governed-endpoint-status-bivalence}, governed endpoint use is positive or separator. The assumption \(\neg\CnfPositiveEndpoint\) eliminates positive endpoint occupation. The only remaining governed endpoint status is separator status, and separator status is the SAT image-separator branch on the context-bound CNF model. This is the direct branch form formalized by \LeanName{cnfSATImageSeparatorBranch_of_negativeGovernedEndpointUse}.
\end{proof}

Consequently, the objection that the official raw negative may be theorem-bearing while avoiding image-separator occupation is not a remaining third route. Mere proposition-targeted negative theoremhood is harmless but does not by itself resolve the official endpoint. Once the same theorem is offered as the official negative endpoint resolution, Theorem~\ref{thm:official-negative-theoremhood-endpoint-use} makes it endpoint use; on the fixed governed finite CNF-SAT carrier, a negative endpoint is not positive; endpoint-status bivalence forces separator status; separator status occupies the image-separator branch; and separator endpoint-image occupation is the branch closed by the no-independent-discriminator route below.

\begin{definition}[SAT-local independence bridge]\label{def:sat-local-independence-bridge}
\(\SepIndBridge(model)\) denotes the SAT-local same-domain independence bridge: a separating candidate-status classifier for the context-bound CNF-SAT model is independent same-domain endpoint-status discrimination. In the Lean surface this is the predicate \(\CnfSepClassInd(model)\). It is not a source of standing and not an A+ premise; it records the SAT-local condition under which an image-separator branch produces the independent discriminator excluded by kernel-forced rigidity.
\end{definition}

\begin{definition}[Separator discriminator and no-independent closure]\label{def:sepdisc-nosepdisc}
For the context-bound CNF-SAT model, \(\SepDisc(model)\) denotes the independent same-domain endpoint-status discriminator induced by the SAT image-separator branch:
\[
  \SepDisc(model):=\CnfIndSepDisc(model).
\]
The predicate \(\NoSepDisc(model)\) denotes the kernel-forced no-independent-separator-classifier closure for that same model:
\[
  \NoSepDisc(model):=\CnfNoIndSepClass(model).
\]
It is used only at the endpoint-use consequence layer, and its operative content is that no independent same-domain separator discriminator survives on the fixed CNF-SAT endpoint carrier.
\end{definition}

\begin{theorem}[No-independent closure excludes the separator discriminator]\label{thm:no-sepdisc-excludes-sepdisc}
For the context-bound CNF-SAT model,
\[
  \NoSepDisc(model)\Rightarrow\neg\SepDisc(model).
\]
\end{theorem}

\begin{proof}
By Definition~\ref{def:sepdisc-nosepdisc}, \(\SepDisc(model)\) is the formal independent separator endpoint-status discriminator and \(\NoSepDisc(model)\) is the corresponding no-independent-separator-classifier closure. The latter is precisely the negation of the former at the fixed SAT endpoint carrier: independent same-domain separator discrimination is excluded. Thus \(\NoSepDisc(model)\Rightarrow\neg\SepDisc(model)\).
\end{proof}

\subsection{Endpoint-Image Occupation and the Single SAT Interface}\label{subsec:single-sat-interface}

The next results isolate the carrier-typing fact that blocks the remaining coequal-outcome objection. The broad P versus NP problem may be discussed as a two-outcome question. The proof spine here uses a narrower endpoint-image carrier: occupation of the positive encoded CNF-SAT candidate image. On that image, an occupant is an encoded polynomial-time SAT candidate with no SAT failure. Universal candidate failure is not a second occupant of that same image; it is non-occupation of the image.

\begin{theorem}[Single occupation interface of the fixed SAT endpoint image]\label{thm:single-occupation-interface-sat}
On the context-bound unrestricted finite CNF-SAT carrier, endpoint-image occupation is exactly positive candidate occupation:
\[
\begin{aligned}
  &\ClayCtx\wedge\FixedDom(R)\wedge\ModelOfCnf(model,R)\\
  &\qquad\Rightarrow
  \bigl(\SATEndpointImageOcc(model)\leftrightarrow\PositiveCandidateOcc(model)\bigr),\\
  &\PositiveCandidateOcc(model)\leftrightarrow
  \exists c\,\bigl(\PolyCand_{model}(c)\wedge
  \neg\FailsSAT_{model}(\Decode_{model}(c))\bigr).
\end{aligned}
\]
Consequently, every endpoint-image occupation act on this carrier is governance-equivalent to positive encoded non-failing candidate occupation.
\end{theorem}

\begin{proof}
The carrier fixed by \(\ClayCtx\), \(\FixedDom(R)\), and \(\ModelOfCnf(model,R)\) is the encoded deterministic polynomial-time CNF-SAT candidate image. By Definition~\ref{def:sat-endpoint-image-occupation}, occupation of that image is \(\SATEndpointImageOcc(model)\), and the only occupant type is an encoded candidate \(c\) that is polynomial-time in the model and whose decoded procedure has no SAT failure. The encoded positive endpoint exactness theorem, Theorem~\ref{thm:encoded-positive-endpoint-exactness}, identifies this occupant condition with the ordinary positive CNF-SAT endpoint. No finite benchmark, lower-bound template, oracle comparison, proof-complexity fragment, search trace, or universal failure statement is an object of this occupant type. Such objects may be proof support or obstruction data, but they do not occupy the image. Thus the occupation interface is exhausted by positive candidate occupation.
\end{proof}

\begin{corollary}[Candidate-image exclusion is not coequal image occupation]\label{cor:candidate-image-exclusion-not-coequal-occupation}
On the fixed context-bound SAT endpoint image, candidate-image exclusion is non-occupation of the image, not a second occupation interface:
\[
\begin{aligned}
  &\ClayCtx\wedge\FixedDom(R)\wedge\ModelOfCnf(model,R)\\
  &\qquad\wedge\CandImageExcl(model)
  \Rightarrow \neg\SATEndpointImageOcc(model).
\end{aligned}
\]
\end{corollary}

\begin{proof}
Candidate-image exclusion says that every encoded polynomial-time candidate in the fixed image fails SAT. By Theorem~\ref{thm:single-occupation-interface-sat}, image occupation requires an encoded polynomial-time candidate with no SAT failure. Universal candidate failure therefore does not provide a second occupant of the image; it asserts that the image has no occupant.
\end{proof}

\begin{theorem}[Negative coequal resolution requires an endpoint-interface change]\label{thm:negative-coequal-requires-interface-change}
A negative P versus NP outcome may be treated as a coequal outcome on the symmetric two-output carrier \(\OutcomePvNP\). It is not coequal occupation of the fixed positive CNF-SAT candidate-image carrier. Therefore, on the fixed SAT endpoint image,
\[
\begin{aligned}
  &\ClayCtx\wedge\FixedDom(R)\wedge\ModelOfCnf(model,R)\\
  &\qquad\wedge\CandImageExcl(model)
  \Rightarrow \neg\LawfulCoequalImageOcc_R(model).
\end{aligned}
\]
\end{theorem}

\begin{proof}
The symmetric outcome carrier of Definition~\ref{def:symmetric-outcome-carrier} declares two broad problem outcomes, \(\Pclass=\NPclass\) and \(\Pclass\neq\NPclass\). That is a legitimate external presentation of the open problem. The endpoint-image carrier used in this proof is different: by Theorem~\ref{thm:single-occupation-interface-sat}, its occupants are precisely encoded non-failing polynomial-time CNF-SAT candidates. Candidate-image exclusion supplies no such occupant. Treating the negative outcome as coequal therefore changes the endpoint interface from candidate-image occupation to symmetric outcome classification. On the fixed candidate-image carrier, the negative branch is non-occupation, not lawful coequal image occupation.
\end{proof}


\begin{theorem}[SAT separator discrimination is endpoint-admissibility-relevant]\label{thm:sat-sepdisc-endpoint-adm-relevant}
On the fixed context-bound CNF-SAT carrier, separator discrimination induced by global or local separator occupation is endpoint-admissibility-relevant:
\[
\begin{aligned}
  &\ClayCtx\wedge\FixedDom(R)\wedge\ModelOfCnf(model,R)\wedge\SepDisc(model)\wedge
  \bigl(\SepImgOcc\vee\LocalSepImgOcc(R,model)\bigr)\\
  &\qquad\Rightarrow \EndpointAdmRel_R(\SepDisc,\SepBare).
\end{aligned}
\]
\end{theorem}

\begin{proof}
Global separator occupation \(\SepImgOcc\) and local separator occupation \(\LocalSepImgOcc(R,model)\) are not support-level lower-bound facts. They are, respectively, global endpoint-resolution use and local exact-countercase use of candidate-image non-occupation on the fixed SAT carrier. When \(\SepDisc(model)\) is induced by either act, it determines that the negative separator branch may stand as the endpoint-facing result or countercase while positive candidate occupation is excluded. That determination controls reportability, downstream reuse, and admissible continuation of the endpoint-resolution act. By Definition~\ref{def:endpoint-admissibility-relevance}, it is endpoint-admissibility-relevant. The claim is not that every lower-bound theorem is admissibility-relevant; the claim is only that separator discrimination becomes admissibility-relevant when it is induced by global or local endpoint-resolving separator occupation.
\end{proof}

\begin{corollary}[A+ excludes the SAT separator discriminator as admissibility-relevant independent work]\label{cor:aplus-excludes-sat-sepdisc-admrel}
On the fixed context-bound CNF-SAT carrier, endpoint use of separator occupation cannot retain the separator discriminator:
\[
\begin{aligned}
  &\ClayCtx\wedge\FixedDom(R)\wedge\ModelOfCnf(model,R)
  \wedge\SepDisc(model)\wedge\EndpointUse_R(B)\\
  &\quad\wedge\EndpointAdmRel_R(\SepDisc,\SepBare)\wedge\IndDisc_R(\SepDisc,\SepBare)
  \Rightarrow \bot.
\end{aligned}
\]
\end{corollary}

\begin{proof}
This is Corollary~\ref{cor:aplus-excludes-independent-endpoint-admrel} applied to the SAT separator discriminator. The endpoint-admissibility-relevance hypothesis records that \(\SepDisc(model)\) changes endpoint standing, reportability, downstream reuse, or branch exclusion for the resolution act. The independence hypothesis records that it is not governance-equivalent to the admitted fixed-domain interface. The kernel-forced \(\Aplus\) consequence layer excludes exactly such independent admissibility-relevant same-domain endpoint governance.
\end{proof}


\subsection{Same-Domain Incompleteness Objections Do Not Reopen Endpoint Standing}\label{subsec:godel-stability-endpoint}

This subsection is defensive and local. It is not part of the SAT separator contradiction itself, and it does not use incompleteness theory to prove \(\PosRaw\). Its role is to close a later hostile-referee route: the attempt to rehabilitate the separator branch by appealing to a same-domain ``true but unprovable'', self-referential, semantic, infinitary, or extension-based incompleteness distinction after the endpoint-admissibility bridge has classified \(\SepDisc(model)\) as admissibility-relevant endpoint governance.

The companion non-instantiability theorem for G\"odel-style incompleteness objections distinguishes external metamathematical coding from a same-domain faithful coding architecture whose code judgments re-enter the original fixed domain as standing-bearing constraints. External coding may exist as bookkeeping. It matters to the present endpoint only if its code judgments change endpoint standing, reportability, licensed downstream reuse, or coherent endpoint trajectory. Once it does that, it is governance-relevant and falls under the same fixed-domain closure package. If it does not do that, it is endpoint-inert for the SAT closeout \cite{GodelNonInstantiability}.

\begin{definition}[Same-domain incompleteness endpoint objection]\label{def:same-domain-incompleteness-endpoint-objection}
A same-domain incompleteness endpoint objection to the SAT closeout is a G\"odel-style, truth/provability, self-referential, semantic, infinitary, reflection, or extension-based distinction that is claimed to re-enter the fixed CNF-SAT endpoint carrier and change at least one of the following: endpoint standing, reportability, licensed downstream reuse, admissible continuation of the endpoint-resolution act, or branch exclusion. If it changes none of these, it is endpoint-inert relative to the present proof.
\end{definition}

\begin{theorem}[Endpoint incompleteness relevance must be standing or trajectory relevance]\label{thm:endpoint-incompleteness-relevance-standing-trajectory}
Any same-domain incompleteness endpoint objection with force against the SAT closeout must survive into endpoint standing classification or coherent endpoint-trajectory structure. If it does not, it is endpoint-inert.
\end{theorem}

\begin{proof}
Endpoint force means that the objection changes what may stand, be reported, be reused, or continue as the fixed endpoint result. Those are exactly endpoint-admissibility-relevant features in Definition~\ref{def:endpoint-admissibility-relevance}. The companion G\"odel non-instantiability result proves the corresponding fixed-domain bridge: a closure-relevant distinction must survive into standing classification or coherent trajectory structure; external undecidability that does neither is closure-inert \cite{GodelNonInstantiability}. Applying that bridge to the endpoint-resolution act, an incompleteness objection that changes no endpoint standing and no coherent endpoint trajectory has no endpoint force.
\end{proof}

\begin{theorem}[Code-predicate endpoint objections are hidden gate work]\label{thm:code-predicate-endpoint-objections-hidden-gate-work}
If a code-predicate, proof predicate, provability predicate, or self-reference predicate is used to alter endpoint standing, reportability, licensed downstream reuse, or branch exclusion for the fixed SAT endpoint, then it performs hidden endpoint-gate work. It is therefore subject to the same no-independent-discriminator closure as every other same-domain endpoint-admissibility-relevant classifier.
\end{theorem}

\begin{proof}
A code-predicate that never changes endpoint standing or endpoint trajectory is endpoint-inert by Theorem~\ref{thm:endpoint-incompleteness-relevance-standing-trajectory}. If it does change those features, then it determines which endpoint acts may stand or continue. That is hidden gate work on the endpoint-resolution act. The companion G\"odel result proves that governance-relevant code-predicates on a fixed domain are exhausted by inert bookkeeping, extensional gate renaming, same-side refinement, or independent standing-bearing authorization, with no fifth role \cite{GodelNonInstantiability}. In the present endpoint setting, the latter non-inert roles are exactly endpoint-admissibility-relevant governance and are already controlled by Corollary~\ref{cor:aplus-excludes-independent-endpoint-admrel}.
\end{proof}

\begin{corollary}[Endpoint resolution is stable against same-domain incompleteness objections]\label{cor:endpoint-resolution-stable-against-incompleteness}
Once \(\EndpointUse_R(B)\), \(\FixedDom(R)\), and the kernel-forced \(\Aplus_R(B)\) consequence layer are in force for official endpoint resolution on the unrestricted finite CNF-SAT carrier, no same-domain G\"odel-style incompleteness architecture can reopen endpoint standing, restore reportability, or rehabilitate the separator branch.
\end{corollary}

\begin{proof}
Let an alleged same-domain incompleteness architecture be offered against the endpoint closeout. If it does not change endpoint standing, reportability, downstream reuse, admissible continuation, or branch exclusion, then it is endpoint-inert by Definition~\ref{def:same-domain-incompleteness-endpoint-objection}. If it does change one of those features, then by Theorem~\ref{thm:endpoint-incompleteness-relevance-standing-trajectory} it survives into standing or trajectory structure, and by Theorem~\ref{thm:code-predicate-endpoint-objections-hidden-gate-work} any code-predicate form of it is hidden endpoint-gate work. The companion route-exhaustion theorem for G\"odel-style objections shows that code-level readback, semantic or meta-foundational import, reflection or infinitary rule-power, second-equivalence routes, and richer same-domain extensions are each inert, extensional gate renaming, finite-witness collapse, explicit scope change, or inadmissible authority import \cite{GodelNonInstantiability}. None yields a new same-domain admissibility-relevant endpoint role compatible with \(\Aplus\). Thus no such architecture reopens endpoint standing or rehabilitates the separator branch.
\end{proof}

\begin{corollary}[No same-domain true-but-unprovable separator split]\label{cor:no-true-but-unprovable-separator-split}
There is no same-domain endpoint-admissibility-relevant split between a separator branch that is ``true but unprovable'' and a separator branch that is reportable, reusable, or endpoint-standing on the fixed CNF-SAT carrier.
\end{corollary}

\begin{proof}
Such a split would distinguish two endpoint-relevant statuses while preserving the same fixed carrier. If it changes no endpoint standing or trajectory status, it is endpoint-inert. If it does change endpoint standing or trajectory status, it is a same-domain admissibility-relevant refinement of the endpoint interface. The companion G\"odel non-instantiability theorem excludes same-domain admissibility-relevant ``true but unprovable'' positive-side splits in the unique admissible interior \cite{GodelNonInstantiability}; the endpoint-admissibility bridge applies the same standing/trajectory exhaustion to the SAT resolution act. Hence no such separator split survives.
\end{proof}

\begin{remark}[Scope of the G\"odel supplement]\label{rem:godel-supplement-scope}
This subsection does not challenge classical incompleteness theorems in their standard arithmetized settings, and it does not use the G\"odel paper as a hidden premise of the SAT proof. It closes only the same-domain endpoint objection. External arithmetizations, metamathematical codings, or studies of proof systems may remain legitimate. They affect the present endpoint only if they re-enter the fixed SAT carrier as standing-bearing constraints on endpoint acts; if they do, they are endpoint-admissibility-relevant and fall under the kernel-forced closure already used here. If they do not, they are inert relative to endpoint closure.
\end{remark}


\begin{theorem}[Single-interface typing is not sufficient for endpoint closure]\label{thm:single-interface-not-sufficient-for-endpoint-closure}
The single occupation interface of the fixed SAT endpoint image does not by itself exclude candidate-image non-occupation, lower-bound theorems, obstruction arguments, or negative complexity results. In particular, the manuscript does not use an implication of the form
\[
  \CandImageExcl(model)\Rightarrow \bot .
\]
The contradiction requires all of the following additional structure: official endpoint-resolution use promotes the support-level non-occupation content to \(\SepImgOcc\); separator occupation induces \(\SepDisc(model)\); and endpoint use of that same separator occupation activates \(\NoSepDisc(model)\), which excludes \(\SepDisc(model)\).
\end{theorem}

\begin{proof}
Theorem~\ref{thm:single-occupation-interface-sat} identifies the occupant type of the fixed SAT endpoint image. Corollary~\ref{cor:candidate-image-exclusion-not-coequal-occupation} then says that universal candidate failure is non-occupation of that image, not a coequal occupant of it. Those facts classify the role of lower-bound content; they do not make the lower-bound content contradictory. Candidate-image exclusion may remain proof support, ordinary theorem content, local obstruction, or carrier-diagnostic information. It enters the contradiction only after the additional endpoint-use steps are present: Theorem~\ref{thm:candidate-image-exclusion-to-separator-occupation} promotes candidate-image exclusion to \(\SepImgOcc\) only under official endpoint-resolution use; Theorem~\ref{thm:non-positive-status-work-induces-sepdisc} and Theorem~\ref{thm:image-separator-requires-independent-discriminator} yield \(\SepDisc(model)\); and Theorem~\ref{thm:no-independent-sat-separator-classifier} together with Theorem~\ref{thm:no-sepdisc-excludes-sepdisc} yields \(\neg\SepDisc(model)\). Thus the proof is non-explosive with respect to single-interface typing itself.
\end{proof}

\begin{definition}[Non-positive endpoint-status work on the SAT image]\label{def:non-positive-status-work-sat}
\[
  \NonPositiveStatusWork_R(model)
\]
means theorem-bearing endpoint-status work over the fixed context-bound CNF-SAT candidate image. It assigns non-occupation status to that image, preserves the same carrier and encoded candidate model, does not occupy the positive candidate interface, and is not proof support, bookkeeping, carrier shift, repair, or lawful coequal image occupation.
\end{definition}

\begin{definition}[Projected negative role content]\label{def:projected-negative-role-content}
\[
  \ProjectionRoleContent_R(B,model)
\]
means that the theorem-bearing endpoint-use act \(B\), under the accepted official CNF-SAT resolution route, has candidate-image non-occupation \(\CandImageExcl(model)\) as its endpoint-role content on the fixed unrestricted finite CNF-SAT carrier. This is stronger than the separate availability of \(\CandImageExcl(model)\) as support-level lower-bound content: it binds the endpoint-use act itself to the projected negative content supplied by the SAT carrier.
\end{definition}

\begin{lemma}[Projection binds the official negative act to candidate-image non-occupation]\label{lem:projection-binds-negative-act-content}
On the fixed context-bound CNF-SAT endpoint carrier, ordinary official negative resolution through the accepted CNF-SAT route binds the endpoint-use act to candidate-image non-occupation:
\[
\begin{aligned}
&\ClayCtx\wedge\FixedDom(R)\wedge\ModelOfCnf(model,R)
  \wedge\OfficialCNFSATRoute(R,model)\\
&\qquad\wedge\OrdOfficialRes_R(\neg\PosRaw)
  \Rightarrow \ProjectionRoleContent_R(\SepBare,model).
\end{aligned}
\]
\end{lemma}

\begin{proof}
By Theorem~\ref{thm:ordinary-official-resolution-is-endpoint-use}, ordinary official negative resolution is endpoint use of the negative branch rather than support-level theoremhood. By Theorem~\ref{thm:sat-projection-fixes-endpoint-role-content}, the accepted official CNF-SAT route sends the symmetric negative outcome \(\Pclass\neq\NPclass\) to \(\SepBare\), and sends \(\SepBare\) to candidate-image non-occupation \(\CandImageExcl(model)\) on the same context-bound carrier. Hence the endpoint-use act itself is the projected negative act whose endpoint-role content is candidate-image non-occupation. This is precisely \(\ProjectionRoleContent_R(\SepBare,model)\).
\end{proof}

\begin{theorem}[Exhaustiveness of endpoint-role classification on the projected SAT carrier]
\label{thm:projected-sat-role-exhaustiveness}
On the unrestricted finite CNF-SAT endpoint-image carrier, after the standard CNF-SAT projection is accepted for official resolution, no stable middle endpoint role remains between positive occupation, lawful coequal occupation, and independent non-positive status work. More precisely, if the context-bound carrier, official CNF-SAT resolution route, endpoint use of the act, the binding of that act to projected negative role content, and same-domain preservation are all fixed, then the act is non-positive endpoint-status work:
\[
\begin{aligned}
&\ClayCtx\wedge\FixedDom(R)\wedge\ModelOfCnf(model,R)
  \wedge\OfficialCNFSATRoute(R,model)\\
&\quad\wedge\EndpointUse_R(B)
  \wedge\ProjectionRoleContent_R(B,model)
  \wedge\CandImageExcl(model)\\
&\quad\wedge\neg\SATEndpointImageOcc(model)
  \wedge\neg\LawfulCoequalImageOcc_R(model)\\
&\qquad\Rightarrow \NonPositiveStatusWork_R(model).
\end{aligned}
\]
The hypothesis \(\EndpointUse_R(B)\) excludes support-only and bookkeeping readings. The hypothesis \(\ProjectionRoleContent_R(B,model)\) binds the endpoint-use act \(B\) to candidate-image non-occupation, so \(\CandImageExcl(model)\) is not merely a free-standing lower-bound support statement. The fixed-domain hypotheses exclude carrier shift, repair, and interface change.
\end{theorem}

\begin{proof}
The hypotheses place the act on the fixed unrestricted finite CNF-SAT candidate-image carrier and make the CNF-SAT correspondence the route of official endpoint resolution. Definition~\ref{def:projected-negative-role-content} supplies the missing binding: the endpoint-use act \(B\) itself has candidate-image non-occupation as its projected endpoint-role content. Thus \(\CandImageExcl(model)\) is not merely adjacent lower-bound content; it is the role content of the endpoint act being classified.

Theorem~\ref{thm:sat-projection-fixes-endpoint-role-content} gives the endpoint-role content of a negative resolution on that route, and Theorem~\ref{thm:single-occupation-interface-sat} gives the single occupation interface of the image, so positive occupation is exactly encoded non-failing polynomial-time candidate occupation. Corollary~\ref{cor:candidate-image-exclusion-not-coequal-occupation} and Theorem~\ref{thm:negative-coequal-requires-interface-change} exclude the reading of candidate-image exclusion as either positive occupation or coequal image occupation. Endpoint use excludes support-only and bookkeeping roles; fixed domain excludes carrier shift; act-time finality excludes repair. Thus the act is theorem-bearing, endpoint-resolving, same-domain, non-positive, non-coequal, non-support, non-bookkeeping, non-repair, and non-carrier-shifting.

By Theorem~\ref{thm:endpoint-discriminator-trichotomy}, endpoint-status work on a fixed domain is bookkeeping, governance-equivalent to the already admitted interface, independent same-domain work, or domain-shifting. The listed exclusions remove bookkeeping and domain shift. The single-interface theorem removes governance-equivalence to positive occupation because the projected role content assigns non-occupation of the positive candidate image. The coequal-image route is excluded by Theorem~\ref{thm:negative-coequal-requires-interface-change}. Therefore the remaining role is same-domain non-positive endpoint-status work on the fixed SAT image, exactly \(\NonPositiveStatusWork_R(model)\).
\end{proof}

\begin{remark}[Proof legitimacy does not neutralize endpoint role]
Ordinary theoremhood establishes that a proof act meets mathematical standards of justification. It does not determine the endpoint role that act performs once it is used as official resolution on a fixed endpoint image. On the projected SAT carrier, an official negative theorem has the endpoint-role content of candidate-image non-occupation. Proof legitimacy does not cancel, override, or reclassify that role. It only says that the act is theorem-bearing rather than heuristic or unsupported.
\end{remark}

\begin{remark}[Single-interface typing feeds the trichotomy]
The single-interface theorem is not an endpoint-closure theorem by itself. It supplies one premise of the trichotomy elimination: on the fixed positive CNF-SAT candidate-image carrier, the only image-occupation interface is encoded non-failing candidate occupation. Any official endpoint-resolving act that determines non-occupation of that same image is therefore non-positive status work. If it is not support, bookkeeping, repair, carrier shift, or coequal image occupation, the endpoint-discriminator trichotomy classifies it as independent same-domain endpoint-status work.
\end{remark}

\begin{theorem}[Non-positive status work induces the separator discriminator]\label{thm:non-positive-status-work-induces-sepdisc}
On the fixed SAT endpoint image, non-positive endpoint-status work is separator-discriminator work:
\[
\begin{aligned}
  &\ClayCtx\wedge\FixedDom(R)\wedge\ModelOfCnf(model,R)\\
  &\qquad\wedge\NonPositiveStatusWork_R(model)
  \Rightarrow \SepDisc(model).
\end{aligned}
\]
\end{theorem}

\begin{proof}
By Definition~\ref{def:non-positive-status-work-sat}, the act settles endpoint status on the same fixed encoded candidate image while not occupying the positive interface, not remaining support or bookkeeping, not shifting carrier, and not occupying a lawful coequal image role. Theorem~\ref{thm:single-occupation-interface-sat} supplies the single occupation interface, and Theorem~\ref{thm:negative-coequal-requires-interface-change} excludes coequal negative occupation on that image. Thus the act is an endpoint-status discriminator that is not bookkeeping, not governance-equivalent to positive candidate occupation, and not domain-shifting. By Theorem~\ref{thm:endpoint-status-work-criterion-independence}, it is independent same-domain endpoint-status governance. In the SAT image, such independent same-domain non-positive status work is exactly the separator discriminator \(\SepDisc(model)\).
\end{proof}


\begin{theorem}[Ordinary theorem-bearing non-occupation is not a fifth endpoint role after native lower-bound collapse]
\label{thm:ordinary-nonoccupation-not-fifth-role}
On the fixed context-bound SAT endpoint image, ordinary theorem-bearing non-occupation is not a role outside the endpoint-discriminator trichotomy. More precisely, if the official CNF-SAT route is accepted, the ordinary negative branch is used with official endpoint force, and that endpoint-use act is bound by the projection to candidate-image non-occupation, then the induced endpoint role is the separator discriminator:
\[
\begin{aligned}
&\ClayCtx\wedge\FixedDom(R)\wedge\ModelOfCnf(model,R)\\
&\quad\wedge\OfficialCNFSATRoute(R,model)\wedge\OrdThm(\SepBare)\\
&\quad\wedge\OrdOfficialRes_R(\neg\PosRaw)
  \wedge\ProjectionRoleContent_R(\SepBare,model)\\
&\quad\wedge\CandImageExcl(model)\wedge\neg\SATEndpointImageOcc(model)\\
&\quad\wedge\neg\LawfulCoequalImageOcc_R(model)
\Rightarrow \SepDisc(model).
\end{aligned}
\]
\end{theorem}

\begin{proof}
Ordinary theoremhood certifies proof legitimacy; it does not determine endpoint-interface role. By Theorem~\ref{thm:ordinary-official-resolution-is-endpoint-use}, ordinary official endpoint resolution is endpoint use rather than support or bookkeeping. The hypothesis \(\ProjectionRoleContent_R(\SepBare,model)\) binds that endpoint-use act to candidate-image non-occupation, so \(\CandImageExcl(model)\) is the content of the act being classified and not a merely adjacent lower-bound statement. Corollary~\ref{cor:candidate-image-exclusion-not-coequal-occupation} and Theorem~\ref{thm:negative-coequal-requires-interface-change} show that this content is not positive candidate occupation and not coequal image occupation. The fixed context excludes carrier shift, and official endpoint use excludes support-only and bookkeeping roles.

The expanded native lower-bound collapse, Theorem~\ref{thm:ordinary-native-lb-falsification-collapse}, has already exhausted ordinary lower-bound falsification into support, positive bridge-neutralization, bookkeeping, carrier/interface shift, or \(\SepDisc(model)\). After official endpoint use and projected role content are fixed, the support, bridge-neutralized, bookkeeping, and carrier-shift branches are unavailable. Thus the hypotheses of Theorem~\ref{thm:projected-sat-role-exhaustiveness}, with \(B=\SepBare\), are precisely the hypotheses of the alleged middle role. That theorem classifies the act as \(\NonPositiveStatusWork_R(model)\); Theorem~\ref{thm:non-positive-status-work-induces-sepdisc} then identifies that non-positive same-domain endpoint-status work on the SAT image with \(\SepDisc(model)\). Therefore ordinary theorem-bearing non-occupation supplies no fifth endpoint role.
\end{proof}


\begin{theorem}[No stable non-independent classification for projected official negative SAT resolution]
\label{thm:no-stable-nonindependent-classification-projected-sat}
On the projected unrestricted finite CNF-SAT endpoint carrier, there is no stable non-independent classification for ordinary theorem-bearing non-occupation once official endpoint use is fixed. More precisely, if ordinary mathematical practice uses the standard CNF-SAT endpoint route for official resolution, the endpoint-use act is bound to projected candidate-image non-occupation, and the act is theorem-bearing endpoint use on the fixed carrier, then the act cannot remain a non-independent ordinary endpoint role. The displayed endpoint-use hypothesis records official-resolution use rather than support-level lower-bound content, and the projected-role-content hypothesis records that \(\CandImageExcl(model)\) is the endpoint content of that act rather than a separate lower-bound statement. The resulting role is the SAT separator discriminator:
\[
\begin{aligned}
  &\ClayCtx\wedge\FixedDom(R)\wedge\ModelOfCnf(model,R)
  \wedge\OfficialCNFSATRoute(R,model)\\
  &\quad\wedge\EndpointUse_R(B)
  \wedge\ProjectionRoleContent_R(B,model)
  \wedge\CandImageExcl(model)\\
  &\quad\wedge\neg\SATEndpointImageOcc(model)
  \wedge\neg\LawfulCoequalImageOcc_R(model)
  \Rightarrow \SepDisc(model).
\end{aligned}
\]
\end{theorem}

\begin{proof}
Ordinary theorem-bearing status certifies the legitimacy of the proof act; it does not supply an endpoint-role exemption after official use and SAT projection are fixed. The projected-role-content hypothesis binds the act \(B\) to candidate-image non-occupation on the fixed carrier. The fixed context and model keep the act on the same unrestricted finite CNF-SAT candidate image. Official endpoint use excludes support-only and bookkeeping readings; act-time finality excludes repair; fixed-domain preservation excludes carrier shift. The single-interface theorem, Theorem~\ref{thm:single-occupation-interface-sat}, excludes positive candidate occupation, and Theorem~\ref{thm:negative-coequal-requires-interface-change} excludes lawful coequal image occupation on this carrier.

After these exclusions, the proposed middle classification is exactly the endpoint-deployed form of \(\LBNatSAT(R,model)\) after support, bookkeeping, positive bridge occupation, coequal occupation, repair, and carrier shift have been removed. The ordinary native lower-bound falsification collapse, Theorem~\ref{thm:ordinary-native-lb-falsification-collapse}, identifies that residue as \(\SepDisc(model)\). Equivalently, Theorem~\ref{thm:projected-sat-role-exhaustiveness} rules out the middle role by yielding \(\NonPositiveStatusWork_R(model)\), and Theorem~\ref{thm:non-positive-status-work-induces-sepdisc} identifies that role with \(\SepDisc(model)\).
\end{proof}

\begin{remark}[Ordinary practice supplies legitimacy, not a role exemption]
The theorem does not deny ordinary mathematical practice. It grants ordinary theorem-bearing legitimacy and then asks what endpoint role that legitimate act performs after official CNF-SAT projection. Once the act resolves the endpoint through candidate-image non-occupation on the fixed carrier, ordinary theoremhood cannot neutralize the endpoint-role classification. A stable non-independent middle role would have to preserve official endpoint force while avoiding support, bookkeeping, repair, carrier shift, positive occupation, coequal occupation, and independence. No such role remains.
\end{remark}


\begin{lemma}[Exhaustiveness of endpoint-role classification under official CNF-SAT resolution]
\label{lem:official-cnf-sat-resolution-exhaustiveness}
Let \(R\) be the official \(P\) versus \(NP\) endpoint context and let \(model\) be a context-bound adequate encoded CNF-SAT model. If ordinary mathematical practice performs official endpoint resolution of the negative branch through the standard unrestricted finite CNF-SAT correspondence, then the projected negative act has no stable non-independent endpoint classification on the fixed positive candidate-image carrier. In particular,
\[
\begin{aligned}
  &\ClayCtx\wedge\FixedDom(R)\wedge\ModelOfCnf(model,R)
  \wedge\OfficialCNFSATRoute(R,model)\\
  &\qquad\wedge\OrdOfficialRes_R(\Pclass\neq\NPclass)
  \Rightarrow \SepDisc(model).
\end{aligned}
\]
\end{lemma}

\begin{proof}
The hypothesis \(\OfficialCNFSATRoute(R,model)\) fixes the standard unrestricted finite CNF-SAT correspondence as the route of official endpoint resolution, not merely as a source of support-level lower-bound material. By Theorem~\ref{thm:sat-projection-faithfully-carries-outcomes}, the symmetric negative outcome is carried by that route to \(\SepBare\) and then to candidate-image exclusion \(\CandImageExcl(model)\) on the fixed encoded candidate image. By Theorem~\ref{thm:sat-projection-fixes-endpoint-role-content}, this is endpoint-role content on that carrier rather than an uninterpreted symmetric negative answer.

The hypothesis \(\OrdOfficialRes_R(\Pclass\neq\NPclass)\) says that the theorem-bearing negative act is used as the official resolution of the endpoint. By Theorem~\ref{thm:ordinary-official-resolution-is-endpoint-use}, such use is \(\EndpointUse_R\), and hence is not proof support, bookkeeping, or inert proposition-level theoremhood. Corollary~\ref{cor:official-negative-resolution-through-cnf-sat-nonpositive-work} then classifies the projected official negative resolution as \(\NonPositiveStatusWork_R(model)\): it determines non-occupation of the single positive candidate-image interface on the same carrier, is not positive occupation, is not coequal occupation of that image, and is not a carrier or interface shift. Theorem~\ref{thm:non-positive-status-work-induces-sepdisc} identifies that non-positive same-domain endpoint-status work with the SAT separator discriminator \(\SepDisc(model)\).
\end{proof}

\begin{remark}[Ordinary theorem-bearing status does not create a fifth role]
The lemma does not deny that the negative resolution may be ordinary theorem-bearing mathematics. It says that ordinary theoremhood is proof legitimacy, not endpoint-role exemption. Once ordinary theoremhood is used as official endpoint resolution through the CNF-SAT route, its endpoint-role content on the fixed candidate image is non-occupation. Since the act is endpoint-resolving, same-domain, non-positive, non-support, non-bookkeeping, non-repair, non-carrier-shifting, and non-coequal on that image, the endpoint-discriminator trichotomy leaves no stable middle classification. Treating the act as ``ordinary theorem-bearing negative determination'' while refusing the independent classification would require exactly the fifth endpoint role excluded by Theorem~\ref{thm:no-stable-nonindependent-classification-projected-sat}.
\end{remark}


\begin{theorem}[Exhaustion of remaining same-mode objections]
\label{thm:remaining-same-mode-objections-exhausted}
After the official CNF-SAT resolution route, SAT endpoint projection, single-interface typing, ordinary-resolution bridge, exact-complement reductio discipline, and endpoint-discriminator trichotomy are in force, any same-mode objection to the separator contradiction must deny at least one of the following three links:
\begin{enumerate}[label=(\roman*)]
\item the projection link pairing the symmetric outcome carrier with the SAT endpoint-image carrier, including
\[
\begin{gathered}
  \Pclass=\NPclass\leftrightarrow\PosRaw\leftrightarrow\SATEndpointImageOcc(model),\\
  \Pclass\neq\NPclass\leftrightarrow\SepBare\leftrightarrow\CandImageExcl(model).
\end{gathered}
\]
\item the trichotomy link classifying a theorem-bearing, endpoint-resolving, same-domain, non-positive, non-coequal, non-support, non-bookkeeping, non-carrier-shifting act as independent endpoint-status work;
\item the \(\Aplus\) closure link excluding independent same-domain endpoint-status discriminators on the fixed kernel-governed carrier.
\end{enumerate}
Appeal to ordinary theoremhood alone does not supply a fourth endpoint role.
\end{theorem}

\begin{proof}
Let a same-mode objection preserve the fixed carrier, context-bound candidate image, official endpoint use, the official CNF-SAT resolution route, and the accepted CNF-SAT projection. If the objection denies that the symmetric negative outcome becomes candidate-image exclusion on the SAT carrier, it attacks the projection link. If it accepts the projection but claims that official theorem-bearing non-occupation is not support, not bookkeeping, not carrier shift, not positive occupation, not coequal image occupation, and yet also not independent endpoint-status work, then it attacks the endpoint-discriminator trichotomy. If it accepts the independent discriminator classification but denies the contradiction with \(\NoSepDisc(model)\), it attacks the \(\Aplus\) no-independent-discriminator consequence. No further case remains: ordinary theoremhood is proof legitimacy, not endpoint-role classification. Theorem~\ref{thm:ordinary-nonoccupation-not-fifth-role} excludes ordinary theorem-bearing non-occupation as a fifth endpoint role, and Theorem~\ref{thm:no-stable-nonindependent-classification-projected-sat} excludes the final proposed non-independent middle classification after official use and SAT projection are fixed.
\end{proof}

\begin{theorem}[Separator occupation is non-positive status work]\label{thm:sepocc-is-non-positive-status-work}
In the fixed SAT context, separator endpoint-image occupation is non-positive endpoint-status work over the candidate image:
\[
  \ClayCtx\wedge\FixedDom(R)\wedge\ModelOfCnf(model,R)\wedge\SepImgOcc
  \Rightarrow \NonPositiveStatusWork_R(model).
\]
\end{theorem}

\begin{proof}
By Lemma~\ref{lem:separator-occupation-contains-candidate-image-exclusion}, \(\SepImgOcc\) contains \(\CandImageExcl(model)\). By Corollary~\ref{cor:candidate-image-exclusion-not-coequal-occupation}, that content is non-occupation of the fixed positive candidate image. By Definition~\ref{def:image}, \(\SepImgOcc\) is this non-occupation content after it is used as the official endpoint-resolving negative branch on the fixed carrier, so it is not support-only or bookkeeping and does not shift carrier. By Theorem~\ref{thm:negative-coequal-requires-interface-change}, it is not lawful coequal image occupation. Hence it is non-positive endpoint-status work on the fixed SAT image.
\end{proof}

\begin{theorem}[Endpoint-image exclusion is independent only under endpoint-resolving use]\label{thm:endpoint-image-exclusion-independent-only-under-endpoint-use}
Candidate endpoint-image exclusion is not automatically forbidden and is not automatically an independent same-domain discriminator. On the fixed unrestricted finite CNF-SAT carrier, \(\CandImageExcl(model)\) remains legitimate when it is used as proof support, a local obstruction, lower-bound lemma, ordinary theorem content not offered as the official endpoint resolution, or carrier-diagnostic information. It invokes the SAT-local independence bridge only when it is used to resolve the fixed endpoint negatively while preserving the same carrier and while not occupying the positive endpoint:
\[
\begin{aligned}
  &\ClayCtx\wedge\FixedDom(R)\wedge\ModelOfCnf(model,R)\\
  &\qquad\wedge\CandImageExcl(model)\wedge\SepImgOcc\wedge\EndpointUse_R(\SepImgOcc)\\
  &\Rightarrow \SepIndBridge(model).
\end{aligned}
\]
\end{theorem}

\begin{proof}
Proof-support use of candidate-image exclusion does not settle the official endpoint; it remains lawful lower-bound or obstruction material. Carrier-changing use is not the same endpoint. Bookkeeping use changes no endpoint status. The only remaining case in which candidate-image exclusion does endpoint work on the fixed carrier is endpoint-resolving use of the separator endpoint-image occupation \(\SepImgOcc\). By Theorem~\ref{thm:single-occupation-interface-sat}, the fixed SAT endpoint image has a single occupation interface: positive encoded non-failing candidate occupation. By Corollary~\ref{cor:candidate-image-exclusion-not-coequal-occupation}, candidate-image exclusion is non-occupation of that image, not coequal occupation. By Theorem~\ref{thm:negative-coequal-requires-interface-change}, treating the negative outcome as coequal would change to the symmetric outcome carrier rather than preserve the same endpoint image. Therefore endpoint-resolving use of \(\SepImgOcc\) is non-positive status work on the fixed image. Theorem~\ref{thm:sepocc-is-non-positive-status-work} gives \(\NonPositiveStatusWork_R(model)\), and Theorem~\ref{thm:non-positive-status-work-induces-sepdisc} identifies that work as the separator-discriminator role. Equivalently, the SAT-local independence bridge \(\SepIndBridge(model)\) is invoked only under endpoint-resolving use, not by lower-bound content alone.
\end{proof}


\begin{theorem}[SAT-local bridge instantiated by endpoint-resolving separator occupation]\label{thm:sat-local-bridge-instantiated}
In the fixed SAT context, the SAT-local independence bridge is invoked by candidate-image exclusion only when that content is in endpoint-resolving separator occupation:
\[
\begin{aligned}
  &\ClayCtx\wedge\FixedDom(R)\wedge\ModelOfCnf(model,R)\wedge\CandImageExcl(model)\\
  &\qquad\wedge\SepImgOcc\wedge\EndpointUse_R(\SepImgOcc)
  \Rightarrow \SepIndBridge(model).
\end{aligned}
\]
\end{theorem}

\begin{proof}
The model-binding clause fixes the encoded CNF-SAT carrier: arbitrary finite CNF instances over finite Boolean alphabets, evaluated by formula size, with admissible encodings preserving the same SAT endpoint image. The support-level lower-bound content is \(\CandImageExcl(model)\); it may remain legitimate proof support. The hypothesis \(\SepImgOcc\) identifies endpoint-image separator occupation, not mere lower-bound content, and \(\EndpointUse_R(\SepImgOcc)\) makes the endpoint-resolving use explicit. Theorem~\ref{thm:single-occupation-interface-sat} fixes the single occupation interface of the image; Corollary~\ref{cor:candidate-image-exclusion-not-coequal-occupation} classifies candidate-image exclusion as non-occupation; and Theorem~\ref{thm:negative-coequal-requires-interface-change} excludes coequal negative occupation on the same image. Therefore Theorem~\ref{thm:sepocc-is-non-positive-status-work} applies, and the endpoint-resolving separator occupation is non-positive status work over the same image. This is the SAT-local bridge represented in Lean by \(\CnfSepClassInd(model)\). The theorem therefore does not assert an independent discriminator from lower-bound content alone; it states the endpoint-use condition under which the image-separator theorem applies.
\end{proof}


\begin{theorem}[Separator endpoint occupation requires independent separator discriminator]\label{thm:image-separator-requires-independent-discriminator}
For the fixed SAT endpoint carrier,
\[
  \ClayCtx\wedge\FixedDom(R)\wedge\ModelOfCnf(model,R)\wedge\SepIndBridge(model)\wedge\SepImgOcc
  \Rightarrow
  \SepDisc(model).
\]
\end{theorem}

\begin{proof}
The single-interface result, Theorem~\ref{thm:single-occupation-interface-sat}, shows that the fixed SAT endpoint image is occupied only by positive encoded non-failing candidate occupation. Separator occupation is therefore non-positive status work by Theorem~\ref{thm:sepocc-is-non-positive-status-work}. The SAT-local independence bridge records that this non-positive same-domain status work is independent discriminator work on the encoded candidate image. With \(\SepIndBridge(model)\) and \(\SepImgOcc\) in force, the formal SAT operator theorem yields \(\SepDisc(model)\). This is the exact branch asymmetry: the positive branch is direct endpoint occupation, while endpoint-resolving separator occupation induces independent same-domain endpoint-status discrimination.
\end{proof}


\begin{theorem}[Official endpoint-resolving bare negative branch requires independent separator discriminator]\label{thm:bare-negative-requires-independent-discriminator}
For the fixed SAT endpoint carrier,
\[
\begin{aligned}
  &\ClayCtx\wedge\FixedDom(R)\wedge\ModelOfCnf(model,R)\wedge\SepBare\\
  &\qquad\wedge\OfficialEndpointResolution_R(\neg\PosRaw)
  \Rightarrow \SepDisc(model).
\end{aligned}
\]
\end{theorem}

\begin{proof}
By Theorem~\ref{thm:bare-negative-forces-image}, \(\SepBare\) gives \(\CandImageExcl(model)\), not endpoint occupation by itself. The additional endpoint-resolution premise is essential. Together with candidate-image exclusion, Theorem~\ref{thm:candidate-image-exclusion-to-separator-occupation} gives \(\SepImgOcc\). Theorem~\ref{thm:image-separator-endpoint-use} gives \(\EndpointUse_R(\SepImgOcc)\). With endpoint use explicit, Theorem~\ref{thm:sat-local-bridge-instantiated} gives \(\SepIndBridge(model)\), and Theorem~\ref{thm:image-separator-requires-independent-discriminator} gives \(\SepDisc(model)\). Thus the independent discriminator is required only for the official endpoint-resolving use of the bare negative branch, not for lower-bound content or ordinary negative theoremhood as such.
\end{proof}


\begin{theorem}[Separator endpoint-image occupation is endpoint use]\label{thm:image-separator-endpoint-use}
In the fixed SAT context, occupation of the separator endpoint-image branch is endpoint use:
\[
  \ClayCtx\wedge\FixedDom(R)\wedge\ModelOfCnf(model,R)\wedge\SepImgOcc
  \Rightarrow
  \EndpointUse_R(\SepImgOcc).
\]
\end{theorem}

\begin{proof}
The support-level lower-bound content is \(\CandImageExcl(model)\). It may appear as a lower-bound lemma, obstruction, or proof-support statement without endpoint use. The symbol \(\SepImgOcc\), by contrast, denotes endpoint-image separator occupation: candidate-image exclusion after it is used as the official endpoint-resolving negative branch on the fixed carrier. Thus the theorem does not promote arbitrary lower-bound content into endpoint governance. It only unpacks the occupation notation: when separator endpoint-image occupation is present on the fixed endpoint carrier, that presence is theorem-bearing endpoint use of the separator branch.
\end{proof}


\begin{theorem}[Branch-local no independent SAT separator classifier]\label{thm:no-independent-sat-separator-classifier}
For the fixed SAT endpoint carrier, separator endpoint-image occupation in endpoint use triggers the kernel-forced no-independent-discriminator consequence:
\[
  \ClayCtx\wedge\FixedDom(R)\wedge\ModelOfCnf(model,R)\wedge\EndpointUse_R(\SepImgOcc)
  \Rightarrow
  \NoSepDisc(model).
\]
\end{theorem}

\begin{proof}
Endpoint use is governed construction, hence it instantiates the kernel by Theorem~\ref{thm:kernel-necessity}. The fixed-domain A+ consequence layer then supplies the AMetric boundary, unique admissible interior, standing conservation, no-selector discipline, transport closure, and structural rigidity. On a fixed SAT endpoint carrier, an endpoint-status discriminator must be bookkeeping, governance-equivalent, domain-changing, or independent same-domain governance work by Theorem~\ref{thm:endpoint-discriminator-trichotomy}. The first two cases do not yield \(\SepDisc(model)\), the third is not the same SAT endpoint domain, and the fourth is excluded by the kernel-forced consequence layer. Thus separator endpoint-image use on the context-bound SAT model yields \(\NoSepDisc(model)\).
\end{proof}



\begin{corollary}[A+ consequences bind the SAT separator discriminator]
\label{cor:aplus-binds-sat-separator-discriminator}
On the fixed context-bound CNF-SAT carrier, once an endpoint act has been classified as the SAT separator discriminator, endpoint use binds that discriminator to the no-independent closure:
\[
\begin{aligned}
  &\ClayCtx\wedge\FixedDom(R)\wedge\ModelOfCnf(model,R)\wedge\EndpointUse_R(\SepImgOcc)\wedge\SepDisc(model)\\
  &\qquad\Rightarrow \NoSepDisc(model)\wedge\neg\SepDisc(model).
\end{aligned}
\]
\end{corollary}

\begin{proof}
The endpoint-use premise activates the kernel and the \(\Aplus\) consequence layer by Theorems~\ref{thm:kernel-necessity} and~\ref{thm:a-plus}. By Theorem~\ref{thm:sat-sepdisc-endpoint-adm-relevant}, a separator discriminator induced by endpoint-resolving separator occupation is endpoint-admissibility-relevant: it determines branch standing, reportability, downstream reuse, and exclusion on the fixed carrier. Thus the discriminator lies in the exact class of independent admissibility-relevant same-domain endpoint governance controlled by Proposition~\ref{prop:boundary-closure-forced} and Corollary~\ref{cor:aplus-excludes-independent-endpoint-admrel}. The SAT-local form of the no-independent closure is Theorem~\ref{thm:no-independent-sat-separator-classifier}, which yields \(\NoSepDisc(model)\). Theorem~\ref{thm:no-sepdisc-excludes-sepdisc} then gives \(\neg\SepDisc(model)\). Thus, after projection and trichotomy have identified the separator discriminator, ordinary endpoint use cannot retain that admissibility-relevant discriminator while exempting it from the same \(\Aplus\) closure.
\end{proof}

\begin{theorem}[Separator endpoint-image occupation impossible]\label{thm:no-image-sep}
In the fixed SAT context,
\[
  \ClayCtx\wedge\FixedDom(R)\wedge\ModelOfCnf(model,R)
  \Rightarrow \neg\SepImgOcc.
\]
\end{theorem}

\begin{proof}
Assume \(\ClayCtx\wedge\FixedDom(R)\wedge\ModelOfCnf(model,R)\) and suppose \(\SepImgOcc\). Theorem~\ref{thm:image-separator-endpoint-use} gives \(\EndpointUse_R(\SepImgOcc)\). Lemma~\ref{lem:separator-occupation-contains-candidate-image-exclusion} gives the support-level content \(\CandImageExcl(model)\). By Theorem~\ref{thm:single-interface-not-sufficient-for-endpoint-closure}, this support-level content alone is not the contradiction; the contradiction requires endpoint-resolving separator occupation. With endpoint use explicit, Theorem~\ref{thm:sat-local-bridge-instantiated} gives \(\SepIndBridge(model)\). Theorem~\ref{thm:image-separator-requires-independent-discriminator} gives \(\SepDisc(model)\), and Theorem~\ref{thm:sat-sepdisc-endpoint-adm-relevant} records that this discriminator is endpoint-admissibility-relevant for the resolution act rather than a mere descriptive label. Theorem~\ref{thm:no-independent-sat-separator-classifier} gives \(\NoSepDisc(model)\). By Theorem~\ref{thm:no-sepdisc-excludes-sepdisc}, \(\neg\SepDisc(model)\). Contradiction. Therefore \(\neg\SepImgOcc\).
\end{proof}



\paragraph{No fifth negative occupation.}
Thus the alleged third governed negative occupation is not a fifth case. If it has no endpoint-governance effect, it is bookkeeping. If it has endpoint-governance effect on the fixed carrier while being neither positive nor separator, it is hidden selector work or forbidden third-status refinement. If it depends on later reclassification of the same act, it is repair-sensitive middle status. If it avoids those consequences only by changing the admissibility-relevant invariant bundle, it is carrier or domain shift. Hence every purported governed negative endpoint occupation that is not separator either fails to be endpoint-bearing, violates the fixed-domain endpoint-status discipline, reopens repair, or leaves the fixed CNF-SAT carrier. No coherent fifth occupation remains.

\begin{theorem}[Official endpoint-resolving bare negative branch impossible]\label{thm:no-bare-sep}
In the fixed SAT context,
\[
\begin{aligned}
  &\ClayCtx\wedge\FixedDom(R)\wedge\ModelOfCnf(model,R)\\
  &\Rightarrow \neg\bigl(\SepBare\wedge\OfficialEndpointResolution_R(\neg\PosRaw)\bigr).
\end{aligned}
\]
\end{theorem}

\begin{proof}
Assume \(\ClayCtx\wedge\FixedDom(R)\wedge\ModelOfCnf(model,R)\) and suppose \(\SepBare\wedge\OfficialEndpointResolution_R(\neg\PosRaw)\). Theorem~\ref{thm:ordinary-official-negative-sat-resolution-nonpositive-status-work} shows that the ordinary official negative resolution has the non-positive endpoint-status role on the fixed candidate image. Theorem~\ref{thm:bare-negative-forces-image} gives the support-level content \(\CandImageExcl(model)\). By Theorem~\ref{thm:candidate-image-exclusion-to-separator-occupation}, the conjunction of candidate-image exclusion with official endpoint-resolution use gives \(\SepImgOcc\). But Theorem~\ref{thm:no-image-sep} excludes \(\SepImgOcc\) under the same fixed context. Contradiction. Thus the bare negative branch is impossible as the official endpoint-resolving negative branch. Non-endpoint lower-bound content and ordinary negative theoremhood remain governed by the non-explosiveness theorem; ordinary theoremhood becomes endpoint-role disciplined only when used as official endpoint resolution.
\end{proof}



\section{Route-Locus Exhaustion for SAT Separator Use}\label{sec:sat-route-locus-exhaustion}

\begin{definition}[Endpoint-resolving negative SAT route]\label{def:endpoint-resolving-negative-sat-route}
An endpoint-resolving negative SAT route is a proof trace, invariant, obstruction, lower-bound theorem, model comparison, classifier, or report that is used to make \(\SepBare\) the official negative resolution of the unrestricted finite CNF-SAT endpoint on the fixed carrier.
\end{definition}

\begin{definition}[SAT route loci]\label{def:sat-route-loci}
The SAT route-locus basis consists of the following loci:
\begin{center}
\small
\begin{tabularx}{\linewidth}{>{\raggedright\arraybackslash}p{0.10\linewidth}>{\raggedright\arraybackslash}X>{\raggedright\arraybackslash}p{0.30\linewidth}}
\toprule
Locus & SAT form & Disposition \\
\midrule
L1 & Carrier identity: unrestricted finite CNF versus fixed-width, promise, oracle, nonuniform, or benchmark carriers. & Carrier/domain shift unless bridged back. \\
L2 & Candidate image: uniform deterministic polynomial-time machine/bound pairs. & Adequate model by Subsection~\ref{subsec:encoded-candidate-model-adequacy}. \\
L3 & Encoding or presentation artifacts, duplicate codes, and representative choices. & Bookkeeping unless they change endpoint status, in which case selector/import failure. \\
L4 & Positive endpoint occupation by a deterministic polynomial-time SAT decider. & Direct positive branch. \\
L5 & Standard lower-bound normal form: every encoded polynomial-time candidate fails. & Proof support until endpoint-resolving; then enters image exclusion. \\
L6 & Candidate endpoint-image exclusion. & Endpoint-governing if used as official negative resolution. \\
L7 & Candidate-status classifier or theorem-level endpoint discriminator. & Routes to independent-discriminator closure. \\
L8 & Nonuniform, oracle, promise, average-case, parameterized, or restricted-fragment shift. & Not the same endpoint carrier. \\
L9 & Finite benchmark, empirical solver behavior, or heuristic hardness. & Support only; not endpoint occupation. \\
L10 & Repair or later reclassification of an earlier endpoint act. & Violates act-time finality. \\
L11 & Selector or ranking over candidates, encodings, search depths, or benchmarks. & AMetric no-selector violation if endpoint-bearing. \\
L12 & Endpoint-status governance outside positive occupation. & Forbidden independent same-domain authority. \\
L13 & Ordinary native lower-bound falsification: lower-bound theorem offered as endpoint defeat while refusing separator/discriminator classification. & Collapses by Theorem~\ref{thm:ordinary-native-lb-falsification-collapse} into support/bookkeeping, positive bridge-neutralization, carrier shift, or \(\SepDisc(model)\). \\
\bottomrule
\end{tabularx}
\end{center}
\end{definition}

\begin{theorem}[SAT route-locus exhaustion]\label{thm:sat-route-locus-exhaustion}
Every same-carrier endpoint-resolving negative SAT route enters at least one of L1--L13. If it enters none, it does no endpoint work and cannot resolve the official endpoint.
\end{theorem}

\begin{proof}
A negative SAT route that resolves the endpoint must either alter the carrier, quantify over the polynomial-time candidate image, depend on encodings or representatives, deny positive endpoint occupation, express lower-bound normal form, exclude the candidate endpoint image, classify candidate status, shift to nonuniform/oracle/promise scope, rely on finite evidence, repair a prior report, select among candidates, govern endpoint status, or claim ordinary native lower-bound falsification outside the separator taxonomy. These cases cover the ways a same-carrier negative route can do endpoint work. L1 and L8 are carrier shifts; L3 and L11 are bookkeeping or selector failures; L5 and L9 are proof support until endpoint-resolving; L6 and L7 route to endpoint governance; L10 violates irreversibility; L12 is the forbidden same-domain discriminator case; and L13 is collapsed by Theorem~\ref{thm:ordinary-native-lb-falsification-collapse}. If none of the loci is entered, the route preserves carrier, candidate image, encoding status, branch slot, lower-bound status, image occupation, classifier status, auxiliary scope, evidential status, act-time status, selector status, endpoint-status governance, and native-lower-bound use status. No endpoint-relevant work remains.
\end{proof}

\begin{corollary}[No surviving same-carrier separator route]\label{cor:no-surviving-sat-separator-route}
No same-carrier endpoint-resolving negative SAT route survives the route-locus audit outside the image-separator route already excluded by Theorem~\ref{thm:no-image-sep}.
\end{corollary}

\begin{proof}
By Theorem~\ref{thm:sat-route-locus-exhaustion}, every route enters L1--L13 or does no endpoint work. The non-endpoint routes are proof support or bookkeeping. The carrier-shift routes do not preserve the official endpoint. Repair and selector routes violate irreversibility or AMetricity. The endpoint-governing routes enter candidate-image exclusion, theorem-level discrimination, or independent-discriminator closure and are excluded by the image-separator argument. Therefore no additional same-carrier endpoint-resolving separator route survives.
\end{proof}

\section{Endpoint-Image Resolution}\label{sec:final}

\begin{theorem}[Bare SAT endpoint complementarity]\label{thm:bare-complement}
At the raw unrestricted finite CNF-SAT layer,
\[
  \SepBare\leftrightarrow\neg\PosRaw.
\]
\end{theorem}

\begin{proof}
By Definition~\ref{def:bare}, \(\PosRaw:=\CnfSATP\), the assertion that arbitrary finite CNF-SAT is decidable by a deterministic polynomial-time machine. The bare negative branch \(\SepBare\) is the official raw complement branch: failure of polynomial-time decidability for that same unrestricted finite CNF-SAT carrier. Definition~\ref{def:bounded-cnf-clarification} prevents any shift to a restricted CNF fragment. Therefore the raw branches are exact complements on the same carrier.
\end{proof}

\begin{theorem}[Bare endpoint bivalence]\label{thm:bare-bivalence}
\[
  \PosRaw\vee\SepBare.
\]
\end{theorem}

\begin{proof}
Classical bivalence gives \(\PosRaw\vee\neg\PosRaw\). By Theorem~\ref{thm:bare-complement}, \(\neg\PosRaw\) is equivalent to \(\SepBare\). Hence \(\PosRaw\vee\SepBare\).
\end{proof}

\begin{theorem}[Candidate-image carrier adequacy for the official raw split]\label{thm:candidate-image-carrier-adequacy-raw-split}
On the context-bound unrestricted finite CNF-SAT endpoint carrier, the positive candidate image represents the official raw split exactly:
\[
\begin{aligned}
  &\ClayCtx\wedge\FixedDom(R)\wedge\ModelOfCnf(model,R)\\
  &\qquad\Rightarrow
  \bigl[(\PosRaw\leftrightarrow\SATEndpointImageOcc(model))
  \wedge(\SepBare\leftrightarrow\CandImageExcl(model))\bigr].
\end{aligned}
\]
Thus the positive candidate-image carrier is not a partial or one-sided surrogate. The positive branch is image occupation, and the bare negative branch is image non-occupation on the same exact encoded candidate image.
\end{theorem}

\begin{proof}
For the positive half, Definition~\ref{def:bare} identifies \(\PosRaw\) with \(\CnfSATP\). The encoded positive endpoint exactness theorem, Theorem~\ref{thm:encoded-positive-endpoint-exactness}, identifies \(\CnfSATP\) with the existence of an encoded polynomial-time candidate whose decoded procedure has no SAT failure. The single-interface theorem, Theorem~\ref{thm:single-occupation-interface-sat}, identifies that occupant condition with \(\SATEndpointImageOcc(model)\). Hence \(\PosRaw\leftrightarrow\SATEndpointImageOcc(model)\).

For the negative half, Theorem~\ref{thm:bare-complement} gives \(\SepBare\leftrightarrow\neg\PosRaw\). The encoded negative endpoint exactness theorem, Theorem~\ref{thm:encoded-negative-endpoint-exactness}, identifies \(\neg\CnfSATP\) with universal failure of every encoded polynomial-time candidate on the fixed candidate image. That universal failure is exactly \(\CandImageExcl(model)\) by Definition~\ref{def:image} and Theorem~\ref{thm:standard-lower-bound-candidate-image-exclusion}. Therefore \(\SepBare\leftrightarrow\CandImageExcl(model)\). Both equivalences are on the same context-bound model and the same unrestricted finite CNF-SAT carrier, so the carrier represents the official raw split without changing endpoint interface.
\end{proof}



\subsection{Projection and Ordinary Practice}\label{subsec:projection-ordinary-practice}

Ordinary mathematical practice may present \(P\) versus \(NP\) symmetrically. The present argument does not deny the legitimacy of that presentation. It claims only that once the standard unrestricted finite CNF-SAT endpoint route is used for official resolution, the symmetric statement is decomposed on the fixed carrier into candidate-image occupation and candidate-image non-occupation. The negative resolution remains ordinary theoremhood as a matter of proof legitimacy, but its endpoint-role content on the SAT carrier is non-occupation. Accepting the CNF-SAT correspondence for support-level lower-bound work does not trigger this classification; accepting it as the route of official endpoint resolution does.

\begin{theorem}[SAT projection faithfully carries both symmetric outcomes]\label{thm:sat-projection-faithfully-carries-outcomes}
On the context-bound unrestricted finite CNF-SAT endpoint carrier, use of the standard CNF-SAT route gives the faithful projection
\[
\begin{gathered}
  \OfficialCNFSATRoute(R,model)\wedge\ClayCtx\wedge\FixedDom(R)\wedge\ModelOfCnf(model,R)\\
  \Rightarrow
  \bigl[(\Pclass=\NPclass\leftrightarrow\PosRaw\leftrightarrow\SATEndpointImageOcc(model))\\
  \wedge(\Pclass\neq\NPclass\leftrightarrow\SepBare\leftrightarrow\CandImageExcl(model))\bigr].
\end{gathered}
\]
\end{theorem}

\begin{proof}
The Cook--Levin/Karp endpoint correspondence identifies \(\Pclass=\NPclass\) with deterministic polynomial-time decidability of arbitrary finite CNF-SAT, which is \(\PosRaw\). The carrier-adequacy theorem, Theorem~\ref{thm:candidate-image-carrier-adequacy-raw-split}, identifies \(\PosRaw\) with \(\SATEndpointImageOcc(model)\) on the same context-bound candidate image. Negating the same endpoint gives \(\Pclass\neq\NPclass\), which corresponds to \(\SepBare\), and Theorem~\ref{thm:candidate-image-carrier-adequacy-raw-split} identifies \(\SepBare\) with \(\CandImageExcl(model)\). Thus the symmetric outcome carrier and the SAT endpoint-image carrier are not rival interfaces once the CNF-SAT route is used for official resolution; the latter is the fixed endpoint-image projection of the former.
\end{proof}

\begin{theorem}[SAT projection fixes endpoint-role content]\label{thm:sat-projection-fixes-endpoint-role-content}
If the standard CNF-SAT correspondence is used as the route of official resolution for \(P\) versus \(NP\) on the unrestricted finite CNF-SAT carrier, then the symmetric outcomes are evaluated through the fixed CNF-SAT endpoint-image carrier. Under that projection, \(\Pclass=\NPclass\) has the endpoint-role content of positive candidate-image occupation, and \(\Pclass\neq\NPclass\) has the endpoint-role content of candidate-image non-occupation. Consequently, a negative resolution cannot coherently claim to resolve the official endpoint through the CNF-SAT route while treating its content as merely an uninterpreted symmetric negative answer rather than as non-occupation on the fixed positive candidate-image carrier.
\end{theorem}

\begin{proof}
By Definition~\ref{def:official-cnf-sat-resolution-route}, \(\OfficialCNFSATRoute(R,model)\) fixes the CNF-SAT endpoint carrier as the route of official resolution, not merely as a support-level tool. Theorem~\ref{thm:sat-projection-faithfully-carries-outcomes} then carries both symmetric outcomes to their SAT endpoint-image roles. The positive outcome is \(\SATEndpointImageOcc(model)\), i.e. positive encoded non-failing candidate occupation. The negative outcome is \(\CandImageExcl(model)\), i.e. candidate-image non-occupation. Refusing that role while still claiming official resolution through CNF-SAT would accept the correspondence for endpoint force while refusing what the corresponding fixed image represents. That is not same-route endpoint use; it is either support-level use of CNF-SAT, which does not resolve the endpoint, or a change of endpoint interface.
\end{proof}

\begin{corollary}[Official negative resolution through CNF-SAT is non-positive status work]\label{cor:official-negative-resolution-through-cnf-sat-nonpositive-work}
On the fixed context-bound carrier,
\[
\begin{gathered}
  \ClayCtx\wedge\FixedDom(R)\wedge\ModelOfCnf(model,R)\wedge
  \OfficialCNFSATRoute(R,model)\\
  \wedge\OrdOfficialRes_R(\Pclass\neq\NPclass)
  \Rightarrow \NonPositiveStatusWork_R(model).
\end{gathered}
\]
\end{corollary}

\begin{proof}
By Theorem~\ref{thm:sat-projection-fixes-endpoint-role-content}, an official negative resolution through the CNF-SAT route has endpoint-role content \(\CandImageExcl(model)\) on the fixed SAT endpoint image. Since the use is ordinary official endpoint resolution, Theorem~\ref{thm:ordinary-official-resolution-is-endpoint-use} makes it endpoint use rather than support or bookkeeping. Corollary~\ref{cor:candidate-image-exclusion-not-coequal-occupation} and Theorem~\ref{thm:negative-coequal-requires-interface-change} show that this content is non-occupation of the positive candidate image and not coequal image occupation. The fixed context excludes carrier shift. Hence the act is non-positive endpoint-status work in the sense of Definition~\ref{def:non-positive-status-work-sat}.
\end{proof}

\begin{remark}[Projection is not optional once used for official resolution]\label{rem:projection-not-optional-once-used}
The manuscript does not require all mathematical discussion of \(P\) versus \(NP\) to pass through the positive candidate-image carrier. It requires only this: if the CNF-SAT correspondence is the route by which the endpoint is officially resolved, then the branch roles are read through that correspondence. Accepting the CNF-SAT correspondence merely as a tool for lower-bound arguments, examples, or support does not trigger this classification. The classification is triggered only when the correspondence is used as the route of official endpoint resolution.
\end{remark}

\begin{definition}[Truth branch, epistemic status, and endpoint outcome]
\label{def:truth-epistemic-endpoint-outcome}
For the fixed SAT endpoint carrier, three notions are kept separate.
A raw truth branch is one side of the proposition-level split \(\PosRaw\vee\SepBare\). An epistemic or review status records whether a community, proof attempt, or manuscript stage has established either branch. An endpoint outcome is a branch standing as the resolved endpoint of the official evaluation. Thus unresolved, unknown, open, deferred, or under-review status may describe a state of inquiry, but it is not an endpoint outcome and it is not a third truth branch.
\end{definition}

\begin{theorem}[Unresolved status is not an endpoint outcome]
\label{thm:unresolved-not-endpoint-outcome}
On the fixed SAT endpoint carrier, unresolved or open status is an epistemic status of inquiry, not an endpoint outcome:
\[
  \UnresolvedEvalStatus(R,model)\Rightarrow
  \neg\EndpointOutcome_R(\PosRaw)\wedge\neg\EndpointOutcome_R(\SepBare).
\]
This theorem does not assert that an official endpoint evaluation is already endpoint-complete. It only prevents unresolved or open status from being treated as a third endpoint branch.
\end{theorem}

\begin{proof}
An endpoint outcome must be a branch that can stand as the settled endpoint of the fixed raw split. The status ``unresolved'' does not assert positive image occupation and does not assert candidate-image non-occupation. It records lack of settlement or a state of inquiry. Therefore it does not occupy either endpoint branch. Treating unresolved status as a third endpoint outcome would change the official two-branch target into a three-status epistemic ledger. That is not the P versus NP endpoint and not the endpoint-image carrier under audit.
\end{proof}

\begin{definition}[Official SAT endpoint-closure evaluation]
\label{def:official-sat-endpoint-evaluation}
\(\OfficialSATEval(R,model)\) means that the official P versus NP endpoint is being evaluated on the fixed unrestricted finite CNF-SAT carrier through the exact raw split represented by the context-bound candidate image, in the endpoint-closure proof class of this manuscript. It fixes the target, carrier, raw branch slots, and endpoint-counterforce discipline: a raw branch can defeat the endpoint-closure proof only by having endpoint counterforce on the fixed carrier. It is not endpoint-complete evaluation, not a truth assumption, not proof availability, and not a rule that turns a merely hypothetical branch assumption into a theorem.
\end{definition}

\begin{definition}[Official SAT endpoint outcomes]
\label{def:official-sat-endpoint-outcomes}
On the fixed SAT endpoint evaluation, \(\PosEndpointOutcome_R\) means that the positive branch stands as the endpoint outcome. It is an endpoint-outcome status and entails \(\PosRaw\). The negative endpoint outcome \(\NegEndpointOutcome_R\) means that the bare negative branch is not merely a hypothetical case in a proof by cases, but is the branch standing as the official negative endpoint outcome. In formulas, the negative endpoint outcome carries the conjunctive endpoint role
\[
  \NegEndpointOutcome_R \Rightarrow \SepBare\wedge\OfficialEndpointResolution_R(\neg\PosRaw).
\]
The notation distinguishes a raw branch assumption from a branch that has endpoint standing as the target outcome.
\end{definition}

\begin{theorem}[Positive endpoint outcome gives the positive branch]
\label{thm:positive-outcome-gives-posraw}
On the fixed SAT endpoint evaluation,
\[
  \PosEndpointOutcome_R\Rightarrow\PosRaw.
\]
\end{theorem}

\begin{proof}
By Definition~\ref{def:official-sat-endpoint-outcomes}, \(\PosEndpointOutcome_R\) is the endpoint-outcome status in which the positive branch is the resolved endpoint outcome. The positive branch of the raw split is exactly \(\PosRaw\). Therefore \(\PosRaw\) follows.
\end{proof}

\begin{lemma}[Negative endpoint outcome carries endpoint-resolution use]
\label{lem:official-evaluation-negative-branch-endpoint-use}
On the fixed context-bound SAT carrier,
\[
\begin{aligned}
  &\ClayCtx\wedge\FixedDom(R)\wedge\ModelOfCnf(model,R)\wedge\NegEndpointOutcome_R\\
  &\qquad\Rightarrow \SepBare\wedge\OfficialEndpointResolution_R(\neg\PosRaw).
\end{aligned}
\]
\end{lemma}

\begin{proof}
This is the unpacking of the negative endpoint-outcome role in Definition~\ref{def:official-sat-endpoint-outcomes}. The premise is not a raw hypothetical assumption of \(\SepBare\). It says that the negative branch has endpoint standing as the official negative endpoint outcome. Such an outcome must contain the bare negative branch and its official endpoint-resolution use. If the negative proposition is present only as support, background theoremhood, a local lower-bound lemma, or a hypothetical branch in a proof by cases, then \(\NegEndpointOutcome_R\) is not in force.
\end{proof}


\begin{definition}[Endpoint counterforce]
\label{def:endpoint-counterforce}
For a fixed endpoint context, \(\EndpointCounterforce_R(B)\) means that the raw branch \(B\) is asserted with force to defeat or displace the endpoint closure on the same fixed carrier. Endpoint counterforce is not raw truth, not mere proof support, and not epistemic open status. It is the endpoint-facing force a branch must have in order to count against the closure of the official endpoint. A bare hypothetical assumption has no such force merely because it is assumed. When the assumed branch is introduced specifically as the exact complement countercase in a reductio or proof-by-cases subproof for the official endpoint split, it has local derivational standing as that countercase; this local role is recorded separately below.
\end{definition}



\begin{definition}[Local endpoint countercase assumption]
\label{def:local-countercase-assumption}
For the fixed SAT endpoint split, \(\LocalCountercase_R(\SepBare)\) means that \(\SepBare\) has been introduced inside a reductio or proof-by-cases subproof as the live exact complement countercase to \(\PosRaw\) for the official endpoint. This is local derivational standing, not theoremhood of \(\SepBare\), not official endpoint resolution, and not a truth-to-standing principle.

\end{definition}

\begin{definition}[Exact-complement reductio assumption]
\label{def:exact-complement-reductio-assumption}
For the fixed SAT endpoint split, \(\ReductioCountercase_R(\SepBare;\PosRaw)\) is the proof-act annotation generated when the sole local assumption \([\SepBare]^i\) is opened as the live exact-complement assumption in a reductio or proof-by-cases subproof whose target is \(\PosRaw\). This is an ordinary local proof role, not an additional object-level premise: it is not theoremhood of \(\SepBare\), not official negative endpoint resolution, not global endpoint outcome, and not a truth-to-standing principle. It records that the assumption \([\SepBare]^i\) is doing reductio work against the endpoint target rather than being an inert background supposition.
\end{definition}

\begin{lemma}[Exact-complement reductio annotation rule]
\label{lem:exact-complement-reductio-annotation-rule}
In a reductio or proof-by-cases subproof for \(\PosRaw\), opening the sole local assumption \([\SepBare]^i\) as the exact raw complement of the target generates the proof-act annotation
\[
  [\SepBare]^i_{\mathrm{exact\ complement\ for}\ \PosRaw}
  \quad\leadsto\quad
  \ReductioCountercase_R(\SepBare;\PosRaw).
\]
The annotation is not a second object-level hypothesis and is not a conjunct added to \(\SepBare\). It is the local role of the discharged assumption \([\SepBare]^i\) in the endpoint reductio subproof.
\end{lemma}

\begin{proof}
A reductio subproof has an object-level assumption and a proof role for that assumption. When the assumption is the exact raw complement of the target being proved, and the subproof is opened to test whether that target can be denied, ordinary proof practice already treats the assumption as the live countercase to the target. The notation \(\ReductioCountercase_R(\SepBare;\PosRaw)\) records only that proof role. It does not add content beyond the assumed proposition \(\SepBare\); therefore the later contradiction discharges \([\SepBare]^i\), not a strengthened assumption \(\SepBare\wedge\ReductioCountercase_R(\SepBare;\PosRaw)\).
\end{proof}


\begin{theorem}[Exact-complement annotation is admissible for reductio discharge]
\label{thm:exact-complement-annotation-discharge}
In the exact-complement reductio subproof for the fixed SAT endpoint, contradiction derived after the proof-act annotation discharges the original object-level assumption \([\SepBare]^i\). In sequent form,
\[
  \left([\SepBare]^i_{\mathrm{exact\ complement\ for}\ \PosRaw}\vdash_R \bot\right)
  \Rightarrow \vdash_R \neg\SepBare.
\]
The discharged formula is \(\SepBare\), not \(\SepBare\wedge\ReductioCountercase_R(\SepBare;\PosRaw)\).
\end{theorem}

\begin{proof}
The only object-level assumption opened in the subproof is \([\SepBare]^i\). Lemma~\ref{lem:exact-complement-reductio-annotation-rule} supplies \(\ReductioCountercase_R(\SepBare;\PosRaw)\) as a proof-act annotation of that assumption, not as an additional hypothesis in the object language. All downstream local roles are therefore rules governing the exact-complement reductio act: local countercase standing, endpoint counterforce, and local separator occupation are generated from the role of the open assumption in the fixed endpoint proof. Since no additional object-level premise is introduced, a contradiction obtained in the annotated local subproof is a contradiction under the sole discharged assumption \([\SepBare]^i\). Ordinary reductio discharge therefore yields \(\neg\SepBare\). This theorem is a discharge rule for the proof annotation; it does not assert that raw truth implies endpoint standing.
\end{proof}


\begin{theorem}[Exact-complement reductio assumptions have local endpoint-countercase standing]
\label{thm:reductio-countercase-yields-local-countercase}
In the fixed official SAT endpoint-closure proof, an exact-complement reductio assumption carries local endpoint-countercase standing:
\[
\begin{aligned}
  &\ClayCtx\wedge\FixedDom(R)\wedge\ModelOfCnf(model,R)\wedge\OfficialSATEval(R,model)\\
  &\qquad\wedge\ReductioCountercase_R(\SepBare;\PosRaw)\Rightarrow \LocalCountercase_R(\SepBare).
\end{aligned}
\]
\end{theorem}

\begin{proof}
By Lemma~\ref{lem:exact-complement-reductio-annotation-rule}, \(\ReductioCountercase_R(\SepBare;\PosRaw)\) is not a strengthened object-level assumption. It is the local annotation of the sole assumption \([\SepBare]^i\) when that assumption is opened as the live exact-complement countercase to \(\PosRaw\). A reductio assumption may otherwise be inert background syntax; in that case it is not a reductio countercase for the endpoint target. When the exact-complement annotation is present, the assumption has precisely the local endpoint-countercase standing recorded by \(\LocalCountercase_R(\SepBare)\). This theorem therefore does not retype every occurrence of \(\SepBare\); it applies only to the discharged exact-complement assumption as used in the endpoint reductio subproof.
\end{proof}

\begin{theorem}[Exact complement countercase has local endpoint counterforce]
\label{thm:local-countercase-has-endpoint-counterforce}
In the fixed official SAT endpoint-closure evaluation, a local assumption of the exact complement branch has endpoint counterforce in the subproof:
\[
\begin{aligned}
  &\ClayCtx\wedge\FixedDom(R)\wedge\ModelOfCnf(model,R)\wedge\OfficialSATEval(R,model)\\
  &\qquad\wedge\LocalCountercase_R(\SepBare)\Rightarrow \EndpointCounterforce_R(\SepBare).
\end{aligned}
\]
\end{theorem}

\begin{proof}
A local countercase assumption is not a theorem of the negative branch and is not official endpoint resolution. It is nevertheless a governed local proof act. In a reductio or proof-by-cases subproof for the official endpoint split, the exact complement branch is introduced precisely as the branch that would block the positive endpoint if it survived the subproof. That is the local endpoint-facing force recorded by \(\EndpointCounterforce_R(\SepBare)\). The implication therefore uses ordinary local assumption standing, not the false principle that raw truth implies theorem standing.
\end{proof}

\begin{definition}[Local separator occupation]
\label{def:local-separator-occupation}
For the fixed SAT endpoint split, \(\LocalSepImgOcc(R,model)\) means that candidate-image exclusion is used locally inside the countercase subproof as the endpoint-defeating negative branch. It is local separator occupation, not global endpoint outcome, not official negative endpoint resolution, and not a proof that the negative branch is true. It records the endpoint-facing role that the exact complement assumption takes while the reductio subproof is live.
\end{definition}

\begin{theorem}[Local endpoint counterforce induces local separator occupation]
\label{thm:local-counterforce-induces-local-separator-occupation}
On the fixed context-bound SAT carrier, local endpoint counterforce by the bare negative branch induces local separator occupation:
\[
\begin{aligned}
  &\ClayCtx\wedge\FixedDom(R)\wedge\ModelOfCnf(model,R)\wedge\OfficialSATEval(R,model)\\
  &\qquad\wedge\EndpointCounterforce_R(\SepBare)\Rightarrow \LocalSepImgOcc(R,model).
\end{aligned}
\]
\end{theorem}

\begin{proof}
Endpoint counterforce is the endpoint-facing force of the exact raw complement branch inside the official endpoint proof. By the carrier-adequacy theorem, \(\SepBare\) is represented on the context-bound candidate image as \(\CandImageExcl(model)\). Local endpoint counterforce means that this non-occupation content is being used inside the subproof as the endpoint-defeating negative branch, not merely as lower-bound support or raw syntax. That is exactly the local separator-occupation role \(\LocalSepImgOcc(R,model)\). The implication does not promote raw truth to theorem standing and does not create a global negative endpoint outcome; it classifies only the live countercase use inside the governed reductio subproof.
\end{proof}

\begin{lemma}[Local separator occupation contains candidate-image exclusion]
\label{lem:local-separator-occupation-contains-candidate-image-exclusion}
On the fixed context-bound carrier,
\[
\begin{aligned}
  &\ClayCtx\wedge\FixedDom(R)\wedge\ModelOfCnf(model,R)\wedge\LocalSepImgOcc(R,model)\\
  &\qquad\Rightarrow \CandImageExcl(model).
\end{aligned}
\]
\end{lemma}

\begin{proof}
By Definition~\ref{def:local-separator-occupation}, local separator occupation is the local countercase use of candidate-image exclusion as the endpoint-defeating branch. It therefore contains the same support-level lower-bound content \(\CandImageExcl(model)\), without implying global endpoint occupation or official negative endpoint resolution.
\end{proof}

\begin{theorem}[Local separator occupation is non-positive status work]
\label{thm:local-sepocc-is-non-positive-status-work}
In the fixed SAT endpoint proof, local separator occupation is non-positive endpoint-status work over the candidate image:
\[
\begin{aligned}
  &\ClayCtx\wedge\FixedDom(R)\wedge\ModelOfCnf(model,R)\wedge\LocalSepImgOcc(R,model)\\
  &\qquad\Rightarrow \NonPositiveStatusWork_R(model).
\end{aligned}
\]
\end{theorem}

\begin{proof}
By Lemma~\ref{lem:local-separator-occupation-contains-candidate-image-exclusion}, local separator occupation contains \(\CandImageExcl(model)\). By Corollary~\ref{cor:candidate-image-exclusion-not-coequal-occupation}, this is non-occupation of the fixed positive candidate image. The local occupation hypothesis says that the content is used as the endpoint-defeating negative branch inside the countercase subproof, so it is not support-only or bookkeeping in that subproof. The carrier and model remain fixed, so it is not carrier shift. By Theorem~\ref{thm:negative-coequal-requires-interface-change}, it is not lawful coequal image occupation. Hence it is non-positive endpoint-status work on the fixed SAT image.
\end{proof}

\begin{theorem}[Local separator occupation induces the separator discriminator]
\label{thm:local-sepocc-induces-sepdisc}
For the fixed SAT endpoint carrier,
\[
\begin{aligned}
  &\ClayCtx\wedge\FixedDom(R)\wedge\ModelOfCnf(model,R)\wedge\LocalSepImgOcc(R,model)\\
  &\qquad\Rightarrow \SepDisc(model).
\end{aligned}
\]
\end{theorem}

\begin{proof}
Theorem~\ref{thm:local-sepocc-is-non-positive-status-work} gives \(\NonPositiveStatusWork_R(model)\). Theorem~\ref{thm:non-positive-status-work-induces-sepdisc} then identifies that non-positive same-domain endpoint-status work with the separator discriminator \(\SepDisc(model)\).
\end{proof}

\begin{theorem}[Local separator occupation triggers no-independent closure]
\label{thm:local-sepocc-triggers-noindependent}
For the fixed SAT endpoint carrier,
\[
\begin{aligned}
  &\ClayCtx\wedge\FixedDom(R)\wedge\ModelOfCnf(model,R)\wedge\LocalSepImgOcc(R,model)\\
  &\qquad\Rightarrow \NoSepDisc(model).
\end{aligned}
\]
\end{theorem}

\begin{proof}
Local separator occupation is a governed local proof act in the same fixed endpoint regime. It preserves the same carrier and model and performs endpoint-facing separator work inside the reductio subproof. By Theorem~\ref{thm:sat-sepdisc-endpoint-adm-relevant}, any \(\SepDisc(model)\) induced by that occupation is endpoint-admissibility-relevant: it changes the branch standing, reportability, downstream reuse, or exclusion profile of the local endpoint-resolution act. Governed local endpoint use is subject to the same kernel-forced A+ consequence layer as global endpoint use: the AMetric boundary, unique admissible interior, standing conservation, no-selector discipline, transport closure, and structural rigidity do not suspend inside a reductio subproof. Therefore the local separator act cannot license an independent same-domain separator discriminator. This is the local form of Theorem~\ref{thm:no-independent-sat-separator-classifier}, and it yields \(\NoSepDisc(model)\).
\end{proof}


\begin{corollary}[A+ consequences bind local separator countercase use]
\label{cor:aplus-binds-local-separator-countercase-use}
Local separator occupation in the exact-complement reductio subproof is not exempt from the \(\Aplus\) no-independent-discriminator discipline:
\[
\begin{aligned}
  &\ClayCtx\wedge\FixedDom(R)\wedge\ModelOfCnf(model,R)\wedge\LocalSepImgOcc(R,model)\\
  &\qquad\Rightarrow \SepDisc(model)\wedge\NoSepDisc(model)\wedge\neg\SepDisc(model).
\end{aligned}
\]
\end{corollary}

\begin{proof}
Theorem~\ref{thm:local-sepocc-induces-sepdisc} gives \(\SepDisc(model)\). Theorem~\ref{thm:local-sepocc-triggers-noindependent} gives \(\NoSepDisc(model)\). Theorem~\ref{thm:no-sepdisc-excludes-sepdisc} gives \(\neg\SepDisc(model)\). The point is the same as in the global endpoint-use case: local countercase use is a governed proof act on the same fixed endpoint carrier, so the kernel-forced \(\Aplus\) consequence layer does not suspend inside the reductio subproof.
\end{proof}

\begin{corollary}[Local separator occupation is impossible]
\label{cor:local-sepocc-impossible}
In the fixed SAT endpoint context,
\[
\begin{aligned}
  &\ClayCtx\wedge\FixedDom(R)\wedge\ModelOfCnf(model,R)\\
  &\qquad\Rightarrow \neg\LocalSepImgOcc(R,model).
\end{aligned}
\]
\end{corollary}

\begin{proof}
Assume the fixed context and suppose \(\LocalSepImgOcc(R,model)\). Theorem~\ref{thm:local-sepocc-induces-sepdisc} gives \(\SepDisc(model)\). Theorem~\ref{thm:local-sepocc-triggers-noindependent} gives \(\NoSepDisc(model)\), and Theorem~\ref{thm:no-sepdisc-excludes-sepdisc} gives \(\neg\SepDisc(model)\). Contradiction. Hence local separator occupation is impossible.
\end{proof}

\begin{corollary}[Negative endpoint counterforce is excluded locally]
\label{cor:negative-endpoint-counterforce-excluded}
In the fixed SAT endpoint context,
\[
\begin{aligned}
  &\ClayCtx\wedge\FixedDom(R)\wedge\ModelOfCnf(model,R)\wedge\OfficialSATEval(R,model)\\
  &\qquad\Rightarrow \neg\EndpointCounterforce_R(\SepBare).
\end{aligned}
\]
\end{corollary}

\begin{proof}
Assume the fixed context and suppose \(\EndpointCounterforce_R(\SepBare)\). Theorem~\ref{thm:local-counterforce-induces-local-separator-occupation} gives \(\LocalSepImgOcc(R,model)\). Corollary~\ref{cor:local-sepocc-impossible} excludes that local separator occupation on the same fixed carrier. Contradiction. Hence the raw negative branch has no endpoint counterforce.
\end{proof}

\begin{theorem}[Raw negative branch can defeat endpoint closure only by endpoint counterforce]
\label{thm:raw-negative-defeats-only-by-counterforce}
In the fixed official SAT endpoint-closure evaluation, the bare negative branch can oppose the positive endpoint only by endpoint counterforce. Equivalently, a raw negative branch that is not asserted with \(\EndpointCounterforce_R(\SepBare)\) may remain a proposition-level branch, support-level lower-bound content, epistemic possibility, or hypothetical case, but it is not the same-mode endpoint counterbranch against the positive endpoint.
\end{theorem}

\begin{proof}
The theorem is definitional at the endpoint-force level and does not assert that raw truth implies theoremhood. The role \(\EndpointCounterforce_R(\SepBare)\) was introduced precisely for the use in which the raw negative branch is invoked to block or defeat endpoint closure on the fixed carrier. If \(\SepBare\) is not used with that force, then it does not oppose the endpoint closure; it remains raw, support-level, epistemic, or hypothetical. If it is used with that force, it is endpoint counterforce by Definition~\ref{def:endpoint-counterforce}. Thus the negative branch can matter against the endpoint only through endpoint counterforce.
\end{proof}

\begin{remark}[Raw cases, endpoint outcomes, and endpoint closure]
\label{rem:raw-cases-versus-endpoint-outcomes}
The raw bivalence theorem \(\PosRaw\vee\SepBare\) remains a proposition-level statement. A raw assumption of \(\SepBare\) inside a local proof by cases is not by itself official endpoint resolution and is not a global negative endpoint outcome. The stronger premise used above is official endpoint-closure evaluation of the fixed target together with local countercase discipline: if the raw negative branch is introduced as the exact endpoint-defeating countercase, it induces local separator occupation. That local separator route is impossible by the same discriminator contradiction. Thus the reductio discharges the local assumption without treating hypothetical syntax as endpoint standing or as official negative resolution.
\end{remark}



\begin{theorem}[Local countercase contradiction excludes the raw negative branch]
\label{thm:local-countercase-contradiction-excludes-raw-negative}
In the fixed SAT endpoint-closure evaluation, the bare negative branch is impossible by local reductio:
\[
\begin{aligned}
  &\ClayCtx\wedge\FixedDom(R)\wedge\ModelOfCnf(model,R)\wedge\OfficialSATEval(R,model)\\
  &\qquad\Rightarrow \neg\SepBare.
\end{aligned}
\]
\end{theorem}

\begin{proof}
Assume the fixed context hypotheses and open a reductio subproof with the sole object-level assumption \([\SepBare]^i\), introduced as the exact raw complement of the target \(\PosRaw\). By Lemma~\ref{lem:exact-complement-reductio-annotation-rule}, this proof act carries the annotation \(\ReductioCountercase_R(\SepBare;\PosRaw)\). The annotation is not a second assumption and not a conjunct added to \(\SepBare\); it is the local role of the discharged assumption \([\SepBare]^i\). By Theorem~\ref{thm:reductio-countercase-yields-local-countercase}, the exact-complement reductio assumption has local endpoint-countercase standing, \(\LocalCountercase_R(\SepBare)\). By Theorem~\ref{thm:local-countercase-has-endpoint-counterforce}, this local countercase has \(\EndpointCounterforce_R(\SepBare)\). Corollary~\ref{cor:negative-endpoint-counterforce-excluded} excludes \(\EndpointCounterforce_R(\SepBare)\) under the same fixed context and official endpoint-closure evaluation, yielding contradiction inside the annotated subproof. By Theorem~\ref{thm:exact-complement-annotation-discharge}, that contradiction discharges the original object-level assumption \([\SepBare]^i\) and gives \(\neg\SepBare\). No step says that raw truth implies theorem standing; only the exact complement assumption as used in the endpoint reductio receives temporary derivational standing, and the annotation is admissible for discharge without adding a premise.
\end{proof}

\begin{theorem}[Endpoint closure is derived by local countercase discharge]
\label{thm:endpoint-closure-derived-not-assumed}
In the fixed SAT endpoint-closure evaluation, local reductio exclusion of the raw negative branch yields the positive raw branch:
\[
\begin{aligned}
  &\ClayCtx\wedge\FixedDom(R)\wedge\ModelOfCnf(model,R)\wedge\OfficialSATEval(R,model)\\
  &\qquad\Rightarrow \PosRaw.
\end{aligned}
\]
\end{theorem}

\begin{proof}
Assume the fixed context hypotheses. Theorem~\ref{thm:local-countercase-contradiction-excludes-raw-negative} gives \(\neg\SepBare\). By raw bivalence, Theorem~\ref{thm:bare-bivalence}, either \(\PosRaw\) or \(\SepBare\). The second branch contradicts \(\neg\SepBare\). Therefore \(\PosRaw\). Equivalently, raw complementarity \(\SepBare\leftrightarrow\neg\PosRaw\) and classical discharge give the same conclusion.
\end{proof}

\begin{theorem}[Positive SAT endpoint]
\label{thm:positive}
In the fixed SAT endpoint evaluation,
\[
\begin{aligned}
  &\ClayCtx\wedge\FixedDom(R)\wedge\ModelOfCnf(model,R)\wedge\OfficialSATEval(R,model)\\
  &\qquad\Rightarrow \PosRaw.
\end{aligned}
\]
\end{theorem}

\begin{proof}
Assume the four hypotheses displayed in the theorem. By Theorem~\ref{thm:candidate-image-carrier-adequacy-raw-split}, the official raw split is represented on the context-bound candidate image: \(\PosRaw\) is image occupation and \(\SepBare\) is candidate-image exclusion on that same image. The proof now excludes the raw negative branch by local countercase reductio rather than by truth-implies-standing or endpoint-complete evaluation. Theorem~\ref{thm:local-countercase-contradiction-excludes-raw-negative} gives \(\neg\SepBare\) under the same fixed context and official endpoint-closure evaluation. Theorem~\ref{thm:endpoint-closure-derived-not-assumed} then applies and yields \(\PosRaw\).
\end{proof}

\begin{theorem}[Instantiated SAT endpoint context]
\label{thm:rsat-context}
The official unrestricted finite CNF-SAT endpoint carrier instantiates the fixed context, fixed-domain, bound-model, and official endpoint-closure evaluation hypotheses:
\[
\begin{aligned}
  &\ClayCtx_{\RSATctx}\wedge\FixedDom(\RSATctx)\wedge\ModelOfCnf(m_{\mathsf{SAT}},\RSATctx)\\
  &\qquad\wedge\OfficialSATEval(\RSATctx,m_{\mathsf{SAT}}).
\end{aligned}
\]
\end{theorem}

\begin{proof}
The official SAT endpoint carrier is arbitrary finite CNF over finite Boolean alphabets with formula length as input size; the positive branch is deterministic polynomial-time CNF-SAT decidability; the raw negative branch is its complement on the same carrier; and admissible encodings preserve the carrier and endpoint image. This is exactly the fixed context of Definition~\ref{def:context} together with fixed-domain preservation and the context-bound CNF model of Definition~\ref{def:cnf-model-of-context}. The manuscript's proof class is endpoint-closure evaluation of the official SAT target rather than a report of unresolved pre-proof status. Thus the official-evaluation condition supplied here is target-fixing and endpoint-counterforce discipline for the endpoint-closure proof; it is not endpoint-complete evaluation and does not assert either branch.
\end{proof}

\begin{theorem}[Instantiated positive SAT endpoint]
\label{thm:positive-instantiated}
\[
  \CnfSATP.
\]
\end{theorem}

\begin{proof}
By Theorem~\ref{thm:rsat-context}, the official SAT carrier supplies the fixed context, fixed-domain, bound-model, and official endpoint-closure evaluation hypotheses. Applying Theorem~\ref{thm:positive} with \(R=\RSATctx\) and \(model=m_{\mathsf{SAT}}\) yields \(\PosRaw\). By Definition~\ref{def:bare}, \(\PosRaw:=\CnfSATP\).
\end{proof}

\section{Positive Report Classification}\label{sec:positive-classification}

\begin{definition}[Positive Clay report act]\label{def:positive-report}
\(\ReportClay{\PosRaw}\) is the report act that presents \(\PosRaw\) as the positive mathematical endpoint after the separator branch has been excluded.
\end{definition}

\begin{theorem}[Positive report act]\label{thm:positive-report}
\[
  \ClayCtx\wedge\FixedDom(R)\wedge\ModelOfCnf(model,R)\wedge\PosRaw\Rightarrow\ReportClay{\PosRaw}.
\]
\end{theorem}

\begin{proof}
Once the fixed SAT context is in force and the raw positive endpoint has been derived after exclusion of the separator branch, the manuscript's theorem-bearing output is the report of \(\PosRaw\) as the positive mathematical endpoint. That report act does not supply standing; it records the already derived endpoint under the report-preserving interface forced by the kernel.
\end{proof}

\begin{theorem}[Positive endpoint classification]\label{thm:pos-class}
\[
  \ClayCtx\wedge\FixedDom(R)\wedge\ModelOfCnf(model,R)\wedge\PosRaw\wedge\ReportClay{\PosRaw}
  \Rightarrow \Apos.
\]
\end{theorem}

\begin{proof}
The positive branch is a theorem-bearing endpoint report on the fixed SAT carrier. By Theorem~\ref{thm:positive-direct-occupation}, the positive branch is direct constructive endpoint occupation rather than separator classification. By Theorem~\ref{thm:kernel-necessity}, endpoint use is kernel-instantiated; by Theorem~\ref{thm:a-plus}, the report-preserving A+ consequence layer applies. The positive classification \(\Apos\) records that downstream status. It is not a source of standing and is not an independent separator discriminator.
\end{proof}

\begin{corollary}[Positive Clay mathematical endpoint]\label{cor:pos-clay}
\[
  \ClayCtx\wedge\FixedDom(R)\wedge\ModelOfCnf(model,R)\Rightarrow\PosClay.
\]
\end{corollary}

\begin{proof}
Theorem~\ref{thm:positive} gives \(\PosRaw\). Theorem~\ref{thm:positive-report} gives \(\ReportClay{\PosRaw}\). Theorem~\ref{thm:pos-class} gives \(\Apos\). Therefore \(\PosClay:=\PosRaw\wedge\Apos\).
\end{proof}

\section{AASC Closure and Endpoint Correspondence}

The manuscript does not redefine \(P\), \(NP\), CNF, satisfiability, polynomial time, reduction, or NP-completeness. AASC proves the positive raw SAT endpoint by excluding the separator endpoint image through kernel-forced fixed-domain discriminator closure. Once \(\CnfSATP\) is obtained, the final move to \(P=NP\) is the standard Cook--Levin/Karp endpoint correspondence. Cook--Levin/Karp is not the branch-exclusion method; it identifies the official endpoint consequence after the AASC proof has delivered unrestricted finite CNF-SAT in deterministic polynomial time.

\begin{theorem}[Cook--Levin/Karp endpoint correspondence]\label{cor:pnp}
\[
  \CnfSATP\Rightarrow \Pclass=\NPclass.
\]
\end{theorem}

\begin{proof}
Arbitrary finite CNF-SAT is NP-complete under the standard Cook--Levin/Karp polynomial-time many-one reduction framework. If CNF-SAT has a deterministic polynomial-time decider, then every language in \(\NPclass\) is decided in deterministic polynomial time by composing its polynomial-time reduction to CNF-SAT with that decider. Hence \(\NPclass\subseteq\Pclass\). Since \(\Pclass\subseteq\NPclass\) by definition, \(\Pclass=\NPclass\).
\end{proof}

\begin{corollary}[Positive P versus NP endpoint]\label{cor:pnp-final}
\[
  \Pclass=\NPclass.
\]
\end{corollary}

\begin{proof}
Theorem~\ref{thm:positive-instantiated} gives \(\CnfSATP\). Theorem~\ref{cor:pnp} gives \(\Pclass=\NPclass\).
\end{proof}


\section{Denial Burden and Exhaustion of Remaining Routes}\label{sec:denial-burden}

\subsection*{Dependency-inversion objections}
An objection that treats the kernel as an optional AASC overlay is not a separate denial route. It is a dependency-inversion objection. Such an objection has standing only if it identifies a failure in the kernel-neutral adequacy argument: target determinacy, step evaluability, act-time finality, or same-regime fidelity must be shown not to be required for theorem-bearing endpoint use, or a same-regime faithful counterexample must be supplied that preserves all four while omitting at least one kernel role. Without that showing, the objection merely renames a forced precondition as optional vocabulary.

\subsection*{Cost of abandoning the kernel}
A denial of the kernel is not a harmless refusal of terminology. The kernel is the role-decomposition of the four adequacy conditions required for determinate same-domain theorem-bearing construction. To deny the kernel while still claiming to assess the same endpoint, an objector must identify which adequacy condition is being abandoned and accept the corresponding loss.

\begin{center}
\small
\begin{tabularx}{\linewidth}{>{\raggedright\arraybackslash}p{0.27\linewidth}>{\raggedright\arraybackslash}p{0.18\linewidth}X}
\toprule
Denied adequacy condition & Kernel role lost & Consequence of denial \\
\midrule
Target determinacy & Reference & No fixed endpoint object, carrier, candidate image, or branch slot is preserved. \\
Step evaluability & Standing & No transition has theorem-bearing stephood; lower-bound traces are inscriptions rather than endpoint acts. \\
Act-time finality & Irreversibility & Same-act repair, later reclassification, and retrospective status change remain open. \\
Same-regime fidelity & Admissibility boundary & Carrier, domain, encoding image, or endpoint classification may shift under redescription. \\
\bottomrule
\end{tabularx}
\end{center}

Thus kernel denial has only two possible forms. Either it denies none of the four adequacy conditions, in which case the kernel has already been conceded under different vocabulary; or it denies at least one of them, in which case the objector is no longer preserving determinate same-domain theorem-bearing endpoint use. There is no third position from which one may reject the kernel while retaining the same official endpoint-resolution act.


\subsection*{Remaining same-mode denial fork}
After the official CNF-SAT resolution-route, carrier-projection, ordinary-resolution, local-countercase, trichotomy, endpoint-admissibility, and A+ binding results are in force, the ordinary-practice objection has no remaining middle role or exemption. A same-domain incompleteness objection also has no independent endpoint force: if it changes endpoint standing, reportability, or downstream reuse, it is endpoint-admissibility-relevant gate work; if it does not, it is endpoint-inert. A critic who accepts the standard CNF-SAT projection, the fixed candidate-image carrier, and the trichotomy classification must now deny one of the named links below.
\begin{center}
\small
\begin{tabularx}{\linewidth}{>{\raggedright\arraybackslash}p{0.24\linewidth}X}
\toprule
Remaining route & What must be shown \\
\midrule
Projection failure & The symmetric \(P=NP/P\neq NP\) outcome carrier does not faithfully project to the SAT endpoint-image carrier; in particular, \(P=NP\) is not \(\PosRaw\) / image occupation, or \(P\neq NP\) is not \(\SepBare\) / candidate-image exclusion. \\
Trichotomy failure & There exists a theorem-bearing, endpoint-resolving, same-domain, non-positive, non-coequal, non-support, non-bookkeeping, non-repair, non-carrier-shifting role that is nevertheless not independent endpoint-status work. \\
\(\Aplus\) failure & Endpoint use on a non-degenerate fixed SAT carrier does not bind the identified independent discriminator to the kernel-forced \(\Aplus\) no-independent-discriminator consequence. \\
Same-domain incompleteness rescue & A G\"odel-style, true-but-unprovable, semantic, infinitary, or extension-based distinction changes endpoint standing, reportability, or downstream reuse without becoming inert bookkeeping, extensional gate renaming, explicit scope change, same-side refinement, independent authorizer, or inadmissible authority import. \\
\bottomrule
\end{tabularx}
\end{center}
Appeals to ordinary mathematical practice do not by themselves meet any of these burdens. Ordinary theoremhood is granted; the remaining question is the endpoint role of that theorem-bearing act after the accepted SAT projection and official endpoint use.

\subsection*{Right-mode objection targets}
A standing mathematical objection must attack one of the following load-bearing claims.
\begin{enumerate}
\item AASC mathematical status: identify a precise failure of a definition, theorem, inference, kernel dependency, or formal constraint.
\item Lean support boundary: identify a mismatch between a Lean-facing support claim and the manuscript theorem chain, if Lean support is invoked.
\item Kernel-neutral target adequacy: show that target determinacy, step evaluability, act-time finality, or same-regime fidelity is not required for theorem-bearing endpoint use.
\item Non-degenerate SAT endpoint regime forcing the kernel: preserve fixed target, carrier, theorem-status distinction, downstream reuse, redescription stability, and act-time finality while lacking one of \(\Adm,\St,\RefK,\Irr\).
\item Local SAT packet: show that a used A+ component is not forced by the kernel plus fixed-domain closure.
\item Strict weakening of K5, K6, K11, or K13: define a weaker regime preserving the same endpoint roles while allowing separator governance.
\item Ordinary official endpoint resolution bridge: produce an ordinary mathematical act accepted as the official resolution of a fixed endpoint while fixing target, carrier, branch, theorem-status, downstream reuse, act-time finality, and redescription-stable endpoint role, but not satisfying \(\EndpointUse_R\).
\item Endpoint-counterforce discipline: show that a raw branch can defeat endpoint closure on the fixed carrier while lacking endpoint counterforce, or that endpoint counterforce avoids the local separator route.
\item Epistemic-status separation: show that open, unresolved, under-review, or proof-unavailable status is a third endpoint outcome of the official SAT target rather than an epistemic status outside endpoint outcome.
\item Raw-case promotion discipline: show that the manuscript infers official endpoint resolution from a merely hypothetical assumption of \(\SepBare\), rather than routing the exact-complement reductio assumption through \(\ReductioCountercase_R(\SepBare;\PosRaw)\), \(\LocalCountercase_R(\SepBare)\), and \(\LocalSepImgOcc(R,model)\).
\item Exact-complement reductio bridge: produce a reductio proof of the official positive endpoint in which the live exact complement assumption blocks the target while failing to count as the local endpoint countercase.
\item Reductio annotation discipline: show that \(\ReductioCountercase_R(\SepBare;\PosRaw)\) functions as an additional object-level premise rather than as the proof-act annotation of the sole discharged assumption \([\SepBare]^i\).
\item Annotation-discharge discipline: show that contradiction derived in the annotated exact-complement subproof discharges only a strengthened assumption rather than the original object-level assumption \([\SepBare]^i\).
\item Ordinary theoremhood shield: show that ordinary theoremhood used as official endpoint resolution is exempt from endpoint-role discipline while retaining official endpoint force.
\item Projection-role content: show that ordinary practice can use the standard CNF-SAT correspondence as the route of official resolution while refusing the endpoint-role content assigned by that projection, namely positive candidate-image occupation for \(P=NP\) and candidate-image non-occupation for \(P\neq NP\).
\item No one-step retyping: show that the manuscript derives \(\SepDisc(model)\) directly from raw \(\SepBare\), \(\neg\PosRaw\), or lower-bound normal form, bypassing candidate-image exclusion, endpoint-use promotion, and separator occupation.
\item Native lower-bound falsification package: produce a same-carrier deployed \(\LBNatSAT(R,model)\) that defeats the official endpoint while not being endpoint use.
\item Native lower-bound falsification collapse: produce a same-carrier endpoint-defeating lower-bound role outside raw/support theoremhood, positive bridge-neutralization, bookkeeping, carrier shift, and \(\SepDisc(model)\).
\item Bridge-neutralized lower-bound package: show a live negative branch compatible with an encoded non-failing polynomial-time candidate on the same context-bound candidate image.
\item Ordinary lower-bound theoremhood: show that ordinary theoremhood used as official endpoint resolution is exempt from the endpoint-use features ordinary resolution itself fixes.
\item Ordinary theorem-bearing non-occupation: show that ordinary theorem-bearing non-occupation, after accepted SAT projection and official endpoint use, supplies a fifth endpoint role that is not support, bookkeeping, carrier shift, positive occupation, coequal image occupation, positive bridge-neutralization, or independent endpoint-status work.
\item Remaining denial fork: show that an objection to the separator contradiction avoids all three remaining targets--projection failure, trichotomy failure, and \(\Aplus\) failure.
\item \(\Aplus\) binding for ordinary official resolution: show an ordinary official endpoint-resolution act that is \(\EndpointUse_R\) and is classified as independent by the trichotomy while nevertheless being exempt from the \(\Aplus\) no-independent-discriminator consequence.
\item Standard encoded CNF model adequacy: identify a concrete failure of coverage, soundness, failure adequacy, or candidate-image exactness.
\item Raw complementarity: show \(\SepBare\not\leftrightarrow\neg\PosRaw\) on the unrestricted finite CNF-SAT carrier.
\item SAT-native bridge: show that \(\SepBare\) does not yield standard lower-bound normal form, or that standard lower-bound normal form is not candidate endpoint-image exclusion.
\item Theorem-level discriminator: show that candidate endpoint-image exclusion used at endpoint level does not classify endpoint status.
\item Ordinary negation role-classification: show that the proof rejects ordinary negative theoremhood merely for being negative.
\item Single-interface theorem: show an endpoint-image occupant of the fixed CNF-SAT carrier that is not an encoded non-failing polynomial-time candidate, or show a same-carrier lawful coequal negative occupation interface inside the positive candidate image.
\item Single-interface non-explosiveness: show that the proof derives contradiction from \(\CandImageExcl(model)\) alone, rather than from \(\SepImgOcc\Rightarrow\SepDisc(model)\) together with \(\EndpointUse_R(\SepImgOcc)\Rightarrow\NoSepDisc(model)\).
\item Coequal-role denial: show that \(\SepBare\) is a lawful coequal endpoint role inside the same positive endpoint-image carrier rather than endpoint-status governance, without moving to the symmetric outcome carrier \(\OutcomePvNP\).
\item Status-work typing: show same-domain non-positive endpoint-status work on the fixed candidate image that is neither positive occupation, support, bookkeeping, carrier shift, repair, lawful coequal image occupation, nor independent separator-discriminator work.
\item No hidden fifth case: inhabit endpoint-resolving, same-carrier, non-governing candidate-image exclusion.
\item Endpoint governance: show that official negative endpoint use can remain theorem-bearing while avoiding the occupation-exhaustion and endpoint-status-bivalence locks.
\item Independent discriminator: show that the image-separator branch can avoid inducing independent same-domain endpoint-status discrimination under the SAT-local independence bridge.
\item No independent discriminator: show that such an independent same-domain endpoint-status discriminator survives AMetricity, unique interior, standing conservation, transport closure, and structural rigidity.
\item Non-explosiveness: show that the proof globally bans negative complexity mathematics rather than excluding only same-carrier endpoint-resolving separator governance.
\item Route-locus exhaustion: show a same-carrier endpoint-resolving negative SAT route outside L1--L13, including the ordinary native lower-bound falsification locus.
\item Cook--Levin/Karp endpoint correspondence: show that \(\CnfSATP\) does not imply \(P=NP\) under the standard NP-completeness correspondence.
\end{enumerate}

\subsection*{Faithful same-regime counterexample worksheet}
A successful same-mode objection must complete at least one of the following tasks.
{\footnotesize
\begin{longtable}{@{}>{\raggedright\arraybackslash}p{0.30\linewidth}>{\raggedright\arraybackslash}p{0.66\linewidth}@{}}
\toprule
Target & Successful objection must produce \\
\midrule
Kernel denial & A theorem-bearing SAT endpoint act preserving fixed carrier, candidate image, branch slots, target-slot fixation, report evaluability, same-regime fidelity, and act-time finality while abandoning a named kernel role. \\
Strict weakening & A weaker regime \(\WeakSep\) preserving the unrestricted finite CNF carrier, encoded candidate image, target-slot fixation, minimal report evaluability, standing-admissibility discipline, act-time finality, and no carrier shift while permitting endpoint-resolving separator governance that is not positive occupation, proof support, bookkeeping, carrier shift, repair, lawful coequal target role, or second endpoint authority. \\
Model adequacy & A deterministic polynomial-time SAT candidate lost by the encoded image, an unlawful candidate admitted as polynomial-time, an encoding artifact that changes SAT-failure status, or a mismatch between \(\CnfSATP\) and the encoded non-failing candidate image. \\
Raw negative route & A raw negative SAT branch on the unrestricted finite CNF carrier that does not reduce to standard lower-bound normal form over the fixed candidate image. \\
Candidate-image bridge & A standard lower-bound normal form that is not candidate endpoint-image exclusion. \\
Ordinary negation & A proof that the manuscript treats negative theoremhood as defective merely because it is negative. \\
Ordinary official resolution & An ordinary endpoint-resolution act that fixes target, carrier, branch, theorem-status, downstream reuse, act-time finality, and redescription-stable endpoint role while nevertheless not being \(\EndpointUse_R\). \\
Endpoint counterforce without local separator route & A raw branch that defeats endpoint closure on the fixed carrier while avoiding local separator occupation. \\ 
Raw-case promotion & A proof that the manuscript treats a merely hypothetical assumption of \(\SepBare\) as official negative endpoint resolution, rather than using the local countercase-to-local-separator route. \\
Ordinary theoremhood shield & A proof that ordinary theoremhood used as official endpoint resolution remains exempt from endpoint-role discipline while retaining official endpoint force. \\
No one-step retyping & A derivation from raw \(\SepBare\), \(\neg\PosRaw\), or lower-bound normal form directly to \(\SepDisc(model)\) that bypasses candidate-image exclusion, endpoint-use promotion, and separator occupation. \\
Native lower-bound falsification package & A same-carrier deployed \(\LBNatSAT(R,model)\) that defeats the official endpoint while not being endpoint use. \\
Native lower-bound falsification collapse & A same-carrier endpoint-defeating lower-bound role outside raw/support theoremhood, positive bridge-neutralization, bookkeeping, carrier shift, and \(\SepDisc(model)\). \\
Bridge-neutralized lower-bound package & A live negative branch compatible with an encoded non-failing polynomial-time candidate on the same context-bound candidate image. \\
Ordinary theorem-bearing non-occupation & A theorem-bearing, endpoint-resolving, same-domain, non-positive, non-coequal, non-support, non-bookkeeping, non-repair, non-carrier-shifting role that is not independent same-domain endpoint-status work and is not covered by the native lower-bound falsification collapse. \\
Remaining denial fork & A same-mode objection that accepts the SAT projection, the endpoint-discriminator trichotomy, and the \(\Aplus\) no-independent-discriminator consequence while still avoiding the separator contradiction. \\
Single occupation interface & An occupant of the fixed CNF-SAT endpoint image that is not an encoded non-failing polynomial-time candidate, or a proof that universal candidate failure occupies that same image as a coequal occupant. \\
Single-interface non-explosiveness & A proof that the manuscript infers contradiction from single-interface typing or \(\CandImageExcl(model)\) alone, rather than from endpoint-resolving separator occupation plus the no-independent-discriminator contradiction. \\
Coequal endpoint & A proof that \(\SepBare\) is a lawful coequal endpoint role inside the same positive endpoint-image carrier, without changing to \(\OutcomePvNP\). \\
Endpoint governance & Endpoint-resolving candidate-image exclusion that is not proof support, not bookkeeping, not carrier shift, not positive occupation, not lawful coequal image occupation, and not endpoint-status governance. \\
Independent discriminator & Image-separator endpoint use that does not induce \(\SepDisc(model)\), or an induced \(\SepDisc(model)\) compatible with \(\NoSepDisc(model)\). \\
Final endpoint & A failure of \(\neg(\SepBare\wedge\OfficialEndpointResolution_R(\neg\PosRaw))\) in the official endpoint evaluation, or of \((\SepBare\leftrightarrow\neg\PosRaw)\Rightarrow\PosRaw\), or a failure of the Cook--Levin/Karp correspondence from \(\CnfSATP\) to \(P=NP\). \\
\bottomrule
\end{longtable}
}

Objections outside these classes are non-standing unless they identify a theorem-level failure in one of the cited bridges. The eliminated negative route is the raw SAT lower-bound/no-positive branch, because that branch transports directly to candidate endpoint-image exclusion and, under endpoint-resolution use, separator endpoint-image occupation; endpoint-resolving separator occupation then induces independent same-domain endpoint-status discrimination under the SAT-local bridge; and the kernel-forced consequence layer excludes such discrimination.


\section{Conclusion}
The branch-exclusion argument is AASC-internal and kernel-first. The kernel is not an optional overlay; it is forced by the conditions under which any proof, counterproof, endpoint resolution, or theorem-bearing act remains a determinate same-carrier mathematical act. Raw traces do not become constructions merely by being written. Genuine theoremhood is governed derivation and is therefore kernel-internal. Negative theoremhood is not rejected for being negative.

On the fixed unrestricted finite CNF-SAT carrier, the positive branch is direct constructive endpoint occupation. The raw negative branch is ordinary logical negation at the bare level, but endpoint-bearing use of that negation is role-classified. If it remains proof support or ordinary proposition-targeted theoremhood, it does not resolve the official endpoint. If it changes carrier, it is not the same endpoint. If it is offered as the official negative endpoint resolution, it becomes endpoint use of the negative branch. The SAT-native lower-bound/no-positive bridge then transports it to standard lower-bound normal form, candidate endpoint-image exclusion, and the image-separator branch.

The fixed SAT endpoint image has a single occupation interface: an encoded non-failing polynomial-time CNF-SAT candidate. The candidate-image carrier is adequate for the official raw split: the positive branch is occupation of that image, and the bare negative branch is non-occupation of that same exact image. Universal candidate failure is therefore non-occupation of that image, not a lawful coequal occupant of it. This single-interface fact is not itself an endpoint-closure theorem and does not ban lower-bound or non-occupation statements. Treating the negative outcome as coequal belongs to the symmetric two-output problem carrier, not to the fixed positive candidate-image carrier used in the proof. When candidate-image non-occupation is used to resolve the endpoint on the same carrier, it performs non-positive endpoint-status work. Ordinary theoremhood supplies proof legitimacy, not an exemption from endpoint-role classification. Once the ordinary negative theorem is used as official endpoint resolution after the SAT projection, it is not support, not bookkeeping, not carrier shift, not positive occupation, and not coequal image occupation; the endpoint-discriminator trichotomy therefore leaves only the independent separator-discriminator role. The kernel-forced consequence layer excludes such discriminators. Therefore separator endpoint-image occupation is unavailable, the bare negative branch is excluded in the live exact-complement reductio route, and raw complementarity yields \(\PosRaw\). Since \(\PosRaw\) is \(\CnfSATP\), arbitrary finite CNF-SAT is decidable in deterministic polynomial time. The standard Cook--Levin/Karp endpoint correspondence then gives \(\Pclass=\NPclass\). Ordinary mathematical practice is not asked to adopt AASC vocabulary; it is only asked not to split official endpoint force from the fixed-domain consequences of endpoint use. Once ordinary official resolution supplies endpoint use and the trichotomy identifies an independent discriminator, Theorems~\ref{thm:endpoint-status-work-adm-relevant} and~\ref{thm:sat-sepdisc-endpoint-adm-relevant} show that the discriminator is endpoint-admissibility-relevant: it affects endpoint standing, reportability, downstream reuse, and branch exclusion for the resolution act. Corollary~\ref{cor:aplus-binds-ordinary-official-resolution} and Corollary~\ref{cor:aplus-excludes-independent-endpoint-admrel} then bind that act to the \(\Aplus\) no-independent-discriminator consequence. The remaining same-mode denial burden is exhausted by failure of the official CNF-SAT projection, failure of the endpoint-discriminator trichotomy, failure of the kernel-to-\(\Aplus\) no-independent-discriminator consequence, or a same-domain incompleteness architecture that changes endpoint standing while evading the route-exhaustion results cited in Subsection~\ref{subsec:godel-stability-endpoint}.



\appendix

\section{Lean4 SAT Operator Bridge Appendix}\label{app:lean-sat-operator}

The manuscript proof above remains the primary proof. The Lean4 material in this appendix is support for the local SAT-operator bridge used by the AASC endpoint proof spine:
\[
\begin{aligned}
  \SepBare &\Rightarrow \CandImageExcl(model),\\
  \CandImageExcl(model)\wedge\OfficialEndpointResolution_R(\neg\PosRaw)
    &\Rightarrow \SepImgOcc \Rightarrow \SepDisc(model) \Rightarrow \bot .
\end{aligned}
\]
It is not used as a substitute for the external Cook--Levin/Karp endpoint correspondence from \(\CnfSATP\) to \(P=NP\). Its role is reproducibility support for the AASC foundation-to-SAT endpoint proof spine: the release archive includes the Minimal Conditions/A+ and non-degenerate-kernel audit layers together with the SAT endpoint-operator bridge and negative-branch closure formalized in the release archive \cite{LeanAuditZenodo,LeanAuditGitHub}.

\subsection{Release, environment, and audit runner}

The Lean audit archive included with this project is the following source-material zip:
\begin{lstlisting}[basicstyle=\ttfamily\scriptsize]
source_material/AASC-PvsNP-SAT-Endpoint-Lean-Audit-1.0.0.zip
\end{lstlisting}
The pinned build data recorded in the release are:
\begin{center}
\begin{tabularx}{0.94\linewidth}{>{\raggedright\arraybackslash}p{0.28\linewidth}X}
\toprule
Item & Value \\
\midrule
Lean toolchain & \LeanName{leanprover/lean4:v4.28.0} \\
Lake project & \LeanName{AASCPvsNPSATEndpointLeanAudit} \\
mathlib requirement & \LeanName{mathlib}, revision tag \LeanName{v4.28.0} \\
Pinned mathlib revision & \LeanName{8f9d9cff6bd728b17a24e163c9402775d9e6a365} \\
Main proof queue & \LeanName{MaleyLean/Papers/PvsNP/SATOperatorProofQueue.lean} \\
Focused audit runner & \LeanName{scripts/check-pvsnp-sat-operator-bridge-audit.ps1} \\
Release root suffix & \LeanName{062a2a6} \\
\bottomrule
\end{tabularx}
\end{center}

The verification command supplied by the release is:
\begin{lstlisting}[basicstyle=\ttfamily\scriptsize]
powershell -ExecutionPolicy Bypass -File scripts/check-pvsnp-sat-operator-bridge-audit.ps1
\end{lstlisting}
The runner prints the Lean toolchain, prints the pinned mathlib manifest revision, builds \LeanName{MaleyLean.Papers.PvsNP.SATOperatorAuditRunners}, evaluates the SAT-operator formalization-status and progress summaries, scans the audited AASC foundation plus P vs NP SAT endpoint surface for live \LeanName{axiom}, \LeanName{sorry}, \LeanName{admit}, or \LeanName{unsafe} declarations, builds the AASC foundation/kernel support modules, builds the SAT proof queue, antichecker/model-adequacy bridge, and bridge-callability modules, and runs ten audit files: the Minimal Conditions/A+ audit, the non-degenerate-kernel audit, the SAT endpoint checks, and the combined full-stack AASC/SAT check. The release notes record complete closure for the AASC SAT endpoint and CNF-SAT-in-polynomial-time status summaries. Focused axiom output for the bridge anchors is reported against only the standard Lean/classical/mathlib base dependencies \LeanName{propext}, \LeanName{Classical.choice}, and \LeanName{Quot.sound}.

\subsection{Declaration locations}

The following declaration locations are taken from the released Lean files. They are included to make the manuscript-to-Lean correspondence checkable without treating Lean as a replacement for the manuscript proof.
\begin{lstlisting}[basicstyle=\ttfamily\scriptsize]
SameDomainEndpoint.lean:22
  abbrev CnfPositiveEndpoint : Prop := CnfSATInPolyTime

FoundationalClassifierBridge.lean:44
  def CnfSeparatingClassifierIsIndependentSameDomain

FoundationalClassifierBridge.lean:70
  def CnfNoIndependentSeparatingClassifier

SATOperatorProofQueue.lean:10031
  def CnfSATImageSeparatorBranch

SATOperatorProofQueue.lean:10042
  def CnfSATBareNegativeBranch

SATOperatorProofQueue.lean:10049
  def CnfSATBareSeparator

SATOperatorProofQueue.lean:10142
  def CnfSATImageSeparatorEndpointUse

SATOperatorProofQueue.lean:10161
  def CnfSATDirectPositiveEndpointOccupation

SATOperatorProofQueue.lean:10169
  def CnfSATIndependentSeparatorEndpointStatusDiscriminator

SATOperatorProofQueue.lean:10184
  def CnfSATBookkeepingNegativeOnly
  def CnfSATCarrierShift
  def CnfSATOfficialNegativeEndpointUse
  def CnfSATNegativeOccupationExhaustion

SATOperatorProofQueue.lean:10313
  inductive CnfSATEndpointStatus
  def CnfSATEndpointStatusOccupation
  def CnfSATGovernedEndpointUse
  def CnfSATGovernedEndpointStatusUse
\end{lstlisting}

The proof-spine theorem locations are:
\begin{lstlisting}[basicstyle=\ttfamily\scriptsize]
SATOperatorProofQueue.lean:10075
  cnfBareNegativeBranch_forces_imageSeparatorBranch

SATOperatorProofQueue.lean:10184
  cnfSATNegativeOccupation_exhaustion
  cnfSATNoCarrierShift_of_fixedEndpointDomain
  cnfSATOfficialNegativeEndpointUse_not_bookkeepingOnly
  cnfSATNoOrdinaryNegativeAlternative_of_ametricBoundary
  cnfSATNegativeOccupation_nonoptional

SATOperatorProofQueue.lean:10313
  cnfSATEndpointStatus_exhaustive
  cnfSATGovernedEndpointUse_bivalent
  cnfSATNegativeGovernedEndpointUse_has_separatorStatus
  cnfSATGovernedEndpointStatusUse_eq_separator_of_not_positive
  cnfSATImageSeparatorBranch_of_negativeGovernedEndpointUse

SATOperatorProofQueue.lean:10087
  cnfImageSeparatorBranch_of_not_positiveEndpoint

SATOperatorProofQueue.lean:10104
  cnfSATBareSeparator_forces_imageSeparatorBranch

SATOperatorProofQueue.lean:10148
  cnfSATImageSeparatorBranch_endpointUse

SATOperatorProofQueue.lean:10179
  cnfPositiveEndpoint_directEndpointOccupation

SATOperatorProofQueue.lean:10190
  cnfSATImageSeparatorBranch_requires_independentSeparatorDiscriminator

SATOperatorProofQueue.lean:10212
  cnfBareNegativeBranch_requires_independentSeparatorDiscriminator

SATOperatorProofQueue.lean:10246
  cnfSATImageSeparatorBranch_impossible_of_noIndependentDiscriminator

SATOperatorProofQueue.lean:10267
  cnfBareNegativeBranch_impossible_of_noIndependentDiscriminator

SATOperatorProofQueue.lean:10286
  cnfNotPositiveEndpoint_impossible_of_noIndependentDiscriminator

SATOperatorProofQueue.lean:10308
  cnfPositiveEndpoint_of_noIndependentDiscriminator

SATOperatorProofQueue.lean:10331
  cnfSATInPolyTime_of_noIndependentDiscriminator

SATOperatorProofQueue.lean:10351
  cnfSATInPolyTime_of_context_noIndependentDiscriminator

SATOperatorProofQueue.lean:10372
  cnfSATBareSeparator_impossible_of_context_noIndependentDiscriminator
\end{lstlisting}

\subsection{Manuscript-to-Lean correspondence}

The manuscript-to-Lean symbol correspondence is:
\begin{lstlisting}[basicstyle=\ttfamily\scriptsize]
PosRaw          <-> CnfPositiveEndpoint  (abbrev. CnfSATInPolyTime)
SepBare         <-> CnfSATBareNegativeBranch R model
SepImg          <-> CnfSATImageSeparatorBranch R model (Lean-side image branch)
SepImgOcc       <-> endpoint-occupation use of CnfSATImageSeparatorBranch R model
DirPos          <-> CnfSATDirectPositiveEndpointOccupation
SepDisc(model)  <-> CnfSATIndependentSeparatorEndpointStatusDiscriminator model
NoSepDisc(model)<-> CnfNoIndependentSeparatingClassifier model
SepIndBridge(model)
                <-> CnfSeparatingClassifierIsIndependentSameDomain model
OfficialNegUse  <-> CnfSATOfficialNegativeEndpointUse R model
NegOccExhaust   <-> CnfSATNegativeOccupationExhaustion R model
GovernedUse     <-> CnfSATGovernedEndpointUse R model
EndpointStatus  <-> CnfSATEndpointStatus
\end{lstlisting}

The formal correspondence of the central manuscript chain is:
\begin{lstlisting}[basicstyle=\ttfamily\scriptsize]
CnfPositiveEndpoint
  -> CnfSATDirectPositiveEndpointOccupation

CnfSATBareNegativeBranch R model
  -> CnfSATImageSeparatorBranch R model

CnfSATOfficialNegativeEndpointUse R model
  + CnfSATFixedEndpointDomain R
  + CnfSATContextBoundCNFModel R model
  + AMetricBoundary R
  -> CnfSATImageSeparatorBranch R model

CnfSATGovernedEndpointUse R model
  + Not CnfPositiveEndpoint
  -> CnfSATImageSeparatorBranch R model

CnfSeparatingClassifierIsIndependentSameDomain model
  -> CnfSATImageSeparatorBranch R model
  -> CnfSATIndependentSeparatorEndpointStatusDiscriminator model

CnfNoIndependentSeparatingClassifier model
  -> CnfSeparatingClassifierIsIndependentSameDomain model
  -> Not (CnfSATImageSeparatorBranch R model)

CnfNoIndependentSeparatingClassifier model
  -> CnfSeparatingClassifierIsIndependentSameDomain model
  -> Not (CnfSATBareNegativeBranch R model)

CnfNoIndependentSeparatingClassifier model
  -> CnfSeparatingClassifierIsIndependentSameDomain model
  -> CnfSATInPolyTime
\end{lstlisting}

\subsection{Selected source excerpts}

The following excerpts are reproduced from \LeanName{SATOperatorProofQueue.lean}. They are the Lean-side anchors for the raw-negative bridge, endpoint-use marker, branch asymmetry, and final SAT endpoint closeout.

\begin{lstlisting}[basicstyle=\ttfamily\scriptsize]
def CnfSATImageSeparatorBranch
    {Act Object : Type}
    (_R : MinimalConditionsForAdmissibleConstruction.ConstructionRegime Act Object)
    (_model : CnfEncodedCandidateModel) : Prop :=
  cnfDirectGateLowerBoundResidualTarget

def CnfSATBareNegativeBranch
    {Act Object : Type}
    (_R : MinimalConditionsForAdmissibleConstruction.ConstructionRegime Act Object)
    (_model : CnfEncodedCandidateModel) : Prop :=
  Not CnfPositiveEndpoint

theorem cnfBareNegativeBranch_forces_imageSeparatorBranch
    {Act Object : Type}
    {R : MinimalConditionsForAdmissibleConstruction.ConstructionRegime Act Object}
    {model : CnfEncodedCandidateModel}
    (hBare : CnfSATBareNegativeBranch R model) :
    CnfSATImageSeparatorBranch R model :=
  (cnfDirectGateLowerBoundResidualTarget_iff_noSatInPolyTime).2 hBare

theorem cnfImageSeparatorBranch_of_not_positiveEndpoint
    {Act Object : Type}
    {R : MinimalConditionsForAdmissibleConstruction.ConstructionRegime Act Object}
    {model : CnfEncodedCandidateModel}
    (hNotPositive : Not CnfPositiveEndpoint) :
    CnfSATImageSeparatorBranch R model :=
  cnfBareNegativeBranch_forces_imageSeparatorBranch
    (R := R)
    (model := model)
    hNotPositive
\end{lstlisting}

\begin{lstlisting}[basicstyle=\ttfamily\scriptsize]
def CnfSATDirectPositiveEndpointOccupation : Prop :=
  CnfPositiveEndpoint

def CnfSATIndependentSeparatorEndpointStatusDiscriminator
    (model : CnfEncodedCandidateModel) : Prop :=
  exists classifier : CnfCandidateStatusClassifier model,
    CnfClassifierSeparatesPolynomialCandidates model classifier /\
      Nonempty (CnfClassifierIndependentSameDomain model classifier)

theorem cnfPositiveEndpoint_directEndpointOccupation
    (hPositive : CnfPositiveEndpoint) :
    CnfSATDirectPositiveEndpointOccupation :=
  hPositive

theorem cnfSATImageSeparatorBranch_requires_independentSeparatorDiscriminator
    {Act Object : Type}
    {R : MinimalConditionsForAdmissibleConstruction.ConstructionRegime Act Object}
    {model : CnfEncodedCandidateModel}
    (hIndependent :
      CnfSeparatingClassifierIsIndependentSameDomain model)
    (hSep : CnfSATImageSeparatorBranch R model) :
    CnfSATIndependentSeparatorEndpointStatusDiscriminator model := by
  have hSameDomain : CnfSameDomainSeparator :=
    (cnfDirectGateLowerBoundResidualTarget_iff_sameDomainSeparator).1 hSep
  let classifier := cnfClassifierOfSameDomainSeparator model hSameDomain
  have hSeparates :
      CnfClassifierSeparatesPolynomialCandidates model classifier :=
    cnfClassifierSeparates_of_sameDomainSeparator model hSameDomain
  exact Exists.intro classifier
    (And.intro hSeparates (hIndependent classifier hSeparates))
\end{lstlisting}

\begin{lstlisting}[basicstyle=\ttfamily\scriptsize]
theorem cnfBareNegativeBranch_requires_independentSeparatorDiscriminator
    {Act Object : Type}
    {R : MinimalConditionsForAdmissibleConstruction.ConstructionRegime Act Object}
    {model : CnfEncodedCandidateModel}
    (hIndependent :
      CnfSeparatingClassifierIsIndependentSameDomain model)
    (hBare : CnfSATBareNegativeBranch R model) :
    CnfSATIndependentSeparatorEndpointStatusDiscriminator model :=
  cnfSATImageSeparatorBranch_requires_independentSeparatorDiscriminator
    (R := R)
    hIndependent
    (cnfBareNegativeBranch_forces_imageSeparatorBranch hBare)

theorem cnfSATImageSeparatorBranch_impossible_of_noIndependentDiscriminator
    {Act Object : Type}
    {R : MinimalConditionsForAdmissibleConstruction.ConstructionRegime Act Object}
    {model : CnfEncodedCandidateModel}
    (hNoIndependent : CnfNoIndependentSeparatingClassifier model)
    (hIndependent :
      CnfSeparatingClassifierIsIndependentSameDomain model) :
    Not (CnfSATImageSeparatorBranch R model) := by
  intro hSep
  exact Exists.elim
    (cnfSATImageSeparatorBranch_requires_independentSeparatorDiscriminator
      (R := R)
      hIndependent hSep)
    (fun classifier hClassifier =>
      hNoIndependent classifier hClassifier.1 hClassifier.2)
\end{lstlisting}

\begin{lstlisting}[basicstyle=\ttfamily\scriptsize]
theorem cnfBareNegativeBranch_impossible_of_noIndependentDiscriminator
    {Act Object : Type}
    {R : MinimalConditionsForAdmissibleConstruction.ConstructionRegime Act Object}
    {model : CnfEncodedCandidateModel}
    (hNoIndependent : CnfNoIndependentSeparatingClassifier model)
    (hIndependent :
      CnfSeparatingClassifierIsIndependentSameDomain model) :
    Not (CnfSATBareNegativeBranch R model) := by
  intro hBare
  exact
    (cnfSATImageSeparatorBranch_impossible_of_noIndependentDiscriminator
      (R := R)
      hNoIndependent hIndependent)
      (cnfBareNegativeBranch_forces_imageSeparatorBranch hBare)

theorem cnfNotPositiveEndpoint_impossible_of_noIndependentDiscriminator
    {Act Object : Type}
    {R : MinimalConditionsForAdmissibleConstruction.ConstructionRegime Act Object}
    {model : CnfEncodedCandidateModel}
    (hNoIndependent : CnfNoIndependentSeparatingClassifier model)
    (hIndependent :
      CnfSeparatingClassifierIsIndependentSameDomain model)
    (hNotPositive : Not CnfPositiveEndpoint) :
    False :=
  (cnfBareNegativeBranch_impossible_of_noIndependentDiscriminator
    (R := R)
    (model := model)
    hNoIndependent hIndependent)
    hNotPositive
\end{lstlisting}

\begin{lstlisting}[basicstyle=\ttfamily\scriptsize]
theorem cnfSATInPolyTime_of_noIndependentDiscriminator
    {Act Object : Type}
    {R : MinimalConditionsForAdmissibleConstruction.ConstructionRegime Act Object}
    {model : CnfEncodedCandidateModel}
    (hNoIndependent : CnfNoIndependentSeparatingClassifier model)
    (hIndependent :
      CnfSeparatingClassifierIsIndependentSameDomain model) :
    CnfSATInPolyTime :=
  cnfPositiveEndpoint_of_noIndependentDiscriminator
    (R := R)
    (model := model)
    hNoIndependent hIndependent

theorem cnfSATInPolyTime_of_context_noIndependentDiscriminator
    {Act Object : Type}
    {R : MinimalConditionsForAdmissibleConstruction.ConstructionRegime Act Object}
    {model : CnfEncodedCandidateModel}
    (_hCtx : CnfSATClayEndpointImageContext R model)
    (_hFixed : CnfSATFixedEndpointDomain R)
    (_hModel : CnfSATContextBoundCNFModel R model)
    (hNoIndependent : CnfNoIndependentSeparatingClassifier model)
    (hIndependent :
      CnfSeparatingClassifierIsIndependentSameDomain model) :
    CnfSATInPolyTime :=
  cnfSATInPolyTime_of_noIndependentDiscriminator
    (R := R)
    (model := model)
    hNoIndependent hIndependent
\end{lstlisting}


\begin{lstlisting}[basicstyle=\ttfamily\scriptsize]
def CnfSATBookkeepingNegativeOnly
    {Act Object : Type}
    (_R : MinimalConditionsForAdmissibleConstruction.ConstructionRegime Act Object)
    (_model : CnfEncodedCandidateModel) : Prop :=
  Prop

def CnfSATCarrierShift
    {Act Object : Type}
    (_R : MinimalConditionsForAdmissibleConstruction.ConstructionRegime Act Object)
    (_model : CnfEncodedCandidateModel) : Prop :=
  Prop

def CnfSATOfficialNegativeEndpointUse
    {Act Object : Type}
    (_R : MinimalConditionsForAdmissibleConstruction.ConstructionRegime Act Object)
    (_model : CnfEncodedCandidateModel) : Prop :=
  CnfSATBareNegativeBranch _R _model

def CnfSATNegativeOccupationExhaustion
    {Act Object : Type}
    (_R : MinimalConditionsForAdmissibleConstruction.ConstructionRegime Act Object)
    (_model : CnfEncodedCandidateModel) : Prop :=
  CnfSATImageSeparatorBranch _R _model \/
  CnfOrdinaryNegativeAlternative _R _model \/
  CnfSATBookkeepingNegativeOnly _R _model \/
  CnfSATCarrierShift _R _model
\end{lstlisting}

\begin{lstlisting}[basicstyle=\ttfamily\scriptsize]
theorem cnfSATNegativeOccupation_nonoptional
    {Act Object : Type}
    {R : MinimalConditionsForAdmissibleConstruction.ConstructionRegime Act Object}
    {model : CnfEncodedCandidateModel}
    (hOfficial : CnfSATOfficialNegativeEndpointUse R model)
    (hFixed : CnfSATFixedEndpointDomain R)
    (hModel : CnfSATContextBoundCNFModel R model)
    (hAMetric : AMetricBoundary R) :
    CnfSATImageSeparatorBranch R model :=
  -- formal proof in SATOperatorProofQueue.lean
\end{lstlisting}

\begin{lstlisting}[basicstyle=\ttfamily\scriptsize]
inductive CnfSATEndpointStatus where
  | positive
  | separator

def CnfSATGovernedEndpointUse
    {Act Object : Type}
    (_R : MinimalConditionsForAdmissibleConstruction.ConstructionRegime Act Object)
    (_model : CnfEncodedCandidateModel) : Prop :=
  Prop

theorem cnfSATNegativeGovernedEndpointUse_has_separatorStatus
    {Act Object : Type}
    {R : MinimalConditionsForAdmissibleConstruction.ConstructionRegime Act Object}
    {model : CnfEncodedCandidateModel}
    (hGov : CnfSATGovernedEndpointUse R model)
    (hNotPositive : Not CnfPositiveEndpoint) :
    CnfSATImageSeparatorEndpointUse R model :=
  -- formal proof in SATOperatorProofQueue.lean

theorem cnfSATImageSeparatorBranch_of_negativeGovernedEndpointUse
    {Act Object : Type}
    {R : MinimalConditionsForAdmissibleConstruction.ConstructionRegime Act Object}
    {model : CnfEncodedCandidateModel}
    (hGov : CnfSATGovernedEndpointUse R model)
    (hNotPositive : Not CnfPositiveEndpoint) :
    CnfSATImageSeparatorBranch R model :=
  -- formal proof in SATOperatorProofQueue.lean
\end{lstlisting}


\subsection{Focused axiom-output table}

The focused audit files in the release issue \texttt{\#print axioms} commands for the AASC foundation anchors and the SAT endpoint bridge anchors below. The release records the following focused axiom surface for these proof-spine theorems: the AASC foundation anchors report no project axioms, and the SAT bridge anchors report only the standard Lean/classical/mathlib base dependencies \LeanName{propext}, \LeanName{Classical.choice}, and \LeanName{Quot.sound}. No project-specific axiom, \LeanName{sorry}, \LeanName{admit}, or \LeanName{unsafe} declaration is used by these anchors.

The model-adequacy audit in \LeanName{AnticheckerReduction.lean} names the ordinary coding bridge consumed by the SAT-native route. It records coverage, sound decoding, failure adequacy, candidate-image exactness, and the positive/negative endpoint equivalences for encoded deterministic polynomial-time CNF-SAT candidates. The final Lean audit also names the SAT-native form of the remaining bridge objection. The lower-bound and theorem-level discriminator anchors record that the bare negative branch gives lower-bound normal form over the fixed encoded candidate image, and that the resulting discriminator is theorem-level rather than algorithmic. The invariant-use and candidate-image-use classifications separate legitimate proof support from forbidden endpoint-status governance. The closing hinge is formalized by the following Lean anchors:
\begin{lstlisting}[basicstyle=\ttfamily\scriptsize]
CnfSATCandidateImageExclusionUseKind
CnfSATCandidateImageExclusionUseClassification
CnfSATEndpointStatusGovernanceByCandidateImageExclusion
cnfSATEndpointResolvingNonGovernance_hiddenFifthCase_impossible
cnfSATOfficialNegativeEndpointUse_endpointStatusGovernance
cnfSATEndpointResolvingNegativeTheorem_is_endpointStatusGovernance
\end{lstlisting}

\begin{center}
\scriptsize
\begin{tabularx}{0.96\linewidth}{>{\raggedright\arraybackslash}p{0.50\linewidth}>{\raggedright\arraybackslash}X}
\toprule
Lean theorem & Focused axiom output \\
\midrule
\begin{tabular}[t]{@{}l@{}}\LeanName{cnfBareNegativeBranch}\\\LeanName{forces_imageSeparatorBranch}\end{tabular} & \LeanName{[propext, Classical.choice, Quot.sound]} \\
\begin{tabular}[t]{@{}l@{}}\LeanName{cnfImageSeparatorBranch}\\\LeanName{of_not_positiveEndpoint}\end{tabular} & \LeanName{[propext, Classical.choice, Quot.sound]} \\
\begin{tabular}[t]{@{}l@{}}\LeanName{cnfSATImageSeparatorBranch}\\\LeanName{endpointUse}\end{tabular} & \LeanName{[propext, Classical.choice, Quot.sound]} \\
\begin{tabular}[t]{@{}l@{}}\LeanName{cnfSATImageSeparatorBranch}\\\LeanName{requires_independentSeparatorDiscriminator}\end{tabular} & \LeanName{[propext, Classical.choice, Quot.sound]} \\
\begin{tabular}[t]{@{}l@{}}\LeanName{cnfSATImageSeparatorBranch}\\\LeanName{impossible_of_noIndependentDiscriminator}\end{tabular} & \LeanName{[propext, Classical.choice, Quot.sound]} \\
\begin{tabular}[t]{@{}l@{}}\LeanName{cnfBareNegativeBranch}\\\LeanName{requires_independentSeparatorDiscriminator}\end{tabular} & \LeanName{[propext, Classical.choice, Quot.sound]} \\
\begin{tabular}[t]{@{}l@{}}\LeanName{cnfBareNegativeBranch}\\\LeanName{impossible_of_noIndependentDiscriminator}\end{tabular} & \LeanName{[propext, Classical.choice, Quot.sound]} \\
\begin{tabular}[t]{@{}l@{}}\LeanName{cnfNotPositiveEndpoint}\\\LeanName{impossible_of_noIndependentDiscriminator}\end{tabular} & \LeanName{[propext, Classical.choice, Quot.sound]} \\
\begin{tabular}[t]{@{}l@{}}\LeanName{cnfSATInPolyTime}\\\LeanName{of_context_noIndependentDiscriminator}\end{tabular} & \LeanName{[propext, Classical.choice, Quot.sound]} \\
\begin{tabular}[t]{@{}l@{}}\LeanName{cnfSATNegativeOccupation}\\\LeanName{exhaustion}\end{tabular} & \LeanName{[propext, Classical.choice, Quot.sound]} \\
\begin{tabular}[t]{@{}l@{}}\LeanName{cnfSATNegativeOccupation}\\\LeanName{nonoptional}\end{tabular} & \LeanName{[propext, Classical.choice, Quot.sound]} \\
\begin{tabular}[t]{@{}l@{}}\LeanName{cnfSATGovernedEndpointUse}\\\LeanName{bivalent}\end{tabular} & \LeanName{[propext, Classical.choice, Quot.sound]} \\
\begin{tabular}[t]{@{}l@{}}\LeanName{cnfSATNegativeGovernedEndpointUse}\\\LeanName{has_separatorStatus}\end{tabular} & \LeanName{[propext, Classical.choice, Quot.sound]} \\
\begin{tabular}[t]{@{}l@{}}\LeanName{cnfSATImageSeparatorBranch}\\\LeanName{of_negativeGovernedEndpointUse}\end{tabular} & \LeanName{[propext, Classical.choice, Quot.sound]} \\
\begin{tabular}[t]{@{}l@{}}\LeanName{cnfSATCandidateModelCoverage}\end{tabular} & \LeanName{[propext, Classical.choice, Quot.sound]} \\
\begin{tabular}[t]{@{}l@{}}\LeanName{cnfSATCandidateModelSoundness}\end{tabular} & \LeanName{[propext, Classical.choice, Quot.sound]} \\
\begin{tabular}[t]{@{}l@{}}\LeanName{cnfSATCandidateFailureAdequacy}\end{tabular} & \LeanName{[propext, Classical.choice, Quot.sound]} \\
\begin{tabular}[t]{@{}l@{}}\LeanName{cnfSATCandidateImageExactness}\end{tabular} & \LeanName{[propext, Classical.choice, Quot.sound]} \\
\begin{tabular}[t]{@{}l@{}}\LeanName{cnfEncodedCandidateModelAdequate}\end{tabular} & \LeanName{[propext, Classical.choice, Quot.sound]} \\
\begin{tabular}[t]{@{}l@{}}\LeanName{cnfSATInPolyTime_iff_exists}\\\LeanName{_encoded_nonfailing_candidate}\end{tabular} & \LeanName{[propext, Classical.choice, Quot.sound]} \\
\begin{tabular}[t]{@{}l@{}}\LeanName{cnfSATNotInPolyTime_iff}\\\LeanName{_all_encoded_candidates_fail}\end{tabular} & \LeanName{[propext, Classical.choice, Quot.sound]} \\
\begin{tabular}[t]{@{}l@{}}\LeanName{cnfSATBareNegativeBranch}\\\LeanName{standardLowerBoundNormalForm}\end{tabular} & \LeanName{[propext, Classical.choice, Quot.sound]} \\
\begin{tabular}[t]{@{}l@{}}\LeanName{cnfSATBareNegativeBranch}\\\LeanName{theoremLevelDiscriminator}\end{tabular} & \LeanName{[propext, Classical.choice, Quot.sound]} \\
\begin{tabular}[t]{@{}l@{}}\LeanName{cnfSATOfficialNegativeEndpointUse}\\\LeanName{endpointStatusGovernance}\end{tabular} & \LeanName{[propext, Classical.choice, Quot.sound]} \\
\begin{tabular}[t]{@{}l@{}}\LeanName{cnfSATEndpointResolvingNonGovernance}\\\LeanName{hiddenFifthCase_impossible}\end{tabular} & \LeanName{[propext, Classical.choice, Quot.sound]} \\
\begin{tabular}[t]{@{}l@{}}\LeanName{cnfSATEndpointResolvingNegativeTheorem}\\\LeanName{is_endpointStatusGovernance}\end{tabular} & \LeanName{[propext, Classical.choice, Quot.sound]} \\
\bottomrule
\end{tabularx}
\end{center}

\subsection{Audit scope and boundary}

The Lean archive formalizes the AASC foundation audit surface consumed by the SAT route together with the SAT-operator bridge and negative-branch closeout used in the manuscript proof spine. The SAT-operator proof queue also includes negative-occupation-exhaustion anchors and endpoint-status-bivalence anchors, corresponding to commits \LeanName{282d4ad} and \LeanName{a938648} in the public Lean repository; commit \LeanName{f2df564} exposes the combined full-stack AASC/SAT audit check. It does not attempt to formalize the external Cook--Levin/Karp theorem inside Lean.

The no-one-step-retyping lock and the native lower-bound falsification collapse are manuscript-level defensive expansions of route roles already represented on the SAT audit surface: lower-bound normal form, candidate-image exclusion, image-separator branch, official negative endpoint use, endpoint-status governance, separator discriminator induction, no-independent closure, and local reductio discharge. They add no new Cook--Levin/Karp formalization burden and do not alter the Lean proof-spine endpoint. Optional Lean-side mirror names, not required for the manuscript proof, include:
\begin{lstlisting}[basicstyle=\ttfamily\scriptsize]
CnfSATNativeLowerBoundFalsificationPackage
cnfSATNativeLowerBoundBelowEndpointUse
cnfSATDeployedNativeLowerBound_endpointUse
cnfSATLowerBoundBridgeNeutralized
cnfSATOrdinaryNativeLowerBoundFalsificationCollapse
cnfSATNoSixthLowerBoundRole
cnfSATRouteLocusOrdinaryNativeLowerBound
\end{lstlisting}

Its support role is local and reproducibility-oriented: it checks that the AASC foundation anchors and the formal SAT-side branch chain used by the manuscript have the advertised Lean anchors and that the audit surface is free of live project-level incomplete declarations.

\section{Standard CNF-SAT to P=NP Correspondence}
The manuscript proof derives \(\CnfSATP\). Here ``bounded-CNF'' means finitely encoded CNF formulas of arbitrary instance size over a finite Boolean alphabet; it does not mean a fixed-width, fixed-size, bounded-variable, fixed-clause, bounded-depth, promise-restricted, parameterized, or otherwise tractable restricted CNF fragment. The equality \(\Pclass=\NPclass\) then uses the standard Cook--Levin/Karp endpoint correspondence: arbitrary finite CNF-SAT is NP-complete under polynomial-time many-one reductions; every \(L\in \NPclass\) reduces to CNF-SAT; composition of a polynomial reduction with a polynomial SAT decider gives a polynomial decider for \(L\); and \(\Pclass\subseteq\NPclass\). Hence \(\NPclass\subseteq\Pclass\) and \(\Pclass=\NPclass\). This correspondence is standard Cook--Levin/Karp background. It is not the proof of separator exclusion; AASC supplies that proof.

\section{Reference Integrity Appendix}
The proof ladder is ordered so that each theorem cites prior results only. The final endpoint theorem cites the context definition, official endpoint-evaluation definition, the standard encoded-model instantiation, the encoded candidate model-adequacy packet, positive direct-occupation theorem, model binding, the branch-local SAT image-separator bridge, raw complementarity, the raw-negative-to-image bridge, image-separator exclusion, and the standard CNF-SAT completeness bridge. No theorem cites itself.


\section{Theorem Ladder and Anti-Circularity Audit}

\subsection*{Theorem ladder}
\begin{enumerate}
\item Mathematical-status lock and Lean support boundary fix the proof as a manuscript theorem chain with Lean as audit support.
\item Kernel-neutral target adequacy fixes target determinacy, step evaluability, act-time finality, and same-regime fidelity.
\item Non-degenerate construction forces \(\Adm,\St,\RefK,\Irr\), and endpoint use is governed construction.
\item Ordinary official endpoint resolution is endpoint use: ordinary theoremhood used as official resolution fixes the target, carrier, branch, theorem-status, downstream reuse, act-time finality, and redescription-stable endpoint role.
\item Ordinary theoremhood does not shield official endpoint resolution from endpoint-role discipline; it certifies proof legitimacy but not endpoint-interface equivalence.
\item Exhaustiveness of endpoint-role classification on the projected SAT carrier shows that theorem-bearing, endpoint-resolving, same-domain, non-positive, non-coequal, non-support, non-bookkeeping, non-repair, non-carrier-shifting work has no stable middle role.
\item No one-step retyping blocks any direct passage from raw \(\SepBare\), \(\neg\PosRaw\), or lower-bound normal form to \(\SepDisc(model)\) without candidate-image exclusion and endpoint-use promotion.
\item The native lower-bound falsification package \(\LBNatSAT(R,model)\) isolates ordinary lower-bound falsification as a named hostile role.
\item The SAT endpoint-impact bundle \(\EndpointImpactSAT(R,model)\) separates inert lower-bound bookkeeping from endpoint-impacting negative use.
\item The ordinary native lower-bound falsification collapse exhausts endpoint-deployed lower-bound use into support, positive bridge-neutralization, bookkeeping, carrier shift, or \(\SepDisc(model)\).
\item The no-sixth-lower-bound-role corollary closes the live exact-complement subproof against ordinary native lower-bound escape.
\item Ordinary theorem-bearing non-occupation is not a fifth endpoint role after native lower-bound collapse, SAT projection, and official endpoint use; proof legitimacy does not neutralize endpoint-role classification.
\item Exhaustiveness of endpoint-role classification under official CNF-SAT resolution packages the projection, official-use, single-interface, and trichotomy results into a one-step referee closure: ordinary official negative resolution through CNF-SAT yields \(\SepDisc(model)\).
\item Endpoint-status work is endpoint-admissibility-relevant; separator discrimination induced by global or local endpoint-resolving separator occupation changes endpoint standing, reportability, downstream reuse, or branch exclusion and is therefore in the exact class governed by the kernel consequence layer.
\item A+ consequences bind ordinary official endpoint resolution: once ordinary practice supplies endpoint use and the trichotomy identifies an independent role, the no-independent-discriminator closure applies.
\item Remaining same-mode objections are exhausted by the projection link, the endpoint-discriminator trichotomy, or the \(\Aplus\) no-independent-discriminator consequence.
\item Raw bivalence is separated from endpoint standing: a hypothetical negative case is not official endpoint resolution.
\item Truth branch, epistemic status, and endpoint outcome are separated: unresolved status is not a third endpoint outcome.
\item Endpoint completion is derived, not assumed: raw bivalence is combined with local countercase discipline and exclusion of local separator occupation.
\item Exact-complement reductio annotation is proof-rule discipline: opening \([\SepBare]^i\) as the exact complement of \(\PosRaw\) generates \(\ReductioCountercase_R(\SepBare;\PosRaw)\) as a proof-act annotation, not as an added premise; the annotation-discharge theorem discharges the original assumption, not a strengthened conjunction.
\item Endpoint closure is not proof-availability bivalence: raw bivalence is used, and the local negative countercase is routed directly to local separator occupation, not to global endpoint outcome.
\item Endpoint counterforce by the raw negative branch induces local separator occupation \(\LocalSepImgOcc(R,model)\), which is locally impossible by the same separator-discriminator contradiction.
\item The local SAT kernel packet records the SAT-facing K1--K13 consequences.
\item Weakening-resistance shows that K5, K6, K11, and K13 cannot be strictly weakened while preserving same-carrier endpoint use and allowing separator governance.
\item Encoded candidate model adequacy fixes coverage, soundness, failure adequacy, and candidate-image exactness for the unrestricted finite CNF-SAT carrier.
\item Raw complementarity gives \(\SepBare\leftrightarrow\neg\PosRaw\).
\item Official SAT endpoint-closure evaluation supplies only local countercase standing for the reductio assumption; the negative branch is not promoted to global endpoint outcome.
\item The SAT-native bridge routes \(\SepBare\) to standard lower-bound normal form and candidate endpoint-image exclusion; official endpoint-resolution use promotes that content to \(\SepImgOcc\), and separator occupation explicitly contains the support-level candidate-image exclusion content.
\item The single occupation interface theorem shows that the fixed SAT endpoint image is occupied exactly by encoded non-failing polynomial-time candidate occupation.
\item Candidate-image carrier adequacy pairs the official raw split with the fixed image: \(\PosRaw\leftrightarrow\SATEndpointImageOcc(model)\) and \(\SepBare\leftrightarrow\CandImageExcl(model)\).
\item Single-interface typing is non-explosive: it classifies non-occupation but does not by itself exclude lower-bound theorems, support-level candidate-image exclusion, or negative complexity results.
\item Candidate-image exclusion is non-occupation of that image, not coequal occupation; coequal negative outcome belongs to the symmetric outcome carrier \(\OutcomePvNP\), not to the fixed candidate-image carrier.
\item Ordinary negation remains ordinary negation; endpoint-bearing negation is role-classified, and ordinary official negative SAT resolution is non-positive status work on the fixed candidate image.
\item The separator branch is not a lawful coequal endpoint role inside the same positive endpoint-image carrier.
\item Endpoint-resolving candidate-image exclusion is non-positive endpoint-status work; endpoint-resolving non-governance is the hidden fifth case and is impossible.
\item Candidate-image exclusion remains lawful as proof support; only endpoint-resolving image-separator use invokes \(\SepIndBridge(model)\) and induces \(\SepDisc(model)\), while image-separator endpoint use also gives \(\NoSepDisc(model)\).
\item The contradiction excludes both global separator occupation \(\SepImgOcc\) and local separator occupation \(\LocalSepImgOcc(R,model)\), hence excludes negative endpoint counterforce; it does not treat raw truth or a raw hypothetical branch as endpoint resolution or as a global endpoint outcome.
\item Bare complementarity yields \(\PosRaw\), i.e. \(\CnfSATP\).
\item Cook--Levin/Karp gives \(P=NP\).
\end{enumerate}

\subsection*{Anti-circularity audit}
{\small
\begin{longtable}{>{\raggedright\arraybackslash}p{0.34\linewidth}>{\raggedright\arraybackslash}p{0.60\linewidth}}
\toprule
Potential circle or misreading & Audit response \\
\midrule
AASC is merely philosophical. & The manuscript uses AASC as a mathematical constraint formalism. A same-mode objection must identify a failed definition, theorem, dependency claim, or bridge inference. \\
Kernel governance is optional. & Kernel roles are forced by target determinacy, step evaluability, act-time finality, and same-regime fidelity. Denying the kernel denies non-degenerate endpoint use or supplies a faithful weaker regime. \\
Endpoint target assumes \(P=NP\). & The official target fixes the carrier and endpoint slots. It does not assert \(\PosRaw\). \\
Official endpoint evaluation smuggles in endpoint completeness. & No. \(\OfficialSATEval(R,model)\) fixes target, carrier, raw branch slots, and local countercase discipline for endpoint closure. It does not assert endpoint-complete outcome bivalence or either branch. Raw bivalence is used separately; a local assumption of \(\SepBare\) in the reductio is first marked as the exact-complement reductio assumption \(\ReductioCountercase_R(\SepBare;\PosRaw)\), then as \(\LocalCountercase_R(\SepBare)\), and only that local countercase has endpoint counterforce by Theorem~\ref{thm:local-countercase-has-endpoint-counterforce}. \\
Raw proof-by-cases assumption is promoted to official resolution. & No. The proof separates raw bivalence, exact-complement reductio use, local countercase standing, endpoint counterforce, local separator occupation, and global endpoint outcome. A raw assumption of \(\SepBare\) is not official endpoint resolution; only when it is introduced as \(\ReductioCountercase_R(\SepBare;\PosRaw)\) does it yield \(\LocalCountercase_R(\SepBare)\). That local standing yields local separator occupation and a local discriminator contradiction, without promoting the assumption to \(\NegEndpointOutcome_R\). \\
Merely propositional assumptions are retyped as endpoint countercases. & No. Inert propositional assumptions remain inert. The theorem \(\ReductioCountercase_R(\SepBare;\PosRaw)\Rightarrow\LocalCountercase_R(\SepBare)\) applies only when \(\SepBare\) is introduced as the live exact-complement assumption whose survival would block \(\PosRaw\) in the endpoint reductio. \\
\(\ReductioCountercase_R(\SepBare;\PosRaw)\) is an added premise alongside \(\SepBare\). & No. Lemma~\ref{lem:exact-complement-reductio-annotation-rule} states that \(\ReductioCountercase_R(\SepBare;\PosRaw)\) is the proof-act annotation of the discharged assumption \([\SepBare]^i\). Theorem~\ref{thm:exact-complement-annotation-discharge} states that contradiction in the annotated subproof discharges \([\SepBare]^i\), not a strengthened assumption \(\SepBare\wedge\ReductioCountercase_R(\SepBare;\PosRaw)\). \\
Endpoint closure assumes proof availability. & No. The proof uses raw bivalence plus exact-complement reductio discipline. Open or unresolved status remains epistemic; it is not a third endpoint branch. The proof does not assume that an endpoint outcome was available before the closure argument. \\
\(\ClayCtx\) smuggles positivity. & The context non-smuggling theorem states that \(\ClayCtx\) fixes only carrier and branch slots. \\
Bounded-CNF means a restricted SAT fragment. & Definition~\ref{def:bounded-cnf-clarification} states that it means arbitrary finite encoded CNF-SAT, not fixed-width or promise SAT. \\
The encoded model omits candidates. & Model adequacy requires coverage, soundness, failure adequacy, and image exactness. A valid objection must identify a concrete failure of one of these clauses. \\
Ordinary negative theoremhood is banned. & Theorems~\ref{thm:negative-theorem-not-defective} and~\ref{thm:ordinary-sat-falsity-role-classified} explicitly allow ordinary negative theoremhood. \\
Ordinary negative resolution stays inside ordinary mathematical practice and therefore avoids \(\Aplus\). & Ordinary theoremhood as theoremhood is harmless. But ordinary theoremhood used as official endpoint resolution fixes target, carrier, branch, theorem-status, downstream reuse, act-time finality, and redescription-stable endpoint role. By Theorem~\ref{thm:ordinary-official-resolution-is-endpoint-use}, that act is \(\EndpointUse_R\). \\
The manuscript retypes ordinary negative proof as forbidden governance. & No. It first distinguishes ordinary theoremhood from official endpoint resolution. The negative theorem remains ordinary theoremhood as proof legitimacy; when used as official endpoint resolution on the fixed candidate image, its endpoint role is non-positive status work by Theorem~\ref{thm:ordinary-official-negative-sat-resolution-nonpositive-status-work}. \\
\(\SepBare\) is retyped by fiat. & No one-step retyping is used: \(\SepBare\) passes through standard lower-bound normal form and candidate endpoint-image exclusion before endpoint-governance classification. \\
\(\SepImgOcc\) is an arbitrary extra layer. & \(\SepImgOcc\) is the endpoint-occupation form of candidate-image exclusion under official endpoint-resolution use; support-level lower-bound content remains \(\CandImageExcl(model)\). \\
Candidate-image exclusion is an algorithmic separator. & The theorem-level discriminator is not a witness selector or algorithmic separator; it classifies endpoint status for already-fixed candidate roles. \\
Candidate-image exclusion is automatically forbidden. & No. Candidate-image exclusion may be proof support, a lower-bound lemma, or ordinary theorem content. It invokes the independence bridge only under same-carrier endpoint-resolving use. The separate unpacking lemma says only that \(\SepImgOcc\) contains \(\CandImageExcl(model)\), not conversely. \\
\(\SepImgOcc\) collapses support use into endpoint use. & No. Support-level content is \(\CandImageExcl(model)\). The symbol \(\SepImgOcc\) is reserved for endpoint-image branch occupation under official endpoint-resolution use on the fixed carrier. \\
The negative branch is ordinary coequal resolution. & Correct only on the symmetric two-output problem carrier \(\OutcomePvNP\). On the fixed SAT endpoint-image carrier, the negative branch is non-occupation of the positive candidate image, not coequal occupation of that image. \\
The positive route is assumed as the sole admitted interface. & Not assumed. The single-interface theorem proves it from endpoint-image typing: occupation of the fixed SAT image is exactly encoded non-failing polynomial-time candidate occupation. \\
Standard practice does not call this governance. & The proof uses governance extensionally: status-determining work on the same fixed image that is not equivalent to positive occupation, support, bookkeeping, carrier shift, repair, or coequal image occupation. \\
Ordinary theorem-bearing non-occupation avoids the discriminator classification. & No. Ordinary theoremhood gives proof legitimacy. Once the theorem is used as official endpoint resolution and projected onto the fixed SAT candidate image, its endpoint content is candidate-image non-occupation. Since it is not support, bookkeeping, repair, carrier shift, positive occupation, or coequal image occupation, Theorem~\ref{thm:ordinary-native-lb-falsification-collapse} first exhausts the ordinary native lower-bound reading; after support, positive bridge-neutralization, bookkeeping, and carrier shift are removed, Theorem~\ref{thm:projected-sat-role-exhaustiveness} classifies the remaining act as non-positive endpoint-status work, and Theorem~\ref{thm:non-positive-status-work-induces-sepdisc} identifies that role as independent same-domain separator-discriminator work. Theorem~\ref{thm:sat-sepdisc-endpoint-adm-relevant} further shows that this discriminator is endpoint-admissibility-relevant, so it lies in the class controlled by the kernel consequence layer rather than merely in a descriptive label layer. \\
SepDisc is only a descriptive endpoint label, not an admissibility-relevant invariant. & No. When induced by global or local endpoint-resolving separator occupation, \(\SepDisc(model)\) determines which branch may stand, be reported, reused downstream, or excluded as the endpoint result. By Theorem~\ref{thm:sat-sepdisc-endpoint-adm-relevant}, it is endpoint-admissibility-relevant in the sense of Definition~\ref{def:endpoint-admissibility-relevance}. \\
Official CNF-SAT negative resolution remains an uninterpreted symmetric answer. & No. Once the standard CNF-SAT route is used for official endpoint resolution, the symmetric negative answer is projected to candidate-image non-occupation on the fixed carrier. Lemma~\ref{lem:official-cnf-sat-resolution-exhaustiveness} packages the resulting elimination: official negative resolution through CNF-SAT yields \(\SepDisc(model)\), unless projection, trichotomy, or the \(\Aplus\) consequence layer is denied. \\
Ordinary practice accepts endpoint use and independence but rejects \(\Aplus\). & No. Ordinary practice need not use AASC vocabulary, but ordinary official resolution supplies \(\EndpointUse_R\). By Corollary~\ref{cor:aplus-binds-ordinary-official-resolution}, endpoint use on a fixed carrier triggers the kernel-forced \(\Aplus\) no-independent-discriminator consequence for the identified independent role. \\
Proof legitimacy neutralizes endpoint-role classification. & No. Proof legitimacy and endpoint role are different dimensions. Legitimacy makes the act theorem-bearing; endpoint-role classification determines what the act does on the fixed carrier once it is used for official resolution. \\
The manuscript forces standard practice to accept AASC typing. & No. It asks only for role classification after standard SAT projection and official endpoint use. The vocabulary is AASC; the load-bearing facts are projection to candidate-image non-occupation, exclusion of non-governing roles, and the endpoint-discriminator trichotomy. \\
Single-interface typing proves every positive existence endpoint. & No. The single-interface theorem only identifies the occupant type of this SAT endpoint image. Endpoint closure additionally requires endpoint-resolving separator occupation to induce \(\SepDisc(model)\) while endpoint use yields \(\NoSepDisc(model)\). \\
The candidate-image carrier is a one-sided surrogate. & No. The carrier-adequacy theorem proves that the same image represents the official raw split: positive occupation corresponds to \(\PosRaw\), and image non-occupation corresponds to \(\SepBare\). \\
The proof bans lower bounds globally. & The non-explosiveness theorem permits lower bounds, barriers, obstructions, and proof support. Only same-carrier endpoint-resolving separator governance is excluded. \\
K5 bans multiple candidate procedures. & K5 bans plural endpoint-governing interiors, not multiple candidates, encodings, proof routes, or support structures. \\
K13 bans open problem status. & K13 classifies open, unknown, deferred, and support statuses as epistemic or report statuses, not endpoint truth branches. \\
Target-slot fixation smuggles in K13. & Target-slot fixation identifies what is evaluated. K13 separately proves that hidden-fifth, deferred, and bookkeeping statuses cannot occupy endpoint truth. \\
Minimal report evaluability smuggles in K11. & Minimal report evaluability classifies a report by role. K11 separately imposes no-overreporting: endpoint force may not exceed carrier and bridge support. \\
Cook--Levin/Karp is used as branch-exclusion machinery. & Cook--Levin/Karp appears only after \(\CnfSATP\) is proved, as the standard endpoint correspondence to \(P=NP\). \\
G\"odel-style true-but-unprovable separator objections survive the endpoint proof. & No. Same-domain incompleteness objections matter only if they change endpoint standing, reportability, downstream reuse, admissible continuation, or branch exclusion. Then they are endpoint-admissibility-relevant gate work and are exhausted by the G\"odel non-instantiability route analysis; otherwise they are endpoint-inert. \\
Lean is substituted for the proof. & The manuscript is the proof. Lean records a reproducibility audit of the SAT bridge and kernel consequence surface. \\
\bottomrule
\end{longtable}
}

\section{Notation Ledger}
\begin{longtable}{>{\raggedright\arraybackslash}p{0.21\linewidth}>{\raggedright\arraybackslash}p{0.54\linewidth}>{\raggedright\arraybackslash}p{0.17\linewidth}}
\toprule
Symbol & Meaning & Layer \\
\midrule
\(\ClayCtx\) & fixed SAT endpoint carrier and branch-slot context & context \\
\(\PosRaw\) & \(\CnfSATP\) & raw positive \\
\(\SepBare\) & raw complement branch \(\neg\PosRaw\) & bare negative \\
\(\SepImg\) & Lean-side formal SAT image-separator branch \(\CnfImageSep(R,model)\); in the manuscript proof it is separated into support-level \(\CandImageExcl(model)\) and endpoint occupation \(\SepImgOcc\) & Lean bridge \\
\(\LBNatSAT(R,model)\) & native lower-bound falsification package preserving the fixed unrestricted finite CNF-SAT carrier, encoded candidate image, universal candidate failure, and claimed endpoint-defeating force & no-sixth-role closure \\
\(\EndpointImpactSAT(R,model)\) & endpoint-impact bundle for same-carrier lower-bound use: target, carrier, candidate image, branch, standing, reportability, reuse, exclusion, act-time status, and lawful redescription & role audit \\
\(\mathsf{BN}_{\mathsf{SAT}}(model)\) & shorthand for \(\BridgeNeutralizedSAT(model)\): positive candidate-image occupation that neutralizes universal candidate failure on the same context-bound candidate image & bridge-neutralization \\
\(\LowerBoundGov(R,model)\) & optional shorthand for endpoint-deployed lower-bound governance when support, bookkeeping, bridge-neutralization, and carrier shift are removed & endpoint governance \\
\(\SepImgOcc\) & separator endpoint-image occupation: candidate-image exclusion used as the official endpoint-resolving negative branch on the fixed carrier; it contains \(\CandImageExcl(model)\) but is not entailed by support-level exclusion alone & endpoint occupation \\
\(\EndpointUse\) & branch \(B\) offered as theorem-bearing endpoint use & endpoint act \\
\(\OrdOfficialRes\) & ordinary mathematical use of a theorem-bearing act as the official resolution of branch \(B\) for a fixed endpoint; equivalent to endpoint use by Theorem~\ref{thm:ordinary-official-resolution-is-endpoint-use} & endpoint act \\
\(\OrdThm\) & ordinary theoremhood/proof legitimacy before endpoint-role classification & theorem status \\
\(\mathsf{NoFifthOrdNonOcc}\) & closure theorem stating that ordinary theorem-bearing non-occupation is not a fifth endpoint role after SAT projection and official endpoint use & endpoint discipline \\
\(\mathsf{RemainingFork}\) & final same-mode objection fork: projection failure, trichotomy failure, or \(\Aplus\) failure & denial burden \\
\begin{tabular}[t]{@{}l@{}}\(\mathsf{OfficialEndpoint}\)\\\(\mathsf{Resolution}\)\end{tabular} & ordinary theoremhood of \(\neg\PosRaw\) offered as the official negative endpoint resolution & endpoint act \\
\begin{tabular}[t]{@{}l@{}}\(\EndpointClosureEval\)\end{tabular} & endpoint-closure counterforce discipline for the two-branch SAT target; open/unresolved status is epistemic, not an endpoint outcome & endpoint evaluation \\
\begin{tabular}[t]{@{}l@{}}\(\UnresolvedEvalStatus\)\end{tabular} & epistemic or review status of the problem before or outside proof; not a third endpoint outcome & epistemic status \\
\begin{tabular}[t]{@{}l@{}}\(\OfficialSATEval\)\end{tabular} & official endpoint-closure evaluation of the raw \(\PosRaw\)/\(\SepBare\) split on the fixed unrestricted CNF-SAT carrier; not endpoint-complete evaluation and not a truth assumption & endpoint evaluation \\
\begin{tabular}[t]{@{}l@{}}\(\mathsf{PosEndpoint}\)\\\(\mathsf{Outcome}_R\)\end{tabular} & positive endpoint outcome in the official SAT endpoint evaluation; entails \(\PosRaw\) & endpoint outcome \\ 
\begin{tabular}[t]{@{}l@{}}\(\mathsf{NegEndpoint}\)\\\(\mathsf{Outcome}_R\)\end{tabular} & official negative endpoint outcome; stronger than a raw hypothetical branch and unpacks to \(\SepBare\wedge\OfficialEndpointResolution_R(\neg\PosRaw)\); not used to discharge the local reductio countercase & endpoint outcome \\
\begin{tabular}[t]{@{}l@{}}\(\mathsf{Annotation}\)\\\(\mathsf{Discharge}\)\end{tabular} & admissibility rule for exact-complement reductio annotation: contradiction in the annotated subproof discharges the original assumption \([\SepBare]^i\), not a strengthened object-level conjunction & proof rule \\
\begin{tabular}[t]{@{}l@{}}\(\mathsf{SATEndpoint}\)\\\(\mathsf{ImageOcc}\)\end{tabular} & occupation of the fixed positive CNF-SAT endpoint image & endpoint image \\
\begin{tabular}[t]{@{}l@{}}\(\mathsf{Positive}\)\\\(\mathsf{CandidateOcc}\)\end{tabular} & existence of an encoded polynomial-time candidate whose decoded procedure has no SAT failure & endpoint image \\
\begin{tabular}[t]{@{}l@{}}\(\mathsf{CNFRoute}\)\\\((R,model)\)\end{tabular} & standard unrestricted finite CNF-SAT correspondence used as the route of official resolution, not merely as lower-bound support & projection discipline \\
\begin{tabular}[t]{@{}l@{}}\(\mathsf{Projection}\)\\\(\mathsf{RoleContent}\)\end{tabular} & endpoint-role content fixed by the CNF-SAT projection: positive occupation for \(P=NP\), candidate-image non-occupation for \(P\neq NP\) & projection discipline \\
\(\ImageAdeq_{\mathsf{SAT}}\) & carrier-adequacy bridge \(\PosRaw\leftrightarrow\SATEndpointImageOcc(model)\) and \(\SepBare\leftrightarrow\CandImageExcl(model)\) on the context-bound candidate image & carrier adequacy \\
\(\OutcomePvNP\) & symmetric two-output problem carrier \(\{P=NP,P\neq NP\}\); distinct from the positive candidate-image carrier & outcome carrier \\
\begin{tabular}[t]{@{}l@{}}\(\mathsf{Endpoint}\)\\\(\mathsf{StatusWork}\)\end{tabular} & functional endpoint-status determination by a theorem-bearing act & consequence layer \\
\begin{tabular}[t]{@{}l@{}}\(\mathsf{EndpointRole}\)\\\(\mathsf{Discipline}\)\end{tabular} & classification of an official endpoint-resolution act by its role on the fixed endpoint image; ordinary theoremhood does not exempt it & consequence layer \\
\begin{tabular}[t]{@{}l@{}}\(\mathsf{NonPositive}\)\\\(\mathsf{StatusWork}\)\end{tabular} & same-domain endpoint-status work assigning non-occupation to the fixed candidate image & SAT asymmetry \\
\begin{tabular}[t]{@{}l@{}}\(\mathsf{LawfulCoequal}\)\\\(\mathsf{ImageOcc}\)\end{tabular} & lawful coequal occupation inside the same endpoint-image carrier; shown absent for SAT negative occupation & endpoint image \\
\begin{tabular}[t]{@{}l@{}}\(\mathsf{Reductio}\)\\\(\mathsf{Countercase}\)\end{tabular} & ordinary local proof role in which \(\SepBare\) is introduced as the live exact-complement assumption against \(\PosRaw\); not theoremhood, official resolution, or global outcome & endpoint discipline \\
\begin{tabular}[t]{@{}l@{}}\(\mathsf{Local}\)\\\(\mathsf{Countercase}\)\end{tabular} & local endpoint-countercase standing induced by the exact-complement reductio assumption; temporary derivational standing, not theoremhood & endpoint discipline \\
\begin{tabular}[t]{@{}l@{}}\(\mathsf{Endpoint}\)\\\(\mathsf{Counterforce}\)\end{tabular} & endpoint-facing force by which a raw branch is asserted to defeat or displace endpoint closure on the fixed carrier; the local countercase form routes to \(\LocalSepImgOcc(R,model)\), not to global endpoint outcome & endpoint discipline \\
\begin{tabular}[t]{@{}l@{}}\(\mathsf{SingleOcc}\)\\\(\mathsf{Interface}_{\mathsf{SAT}}\)\end{tabular} & theorem that the fixed SAT endpoint image is exhausted by positive candidate occupation & endpoint image \\
\begin{tabular}[t]{@{}l@{}}\(\mathsf{SingleInterface}\)\\\(\mathsf{NonExplosive}\)\end{tabular} & single-interface typing classifies non-occupation but does not by itself yield contradiction or ban lower-bound content & endpoint image \\
\(\EndpointDisc\) & endpoint-status discriminator attached to branch \(B\) & consequence layer \\
\(\IndDisc\) & independent same-domain governance work by \(C\) & forbidden case \\
\(\DirPos\) & direct constructive occupation of the positive CNF-SAT endpoint & SAT asymmetry \\
\(\SepDisc\) & independent same-domain endpoint-status discriminator induced by separator endpoint-image occupation & SAT asymmetry \\
\(\SepIndBridge\) & branch-local SAT independence bridge invoked by endpoint-resolving \(\SepImgOcc\) on the context-bound CNF model & SAT asymmetry \\
\(\NoSepDisc\) & kernel-forced no-independent-separator-classifier closure triggered by separator endpoint-image use & consequence layer \\
\(\ModelOfCnf\) & context-bound encoded CNF-SAT model for the fixed SAT endpoint carrier & model adequacy \\
\(\ModelAdeq\) & model-adequacy packet: coverage, soundness, failure adequacy, and candidate-image exactness & model adequacy \\
\(\ModelCoverage\) & every deterministic polynomial-time CNF-SAT candidate has an encoded polynomial-time code & model adequacy \\
\(\ModelSound\) & every encoded polynomial-time code decodes to a polynomial-time CNF-SAT candidate & model adequacy \\
\(\FailureAdeq\) & equal decoded procedures have the same SAT-failure status & model adequacy \\
\(\ImageExact\) & \(\CnfSATP\) is exactly existence of an encoded non-failing polynomial-time candidate & model adequacy \\
\(\mathsf{StdModel}_{\mathsf{CNF}}\) & standard uniform encoded model whose codes are machine/bound pairs \((e,k)\) over arbitrary finite CNF inputs & model adequacy \\
\(\FailsSAT\) & wrong answer, unlawful output, nontermination within the declared bound, or bound failure on some lawful finite CNF input & model adequacy \\
\(\StdLB\) & standard SAT lower-bound normal form: every encoded polynomial-time candidate fails SAT on the fixed CNF carrier & SAT-native bridge \\
\(\CandImageExcl\) & candidate endpoint-image exclusion: no encoded polynomial-time candidate occupies the positive SAT endpoint image & SAT-native bridge \\
\(\ThmLevelDisc\) & theorem-level candidate-status discriminator; not an algorithmic witness selector & SAT-native bridge \\
\(\mathsf{CandImageUse}\) & candidate-image exclusion is lawful support unless used as endpoint-resolving separator governance & closing hinge \\
\(\CandImageGov\) & endpoint-status governance by candidate-image exclusion & closing hinge \\
\begin{tabular}[t]{@{}l@{}}\(\mathsf{CandImage}\)\\\(\mathsf{UseKind}\)\end{tabular} & four-way use classification for candidate-image exclusion & closing hinge \\
\(\mathsf{OfficialNegUse}\) & official negative endpoint use on the context-bound finite CNF-SAT model & SAT negative occupation \\
\(\mathsf{NegOccExhaust}\) & four-way exhaustion of official negative occupation alternatives & SAT negative occupation \\
\(\mathsf{GovUse}\) & governed endpoint use on the finite CNF-SAT endpoint carrier & endpoint-status bivalence \\
\(\mathsf{EndpointStatus}\) & two-status Lean surface: positive or separator & endpoint-status bivalence \\
\(\Aplus\) & kernel-forced consequence package for endpoint use & consequence layer \\
\(\Aplus_R(B)\) binding & ordinary official endpoint resolution supplies endpoint use, so the \(\Aplus\) no-independent-discriminator consequences bind any independent role identified by trichotomy & consequence layer \\
\(\Asep\) & role-specific separator status factor, downstream only & classification \\
\(\Apos\) & positive report status factor, downstream only & classification \\
\(\PosClay\) & \(\PosRaw\wedge\Apos\) & positive Clay mathematical endpoint \\
\(\KSAT\) & local SAT kernel packet forced by non-degenerate fixed-carrier endpoint use & local kernel \\
\(\CoreSAT\) & core SAT endpoint adequacy packet: reference, standing, admissibility, irreversibility, target determinacy, report evaluability, act-time finality, and same-regime fidelity & weakening resistance \\
\(\WeakSep\) & alleged strict weaker same-carrier regime permitting separator governance & weakening resistance \\
\begin{tabular}[t]{@{}l@{}}\(\mathsf{OrdNeg}\)\\\(\mathsf{RoleClass}\)\end{tabular} & ordinary negation is granted; endpoint-bearing negation is role-classified & SAT branch audit \\
\begin{tabular}[t]{@{}l@{}}\(\mathsf{SepNot}\)\\\(\mathsf{Coequal}\)\end{tabular} & separator branch is not a lawful coequal endpoint role inside the positive endpoint-image carrier & SAT branch audit \\
\begin{tabular}[t]{@{}l@{}}\(\mathsf{NegComplexity}\)\\\(\mathsf{Lawful}\)\end{tabular} & negative complexity results remain lawful unless used as forbidden endpoint governance & non-explosiveness \\
\(\RouteSAT\) & route-locus audit for endpoint-resolving negative SAT routes & route exhaustion \\
\bottomrule
\end{longtable}

\begin{thebibliography}{12}
\bibitem{CookClayPvsNP}
S. Cook.
\emph{The P versus NP Problem}.
Clay Mathematics Institute Millennium Problem description.

\bibitem{Kernel}
A. J. Maley.
\emph{Non-Degenerate Construction and the Kernel of Admissibility}.
2026.

\bibitem{UEAP}
A. J. Maley.
\emph{Claim Standing and Legitimacy}.
2026.

\bibitem{ATS}
A. J. Maley.
\emph{Anchor, Tensor, and Skin: A Fixed-Domain Decomposition Theorem for Admissible Description}.
2026.

\bibitem{Cook1971}
S. A. Cook.
The complexity of theorem-proving procedures.
In \emph{Proceedings of the Third Annual ACM Symposium on Theory of Computing},
pp. 151--158, 1971.

\bibitem{Levin1973}
L. A. Levin.
Universal search problems.
\emph{Problemy Peredachi Informatsii}, 9(3):115--116, 1973.
English translation in \emph{Problems of Information Transmission}, 9(3):265--266, 1973.

\bibitem{Karp1972}
R. M. Karp.
Reducibility among combinatorial problems.
In R. E. Miller, J. W. Thatcher, and J. D. Bohlinger, editors,
\emph{Complexity of Computer Computations},
pp. 85--103. Plenum Press, 1972.

\bibitem{AroraBarak}
S. Arora and B. Barak.
\emph{Computational Complexity: A Modern Approach}.
Cambridge University Press, 2009.

\bibitem{LeanAuditZenodo}
A. J. Maley.
\emph{AASC P vs NP SAT Endpoint Lean Audit}, version 1.0.0.
Zenodo, 2026. DOI: \url{https://doi.org/10.5281/zenodo.20593794}.

\bibitem{LeanAuditGitHub}
A. J. Maley.
\emph{AASC-PvsNP-SAT-Endpoint-Lean-Audit}.
GitHub repository, 2026. \url{https://github.com/somamaley-ux/AASC-PvsNP-SAT-Endpoint-Lean-Audit}.

\bibitem{GodelNonInstantiability}
A. J. Maley.
\emph{Non-Instantiability of G\"odel Incompleteness in the Unique Admissible Interior}.
Self-contained audit edition, 2026.
\end{thebibliography}

\end{document}
